\documentclass{article}
\usepackage[utf8]{inputenc}

\usepackage{mathpartir}
\usepackage{amsmath}
\usepackage{amsfonts}
\usepackage{amssymb}
\usepackage[standard,thref,thmmarks]{ntheorem}
\usepackage{listings}
\usepackage{stmaryrd}
\usepackage{parskip}
\usepackage{mathtools}
\usepackage{xcolor}
\usepackage{enumitem}

\newcommand{\num}[1]{\overline{#1}}
\newcommand{\mf}[1]{\mathsf{#1}}
\newcommand{\mc}[1]{\mathcal{#1}}
\newcommand{\mb}[1]{\mathbf{#1}}
\newcommand{\tb}[1]{\textbf{#1}}
\newcommand{\tf}[1]{\textsf{#1}}
\newcommand{\ts}[1]{\textsc{#1}}
\newcommand{\tn}[1]{\textnormal{#1}}
\newcommand{\semi}{\,;\,}
\newcommand{\dom}{\mf{dom}}
\newcommand{\po}{\sqsubseteq}

\newcommand{\mhl}[1]{\colorbox{yellow}{$#1$}}
\newcommand{\hl}[1]{\colorbox{yellow}{#1}}
\newcommand{\grey}[1]{\colorbox{lightgray}{#1}}
\newcommand{\mathgrey}[1]{\colorbox{lightgray}{$#1$}}
\newcommand{\llrr}[1]{\llbracket #1 \rrbracket}

\newcommand{\pair}[2]{\langle #1, #2 \rangle}
\newcommand{\trip}[3]{\langle #1, #2, #3 \rangle}
\newcommand{\four}[4]{\langle #1, #2, #3, #4 \rangle}
\newcommand{\map}[2]{[\num{#1 \mapsto #2}]}

\newcommand{\vdval}{\vdash_\Delta}
\newcommand{\vdexp}{\vdash_\Delta}
\newcommand{\vdtop}{\vdash^{\mf{Top}}_\Delta}
\newcommand{\vdprog}{\vdash^{\mf{Prog}}_\Delta}
\newcommand{\vdrt}{\Vdash_\Delta}
\newcommand{\vdtree}{\vdash^{\mf{Tree}}_\Delta}
\newcommand{\vdmu}{\vdash^{\mf{Mode}}_\Delta}
\newcommand{\vdcfg}{\vdash^{\mf{Cfg}}_\Delta}
\newcommand{\vdm}{\vDash_\Delta}
\newcommand{\vdd}{\vDash}
\newcommand{\wf}{\vdash^{\mf{wf}}_\Delta}
\newcommand{\vdashv}{\vdash_v}
\newcommand{\vdashe}{\vdash_e}

\newcommand{\topstep}[2]{\mapsto^{#1}_{#2}}
\newcommand{\Topstep}[2]{\Mapsto^{#1}_{#2}}
\newcommand{\pstep}{\mapsto^{\mf{Prog}}}
\newcommand{\Pstep}{\Mapsto^{\mf{Prog}}}
\newcommand{\cstep}{\hookrightarrow}
\newcommand{\bsinit}[1]{\Downarrow^{\mf{Init}}_{#1}}
\newcommand{\divinit}[1]{\Uparrow^{\mf{Init}}_{#1}}
\newcommand{\bssucc}[1]{\Downarrow^{\mf{Succ}}_{#1}}
\newcommand{\divsucc}[1]{\Uparrow^{\mf{Succ}}_{#1}}

\newcommand{\one}{\mb{1}}
\newcommand{\elet}[3]{\mf{let}\ #1 = #2\ \mf{in}\ #3}
\newcommand{\eletrec}[3]{\mf{let}\ \mf{rec}\ #1 = #2\ \mf{in}\ #3}
\newcommand{\erun}[1]{\mf{run}\ #1}
\newcommand{\eif}[3]{\mf{if}\ #1\ \mf{then}\ #2\ \mf{else}\ #3}
\newcommand{\eret}[1]{\mf{return}\ #1}
\newcommand{\ecase}[5]{\mf{case}\ #1\ \{#2 \Rightarrow #3 \mid #4 \Rightarrow #5\}}
\newcommand{\estate}[4]{\mf{let}\ (#1, #1_{\mf{set}}) = \mf{useState}^{#2}(#3)\ \mf{in}\ #4}
\newcommand{\vset}[2]{\mf{set}@^{#1}_{#2}}
\newcommand{\eref}[1]{\mf{ref}(#1)}
\newcommand{\eget}[1]{\ !#1}
\newcommand{\eset}[2]{#1 := #2}
\newcommand{\vrec}[4]{\mf{rec}\ #1(#2 : #3).#4}
\newcommand{\vcomp}[4]{\mf{comp}\ #1(#2 : #3).#4}

\title{React-tRanslation}
\author{Toby Ueno}

\begin{document}

\maketitle

\section{React-tRace Core}

\subsection{Program-Level Syntax}

Constructs which only appear in intermediate stages 
of evaluation are in \grey{grey}.
Other non-standard constructs are \hl{highlighted}.
\begin{align*}
\text{Value } v &::= x \mid () \mid s \mid \lambda x.e \mid [] \mid v :: v\\
&\quad\mid \mhl{\mc{C}} \mid \mhl{\mf{text}\ v} \mid \mhl{\mf{tag}(v, v)}
\mid \mhl{\mf{attr}(v, v)} \mid \mhl{\mf{onClick}\ v}\\
&\quad\mid \mathgrey{\pair{\mc{C}}{v}} \mid \mathgrey{\vset{\ell}{p}}\\
\text{Expression } e &::= \eret{v} \mid \elet{x}{e}{e} \mid v\ v \\
&\quad\mid \ecase{v}{[]}{e_1}{x_\mf{hd} :: x_\mf{tl}}{e_2}\\
\text{Top-level expr }\hat{e} &::= e \mid \mhl{\estate{x}{\ell}{v}{\hat{e}}} \mid \elet{x}{e}{\hat{e}}\\
\text{Program } P &::= e \mid \elet{x}{e}{P} \mid \mhl{\elet{\mc{C}(x)}{\hat{e}}{P}}\\
\text{Type } \tau &::= \one \mid \mf{String} \mid \tau \xrightarrow{E} \tau
\mid [\tau] \\
&\quad\mid \mhl{\mf{View}} \mid \mhl{\mf{Attr}}
\mid \mhl{\mf{Component}\ \tau} \mid \mhl{\mf{Setter}\ \tau}\\
\text{Effect } E &::= \epsilon \mid \mhl{\mf{Mut}}
\end{align*}

\emph{\tb{Map Ordering}}. Define partial order $\po$ on maps $M=\map{x}{y}$
(e.g. $\Gamma$, $\Pi$, etc.) such that $M_1 \po M_2$ if, for all $x \in \dom(M_1)$,
$M_1[x] = M_2[x]$. (Implicitly, $\dom(M_1) \subseteq \dom(M_2)$.)

\emph{\tb{Expression Labels}}. 
We define a function $\mc{L}$ on the \emph{syntax} of top-level expressions
to denote the set of labels $\ell$ in the component body:
\begin{align*}
&\mc{L}(e) = \emptyset\\
&\mc{L}(\elet{x}{e}{\hat{e}}) = \mc{L}(\hat{e})\\
&\mc{L}(\estate{x}{\ell}{v}{\hat{e}}) = \{\ell\} \cup \mc{L}(\hat{e})
\end{align*}

We overload $\mc{L}$ on \emph{programs} to extract
the set of component names in the program:
\begin{align*}
&\mc{L}(e) = \emptyset\\
&\mc{L}(\elet{x}{e}{P}) = \mc{L}(P)\\
&\mc{L}(\elet{\mc{C}(x)}{\hat{e}}{P}) = \{\mc{C}\} \cup \mc{L}(P)
\end{align*}

\emph{\tb{Internal Representations}}.
\begin{align*}
\text{Type context }\Gamma &::= \map{x}{\tau}\\
\text{Path context }\Pi &::= \map{p}{\mc{C}}\\
\text{Component context }\Delta &::= \map{\mc{C}}{\pair{\tau}{\map{\ell}{\tau}}}\\
\text{Phase }\phi &::= \mf{Init} \mid \mf{Succ}\\
\text{Realized component }\pi &::= \{\mf{spec}: \pair{\mc{C}}{v}, 
\mf{st}: \map{\ell}{v}, \mf{tree}: u, \mf{rerender}: \mf{tt} \mid \mf{ff}\}\\
\text{Tree memory }m &::= \map{p}{\pi}\\
\text{Component definitions }\mc{D} &::= \map{\mc{C}}{x.\hat{e}}
\end{align*}
We elide the grammar of "runtime values" $u$,
which are distinct from normal values $v$, until a later section.

We overload indexing notation on $\Delta$ for brevity,
i.e. $\Delta[\mc{C}]=\tau$ and $\Delta[\mc{C}][\ell]=\tau$.

We also use the notation $m \mid x = v$
to concisely denote nested map/record updates.
The above should be read as ``$m$ with field $x$ set to $v$''.
For instance, $m \mid m[p].\mf{st}[\ell] = v$ updates the
state field of $m$ at path $p$ and location $\ell$ to value $v$.

\subsection{Program-Level Typing}

Names of noteworthy rules are \hl{highlighted}.
Rules for typing internal representations,
i.e. that are necessary for soundness but not typechecking,
are \grey{greyed}. 
Similarly, $\Pi$ is unnecessary for typechecking.

\subsubsection{Value Typing}

Judgment: $\Gamma \semi \Pi \vdval v: \tau$

\begin{mathpar}
\inferrule*[right=TV-Var] {
  \Gamma[x] = \tau
} {
  \Gamma \semi \Pi \vdval x: \tau
}
\and
\inferrule*[right=TV-Unit] { } {
  \Gamma \semi \Pi \vdval () : \one
}
\and
\inferrule*[right=TV-Str] { } {
  \Gamma \semi \Pi \vdval s : \mf{String}
}
\and
\inferrule*[right=TV-Lam] {
  \Gamma[x \mapsto \tau_1] \semi \Pi \vdexp e : \tau_2 \mid E
} {
  \Gamma \semi \Pi 
  \vdval
  \lambda x.e: \tau_1 \xrightarrow{E} \tau_2
}
\and
\inferrule*[right=TV-Nil] { } {
  \Gamma \semi \Pi \vdval [] : [\tau]
}
\and
\inferrule*[right=TV-Cons] {
  \Gamma \semi \Pi \vdval v_1 : \tau
  \and
  \Gamma \semi \Pi \vdval v_2 : [\tau]
} {
  \Gamma \semi \Pi \vdval v_1 :: v_2 : [\tau]
}
\and
\inferrule*[right=\hl{TV-Comp}] {
  \Delta[\mc{C}] = \tau
} {
  \Gamma \semi \Pi \vdval \mc{C} : \mf{Component}\ \tau
}
\and
\inferrule*[right=\hl{TV-Text}] {
  \Gamma \semi \Pi \vdval v: \mf{String}
} {
  \Gamma \semi \Pi \vdval \mf{text}\ v: \mf{View}
}
\and
\inferrule*[right=\hl{TV-Tag}] {
  \Gamma \semi \Pi \vdval v_1: [\mf{Attr}]
  \and
  \Gamma \semi \Pi \vdval v_2: [\mf{View}]
} {
  \Gamma \semi \Pi \vdval \mf{tag}(v_1, v_2): \mf{View}
}
\and
\inferrule*[right=\hl{TV-Attr}] {
  \Gamma \semi \Pi \vdval v_k: \mf{String}
  \and
  \Gamma \semi \Pi \vdval v_v: \mf{String}
} {
  \Gamma \semi \Pi \vdval \mf{attr}(v_k, v_v): \mf{Attr}
}
\and
\inferrule*[right=\hl{TV-OnClick}] {
  \Gamma \vdval v: \one \xrightarrow{E} \one
} {
  \Gamma \vdval \mf{onClick}\ v: \mf{Attr}
}
\and
\inferrule*[right=\grey{TV-Spec}] {
  \Gamma \semi \Pi \vdval \mc{C}: \mf{Component}\ \tau
  \and
  \Gamma \semi \Pi \vdval v: \tau
} {
  \Gamma \semi \Pi \vdval \pair{\mc{C}}{v}: \mf{View}
}
\and
\inferrule*[right=\grey{TV-Set}] {
  \Delta[\Pi[p]][\ell] = \tau
} {
  \Gamma \semi \Pi \vdval \vset{\ell}{p} : \mf{Setter}\ \tau
}
\end{mathpar}

\subsubsection{Expression Typing}

Judgment: $\Gamma \semi \Pi \vdexp e : \tau \mid E$

\begin{mathpar}
\inferrule*[right=TE-Ret] {
  \Gamma \semi \Pi \vdval v: \tau
} {
  \Gamma \semi \Pi \vdexp \eret{v}: \tau \mid \epsilon
}
\and
\inferrule*[right=TE-Let] {
  \Gamma \semi \Pi \vdexp e_1 : \tau_1 \mid E
  \and
  \Gamma[x \mapsto \tau_1] \semi \Pi \vdexp e_2 : \tau_2 \mid E
} {
  \Gamma \semi \Pi \vdexp \elet{x}{e_1}{e_2}: \tau_2 \mid E
}
\and
\inferrule*[right=TE-AppL] {
  \Gamma \semi \Pi \vdval v_1 : \tau_1 \xrightarrow{E} \tau_2
  \and
  \Gamma \semi \Pi \vdval v_2 : \tau_1
} {
  \Gamma \semi \Pi \vdexp v_1\ v_2: \tau_2 \mid E
}
\and
\inferrule*[right=\hl{TE-AppC}] {
  \Gamma \semi \Pi \vdval v_1 : \mf{Component}\ \tau
  \and
  \Gamma \semi \Pi \vdval v_2 : \tau
} {
  \Gamma \semi \Pi \vdexp v_1\ v_2: \mf{View} \mid \epsilon
}
\and
\inferrule*[right=\hl{TE-AppS}] {
  \Gamma \semi \Pi \vdval v_1 : \mf{Setter}\ \tau
  \and
  \Gamma \semi \Pi \vdval v_2 : \tau
} {
  \Gamma \semi \Pi \vdexp v_1\ v_2: \one \mid \mf{Mut}
}
\and
\inferrule*[right=TE-Case] {
  \Gamma \semi \Pi \vdval v : [\tau_1]
  \and
  \Gamma \semi \Pi \vdexp e_1 : \tau_2 \mid E
  \\
  \Gamma[x_\mf{hd} \mapsto \tau_1, x_\mf{tl} \mapsto [\tau_1]]
  \semi \Pi \vdexp e_2 : \tau_2 \mid E
} {
  \Gamma \semi \Pi \vdexp \ecase{v}{[]}{e_1}{x_\mf{hd} :: x_\mf{tl}}{e_2} : \tau \mid E
}
\and
\inferrule*[right=\hl{TE-Sub}] {
  \Gamma \semi \Pi \vdexp e: \tau \mid \epsilon
} {
  \Gamma \semi \Pi \vdexp e: \tau \mid \mf{Mut}
}
\end{mathpar}

\subsubsection{Top-Level Expression Typing}

Judgment: $\Gamma \semi \Pi \semi \mc{C} \vdtop \hat{e}$
\begin{mathpar}
\inferrule*[right=TH-Lift] {
  \Gamma \semi \Pi \vdexp e: \mf{View} \mid \epsilon
} {
  \Gamma \semi \Pi \semi \mc{C} \vdtop e
}
\and
\inferrule*[right=TH-Let] {
  \Gamma \semi \Pi \vdexp e: \tau_1 \mid \epsilon
  \and
  \Gamma[x \mapsto \tau_1] \semi \Pi \semi \mc{C} \vdtop \hat{e}
} {
  \Gamma \semi \Pi \semi \mc{C} \vdtop \elet{x}{e}{\hat{e}}
}
\and
\inferrule*[right=\hl{TH-State}] {
  \Delta[\mc{C}][\ell] = \tau_1
  \and
  \Gamma \semi \Pi \vdval v: \tau_1
  \\
  \Gamma[x \mapsto \tau_1, x_\mf{set} \mapsto \mf{Setter}\ \tau_1]
  \semi \Pi \semi \mc{C} \vdtop \hat{e}
} {
  \Gamma \semi \Pi \semi \mc{C} \vdtop \estate{x}{\ell}{v}{\hat{e}}
}
\end{mathpar}

\subsubsection{Program Typing}

Judgment: $\Gamma \vdprog P$

\begin{mathpar}
\inferrule*[right=TP-Lift] {
  \Gamma \semi \cdot \vdexp e : \mf{View} \mid \epsilon
} {
  \Gamma \vdprog e
}
\and
\inferrule*[right=TP-Let] {
  \Gamma \semi \cdot \vdexp e: \tau \mid \epsilon
  \and 
  \Gamma[x \mapsto \tau] \vdprog P
} {
  \Gamma \vdprog \elet{x}{e}{P}
}
\and
\inferrule*[right=\hl{TP-Def}] {
  \Gamma[x \mapsto \Delta[\mc{C}]] \semi \cdot \semi \mc{C} \vdtop \hat{e}
  \and
  \Gamma \vdprog P
} {
  \Gamma \vdprog \elet{\mc{C}(x)}{\hat{e}}{P}
}
\end{mathpar}

\subsubsection{Memory Typing}

Define $\Gamma \semi \Pi \vdm m$ if, for all $p \in \mf{dom}(m)$:
\begin{itemize}
\item for $m[p].\mf{spec}=\pair{\mc{C}}{v}$,
$\Pi[p] = \mc{C}$ and $\Gamma \semi \Pi \vdval v: \Delta[\mc{C}]$,
\item for all $\ell$ in $m[p].\mf{st}$,
$\Gamma \semi \Pi \vdval m[p].\mf{st}[\ell] : \Delta[\mc{C}][\ell]$, and
\item if $m[p].\mf{tree}$ is defined, then 
$\Gamma \semi \Pi \vdrt m[p].\mf{tree}: \mf{Tree}$.
\end{itemize}

Note that we have not yet defined what it means
for a runtime value such as $m[p].\mf{tree}$,
to be well-typed; we postpone this notion for later,
as doing so is unnecessary for the program-level proofs,
in which no trees are initialized.

\subsubsection{Component Context Typing}
Define $\Gamma \semi \Delta \vdd \mc{D}$ to say that,
for all $\mc{C} \in \dom(\mc{D})$
(such that $\mc{D}[\mc{C}] = x.\hat{e}$),
we have $\Gamma[x \mapsto \Delta[\mc{C}]] \semi \cdot \semi \mc{C} \vdtop \hat{e}$.

\subsection{Program-Level Semantics}

We \hl{highlight} interesting rules.

\subsubsection{Expression Semantics}

Judgment: $m \semi e \mapsto e' \semi m'$
\begin{mathpar}
\inferrule*[right=E-Let] {
  m \semi e_1 \mapsto e_1' \semi m'
} {
  m \semi \elet{x}{e_1}{e_2} \mapsto \elet{x}{e_1'}{e_2} \semi m'
}
\end{mathpar}
\begin{align*}
m \semi \elet{x}{\eret{v}}{e} &\mapsto e[v/x] \semi m
&&\ts{E-Ret}\\
m \semi (\lambda x.e) v &\mapsto e[v/x] \semi m
&&\ts{E-AppL}\\
m \semi \mc{C}\ v &\mapsto \eret{\pair{\mc{C}}{v}} \semi m
&&\ts{\hl{E-AppC}}\\
m \semi \vset{\ell}{p}\ v &\mapsto \eret{()} \semi m'
&&\ts{\hl{E-AppS}}\\
(m' = m \mid m[p].\mf{st}[\ell] = v\ 
    &\mf{and}\ m[p].\mf{rerender} = \mf{tt})\\
m \semi \ecase{[]}{[]}{e_1}{x_\mf{h} :: x_\mf{t}}{e_2} &\mapsto e_1 \semi m
&&\ts{E-CaseN}\\
m \semi \ecase{v_\mf{h} :: v_\mf{t}}{[]}{e_1}{x_\mf{h} :: x_\mf{t}}{e_2} 
&\mapsto e_2[v_\mf{h} / x_\mf{h}, v_\mf{t} / x_\mf{t}] \semi m
&&\ts{E-CaseC}
\end{align*}

\subsubsection{Top-Level Expression Semantics}

Judgment: $m \semi \hat{e} \topstep{\phi}{p} \hat{e}' \semi m'$
\begin{mathpar}
\inferrule*[right=EH-Lift] {
  m \semi e \mapsto e' \semi m
} {
  m \semi e \topstep{\phi}{p} e' \semi m
}
\and
\inferrule*[right=EH-Let] {
  m \semi e \mapsto e' \semi m
} {
  m \semi \elet{x}{e}{\hat{e}} \topstep{\phi}{p} \elet{x}{e'}{\hat{e}} \semi m
}
\end{mathpar}
\begin{align*}
m \semi \elet{x}{\eret{v}}{\hat{e}} &\topstep{\phi}{p} \hat{e}[v/x] \semi m
&&\ts{EH-Ret}\\
m \semi \estate{x}{\ell}{v}{\hat{e}}
&\topstep{\mf{Init}}{p}
\hat{e}[v/x, \vset{\ell}{p}/x_\mf{set}] \semi m'
&&\ts{\hl{EH-StI}}\\
(m' &= m \mid m[p].\mf{st}[\ell] = v)\\
m \semi \estate{x}{\ell}{v}{\hat{e}}
&\topstep{\mf{Succ}}{p}
\hat{e}[v'/x, \vset{\ell}{p}/x_\mf{set}] \semi m
&&\ts{\hl{EH-StS}}\\
(v' &= m[p].\mf{st}[\ell])
\end{align*}

\subsubsection{Program Expression Semantics}

Judgment: $\mc{D} \semi P \pstep P' \semi \mc{D}'$
\begin{mathpar}
\inferrule*[right=EP-Lift] {
  \cdot \semi e \mapsto e' \semi \cdot
} {
  \mc{D} \semi e \pstep e' \semi \mc{D}
}
\and
\inferrule*[right=EP-Let] {
  \cdot \semi e \mapsto e' \semi \cdot
} {
  \mc{D} \semi \elet{x}{e}{P}
  \pstep
  \elet{x}{e'}{P} \semi \mc{D}
}
\end{mathpar}
\begin{align*}
\mc{D} \semi \elet{x}{\eret{v}}{P}
&\pstep
P[v/x] \semi \mc{D}
&&\ts{EP-Ret}\\
\mc{D} \semi \elet{\mc{C}(x)}{\hat{e}}{P}
&\pstep
P \semi \mc{D}[\mc{C} \mapsto x.\hat{e}]
&&\ts{\hl{EP-Def}}
\end{align*}

\subsection{Program-Level Soundness}

Throughout, \emph{propositions} refer to (as-of-yet) unproved assumptions.

\subsubsection{General Lemmas}

We postulate weakening on $\Gamma$ for all typing judgments involving $\Gamma$.
For instance, the statement for weaking $\Gamma$ on values is:
\begin{proposition}[Type context weakening, values]
\label{thm:wktpv}
If $\Gamma \semi \Pi \vdval v: \tau$ and $\Gamma \po \Gamma'$,
then $\Gamma' \semi \Pi \vdval v: \tau$.
\end{proposition}

Similarly, we require weakening on $\Pi$ for all typing judgments involving $\Pi$:
\begin{proposition}[Path context weakening, values]
\label{thm:wkpiv}
If $\Gamma \semi \Pi \vdval v: \tau$ and $\Pi \po \Pi'$,
then $\Gamma \semi \Pi' \vdval v: \tau$.
\end{proposition}

\begin{proposition}[Path context weakening, top-level expressions]
\label{thm:wkpih}
If both $\Gamma \semi \Pi \semi \mc{C} \vdtop \hat{e}$ and $\Pi \po \Pi'$,
then $\Gamma \semi \Pi' \semi \mc{C} \vdtop \hat{e}$.
\end{proposition}

\begin{proposition}[Path context weakening, memory]
\label{thm:wkpim}
If both $\Gamma \semi \Pi \vdm m$ and $\Pi \po \Pi'$,
then $\Gamma \semi \Pi' \vdm m$.
\end{proposition}

We elide statements of all other weakening lemmas.

\begin{proposition}[Canonical Forms]
\label{thm:canon}
If $\cdot \semi \Pi \vdval v: \tau$, then:
\begin{itemize}
\item if $\tau=\one$, then $v=()$.
\item if $\tau=\mf{String}$, then $v=s$.
\item if $\tau=\tau_1 \xrightarrow{E} \tau_2$, then $v=\lambda x: \tau_1.e$.
\item if $\tau=[\tau']$, then either $v=[]$ or $v = v_1 :: v_2$.
\item if $\tau=\mf{View}$, then either $v=\mf{text}\ v'$,
    $v=\mf{tag}(v' \num{v_i}, \num{v_j})$, or $v=\pair{\mc{C}}{v'}$.
\item if $\tau=\mf{Attr}$, then either $v=\mf{attr}(v_1, v_2)$ or $v=\mf{onClick}\ v'$.
\item if $\tau=\mf{Component}\ \tau'$, then $v=\mc{C}$.
\item if $\tau=\mf{Setter}\ \tau'$, then $v=\vset{\ell}{p}$.
\end{itemize}
\end{proposition}

We postulate substitution lemmas on all typing judgments involving $\Gamma$.
\begin{proposition}[Substitution lemma, expressions.]
\label{thm:subexp}
If $\Gamma[x \mapsto \tau_v] \semi \Pi \vdexp e: \tau \mid E$
and $\Gamma \semi \Pi \vdval v: \tau_v$,
then $\Gamma \semi \Pi \vdexp e[v/x]: \tau \mid E$.
\end{proposition}

\begin{proposition}[Substitution lemma, top-level expressions.]
\label{thm:subtle}
If\\$\Gamma[x \mapsto \tau_v] \semi \Pi \semi \mc{C} \vdtop \hat{e}$
and $\Gamma \semi \Pi \vdval v: \tau_v$,
then $\Gamma \semi \Pi \semi \mc{C} \vdtop \hat{e}[v/x]$.
\end{proposition}

\begin{proposition}[Substitution lemma, programs.]
\label{thm:subprog}
If\\$\Gamma[x \mapsto \tau_v] \vdprog P$
and $\Gamma \semi \cdot \vdval v: \tau_v$,
then $\Gamma \vdprog \hat{e}[v/x]$.
\end{proposition}

Finally, we have the following properties regarding syntactic function $\mc{L}$:
\begin{proposition}[Label-set substitution]
\label{thm:sublab}
For all top-level expressions $\hat{e}$, values $v$, 
and variables $x$, 
$\mc{L}(\hat{e}[v/x]) = \mc{L}(\hat{e})$.
\end{proposition}

\begin{proposition}[Component-name substitution]
\label{thm:subname}
For all programs $P$, all values $v$, 
and all variables $x$, 
$\mc{L}(P[v/x]) = \mc{L}(P)$.
\end{proposition}

\subsubsection{Expressions}

\begin{lemma}[Progress of expressions.]
\label{thm:progexp}
If $\cdot \semi \Pi \vdexp e: \tau \mid E$ and $\cdot \semi \Pi \vdm m$,
then either:
\begin{enumerate}[label=(\roman*)]
\item $e = \eret{v}$, or
\item there exist some $e'$ and $m'$ such that
$m \semi e \mapsto e' \semi m'$.
\end{enumerate}


\begin{Proof}
By induction on the derivation of the expression typing judgment.

\tb{Case} \ts{TE-Ret}: We have 
$\cdot \semi \Pi \vdexp \eret{v}: \tau \mid \epsilon$,
and the theorem holds trivially by case \emph{(i)}.

\tb{Case} \ts{TE-Let}: \begin{mathpar}
\inferrule*[right=TE-Let] {
  \inferrule*[right=(1)] { } {\cdot \semi \Pi \vdexp e_1 : \tau_1 \mid E}
  \and
  \inferrule*[right=(2)] { } {
    [x \mapsto \tau_1] \semi \Pi \vdexp e_2 : \tau_2 \mid E
  }
} {
  \cdot \semi \Pi \vdexp \elet{x}{e_1}{e_2}: \tau_2 \mid E
}
\end{mathpar}
\begin{itemize}
\item By inductive hypothesis on (1), we have two subcases on $e_1$:
either (a) $e_1 = \eret{v}$,
or (b) we have $m \semi e_1 \mapsto e_1' \semi m'$.
\begin{itemize}
\item Subcase (a): \ts{E-Ret} gives 
  $m \semi \elet{x}{\eret{v}}{e_2} \mapsto e_2[v/x] \semi m$.
\item Subcase (b):
  \begin{mathpar}
    \inferrule*[right=E-Let] {
      m \semi e_1 \mapsto e_1' \semi m'
    } {
      m \semi \elet{x}{e_1}{e_2} \mapsto \elet{x}{e_1'}{e_2} \semi m'
    }
  \end{mathpar}
\end{itemize}
\item Thus, in any subcase, $\elet{x}{e_1}{e_2}$ steps,
and we satisfy conclusion \emph{(ii)}.
\end{itemize}

\tb{Case} \ts{TE-AppL}: \begin{mathpar}
\inferrule*[right=TE-AppL] {
  \cdot \semi \Pi \vdval v_1 : \tau_1 \xrightarrow{E} \tau_2
  \and
  \cdot \semi \Pi \vdval v_2 : \tau_1
} {
  \cdot \semi \Pi \vdexp v_1\ v_2: \tau_2 \mid E
}
\end{mathpar}\begin{itemize}
\item By canonical forms, $v_1=\lambda x.e$.
\item We apply \ts{E-AppL} to get that
$m \semi (\lambda x.e)\ v_2 \mapsto e[v_2/x] \semi m$,
satisfying conclusion \emph{(ii)}.
\end{itemize}

\tb{Case} \ts{TE-AppC}: \begin{mathpar}
\inferrule*[right=TE-AppC] {
  \cdot \semi \Pi \vdval v_1 : \mf{Component}\ \tau
  \and
  \cdot \semi \Pi \vdval v_2 : \tau
} {
  \cdot \semi \Pi \vdexp v_1\ v_2: \mf{View} \mid \epsilon
}
\end{mathpar}\begin{itemize}
\item By canonical forms, $v_1 = \mc{C}$.
\item \ts{E-AppC} gives
$m \semi \mc{C}\ v_2 \mapsto \eret{\pair{\mc{C}}{v_2}} \semi m$,
satisfying conclusion \emph{(ii)}.
\end{itemize}

\tb{Case} \ts{TE-AppS}: \begin{mathpar}
\inferrule*[right=TE-AppS] {
  \cdot \semi \Pi \vdval v_1 : \mf{Setter}\ \tau
  \and
  \cdot \semi \Pi \vdval v_2 : \tau
} {
  \cdot \semi \Pi \vdexp v_1\ v_2: \one \mid \mf{Mut}
}
\end{mathpar}\begin{itemize}
\item By canonical forms, $v_1 = \vset{\ell}{p}$.
\item We apply \ts{E-AppS} to get that
$m \semi \vset{\ell}{p}\ v_2 \mapsto \eret{()} \semi m'$,
where $m' = m \mid m[p].\mf{st}[\ell] = v_2$ 
and $m[p].\mf{rerender} = \mf{tt}$,
satisfying conclusion \emph{(ii)}.
\end{itemize}

\tb{Case} \ts{TE-Case}: \begin{mathpar}
\inferrule*[right=TE-Case] {
  \cdot \semi \Pi \vdval v : [\tau_1]
  \and
  \cdot \semi \Pi \vdexp e_1 : \tau_2 \mid E
  \\
  [x_\mf{hd} \mapsto \tau_1, x_\mf{tl} \mapsto [\tau_1]]
  \semi \Pi \vdexp e_2 : \tau_2 \mid E
} {
  \cdot \semi \Pi \vdexp \ecase{v}{[]}{e_1}{x_\mf{hd} :: x_\mf{tl}}{e_2} : \tau \mid E
}
\end{mathpar}
\begin{itemize}
\item By \thref{thm:canon} (canonical forms), either $v=[]$ or $v=v_\mf{hd} :: v_\mf{tl}$, 
giving two subcases.
\begin{itemize}
\item If $v=[]$, then \ts{E-CaseN} gives\\
$m \semi \ecase{[]}{[]}{e_1}{x_\mf{hd} :: x_\mf{tl}}{e_2} \mapsto e_1 \semi m$.
\item If $v=v_\mf{hd} :: v_\mf{tl}$, ththen \ts{E-CaseC} gives\\
$m \semi \ecase{v_\mf{hd} :: v_\mf{tl}}{[]}{e_1}{x_\mf{hd} :: x_\mf{tl}}{e_2}
\mapsto e_2[v_\mf{hd} / x_\mf{hd}, v_\mf{tl} / x_\mf{tl}] \semi m$.
\end{itemize}
\item In either subcase, we satisfy conclusion \emph{(ii)}.
\end{itemize}

\tb{Case} \ts{TE-Sub}: \begin{mathpar}
\inferrule*[right=TE-Sub] {
  \Gamma \semi \Pi \vdexp e: \tau \mid \epsilon
} {
  \Gamma \semi \Pi \vdexp e: \tau \mid \mf{Mut}
}
\end{mathpar}
The result follows immediately from the inductive hypothesis
on the more conservative typing premise.
\end{Proof}
\end{lemma}

\begin{lemma}[Preservation of expressions.]
\label{thm:prezexp}
If
\begin{enumerate}[label=(\roman*)]
\item $\Gamma \semi \Pi \vdexp e: \tau \mid E$,
\item $\Gamma \semi \Pi \vdm m$, and
\item $m \semi e \mapsto e' \semi m'$ for some $e'$ and $m'$,
\end{enumerate}
then
\begin{enumerate}[label=(\roman*)]
\item $\Gamma \semi \Pi \vdexp e': \tau \mid E$ and
\item $\Gamma \semi \Pi \vdm m'$.
\end{enumerate}

\begin{Proof}
By induction on derivation of the step judgment \emph{(iii)}.

\tb{Case} \ts{E-Let}: We have \begin{mathpar}
\inferrule*[right=E-Let] {
  \inferrule*[right=(1)] { } {m \semi e_1 \mapsto e_1' \semi m'}
} {
  m \semi \elet{x}{e_1}{e_2} \mapsto \elet{x}{e_1'}{e_2} \semi m'
}
\and
\inferrule*[right=TE-Let] {
  \inferrule*[right=(2)] { } {\Gamma \semi \Pi \vdexp e_1 : \tau_1 \mid E}
  \and
  \inferrule*[right=(3)] { } {
    \Gamma[x \mapsto \tau_1] \semi \Pi \vdexp e_2 : \tau_2 \mid E
  }
} {
  \Gamma \semi \Pi \vdexp \elet{x}{e_1}{e_2}: \tau_2 \mid E
}
\end{mathpar}\begin{itemize}
\item We apply the inductive hypothesis with (1), (2), and premise \emph{(ii)}
to get: \begin{itemize}
\item $\Gamma \semi \Pi \vdm m'$, satisfying conclusion \emph{(ii)}, and
\item (4) $\Gamma \semi \Pi \vdexp e_1': \tau_1 \mid E$.
\end{itemize}
\item We reapply \ts{TE-Let} as follows to get conclusion \emph{(i)}:
\begin{mathpar}
\inferrule*[right=TE-Let] {
  \inferrule*[right=(4)] { } {\Gamma \semi \Pi \vdexp e_1' : \tau_1 \mid E}
  \and
  \inferrule*[right=(3)] { } {
    \Gamma[x \mapsto \tau_1] \semi \Pi \vdexp e_2 : \tau_2 \mid E
  }
} {
  \Gamma \semi \Pi \vdexp \elet{x}{e_1'}{e_2}: \tau_2 \mid E
}
\end{mathpar}
\end{itemize}

\tb{Case} \ts{E-Ret}: \begin{itemize}
\item We have $m \semi \elet{x}{\eret{v}}{e} \mapsto e[v/x] \semi m$ and:
\begin{mathpar}
\inferrule*[right=TE-Let] {
  \inferrule*[right=(1)] { } {\Gamma \semi \Pi \vdexp \eret{v} : \tau_1 \mid E}
  \ 
  \inferrule*[right=(2)] { } {
    \Gamma[x \mapsto \tau_1] \semi \Pi \vdexp e : \tau_2 \mid E
  }
} {
  \Gamma \vdexp \elet{x}{\eret{v}}{e}: \tau_2 \mid E
}
\end{mathpar}
\item (1) lets us derive (3) $\Gamma \semi \Pi \vdval v: \tau_1$
by subcasing on $E$ as follows:
\begin{mathpar}
\inferrule*[right=TE-Ret] {
  \inferrule*[right=(3)] { } {\Gamma \semi \Pi \vdval v: \tau_1}
} {
  \Gamma \semi \Pi \vdexp \eret{v}: \tau_1 \mid \epsilon
}
\and
\inferrule*[right=TE-Sub] {
  \inferrule*[right=TE-Ret] {
    \inferrule*[right=(3)] { } {\Gamma \semi \Pi \vdval v: \tau_1}
  } {
    \Gamma \semi \Pi \vdexp \eret{v}: \tau_1 \mid \epsilon
  }
} {
  \Gamma \semi \Pi \vdexp \eret{v}: \tau_1 \mid \mf{Mut}
}
\end{mathpar}
\item By \thref{thm:subexp} (substitution for expressions)
on (2) and (3), we get
$\Gamma \semi \Pi \vdexp e[v/x] : \tau_2 \mid E$ and we are done.
\end{itemize}

\tb{Case} \ts{E-AppL}: \begin{itemize}
\item We have $m \semi (\lambda x.e) v \mapsto e[v/x] \semi m$ and:
\begin{mathpar}
\inferrule*[right=TE-AppL] {
  \inferrule*[right=TV-Lam] {
    \inferrule*[right=(1)] { } {
      \Gamma[x \mapsto \tau_1] \semi \Pi \vdexp e: \tau_2 \mid E
    }
  } {
    \Gamma \semi \Pi \vdval \lambda x.e : \tau_1 \xrightarrow{E} \tau_2
  }
  \:
  \inferrule*[right=(2)] { } {
    \Gamma \semi \Pi \vdval v_2 : \tau_1
  }
} {
  \Gamma \semi \Pi \vdexp (\lambda x.e)\ v_2: \tau_2 \mid E
}
\end{mathpar}
\item By \thref{thm:subexp} (substitution for expressions)
on (1) and (2), we get
$\Gamma \semi \Pi \vdexp e[v/x] : \tau_2 \mid E$ and we are done.
\end{itemize}

\tb{Case} \ts{E-AppC}: \begin{itemize}
\item We have $m \semi \mc{C}\ v \mapsto \eret{\pair{\mc{C}}{v}} \semi m$ and:
\begin{mathpar}
\inferrule*[right=TE-AppC] {
  \inferrule*[right=(1)] { } {
    \Gamma \semi \Pi \vdval \mc{C} : \mf{Component}\ \tau
  }
  \and
  \inferrule*[right=(2)] { } {
    \Gamma \semi \Pi \vdval v : \tau
  }
} {
  \Gamma \semi \Pi \vdexp \mc{C}\ v: \mf{View} \mid \epsilon
}
\end{mathpar}
\item We can then derive conclusion \emph{(i)}:
\begin{mathpar}
\inferrule*[right=TE-Ret] {
  \inferrule*[right=TV-Spec] {
    \inferrule*[right=(1)] { } {\Gamma \semi \Pi \vdval \mc{C}: \mf{Component}\ \tau}
    \ 
    \inferrule*[right=(2)] { } {\Gamma \semi \Pi \vdval v: \tau}
  } {
    \Gamma \semi \Pi \vdval \pair{\mc{C}}{v} : \mf{View}
}
} {
  \Gamma \semi \Pi \vdexp \eret{\pair{\mc{C}}{v}}: \mf{View} \mid \epsilon
}
\end{mathpar}
\end{itemize}

\tb{Case} \ts{E-AppS}: \begin{itemize}
\item We have $m \semi \vset{\ell}{p}\ v \mapsto \eret{()} \semi m'$,
where $m' = m \mid m[p].\mf{st}[\ell] = v$ and $m[p].\mf{rerender} = \mf{tt}$.
\item The typing premise gives us:
\begin{mathpar}
\inferrule*[right=TE-AppS] {
  \inferrule*[right=TV-Set] {
    \inferrule*[right=(1)] { } {\Delta[\Pi[p]][\ell] = \tau}
  } {
    \Gamma \semi \Pi \vdval \vset{\ell}{p}: \mf{Setter}\ \tau
  }
  \and
  \inferrule*[right=(2)] { } {\Gamma \semi \Pi \vdval v: \tau}
} {
  \Gamma \semi \Pi \vdexp \vset{\ell}{p}\ v: \one \mid \mf{Mut}
}
\end{mathpar}
\item Conclusion \emph{(i)} follows from the derivation:
\begin{mathpar}
\inferrule*[right=TE-Sub] {
  \inferrule*[right=TE-Ret] {
    \inferrule*[right=TV-Unit] { } {
      \Gamma \semi \Pi \vdval (): \one
    }
  } {
    \Gamma \semi \Pi \vdexp \eret{()}: \one \mid \epsilon
  }
} {
  \Gamma \semi \Pi \vdexp \eret{()}: \one \mid \mf{Mut}
}
\end{mathpar}
\item To show $\Gamma \semi \Pi \vdm m'$ given premise $\Gamma \semi \Pi \vdm m$,
it suffices to show that 
$\Gamma \semi \Pi \vdval m'[p].\mf{st}[\ell] : \Delta[\Pi[p]][\ell]$.
\item By definition of $m'$, we want to show 
$\Gamma \semi \Pi \vdval v : \Delta[\Pi[p]][\ell]$.
\item This fact follows from (1) and (2), and we are done.
\end{itemize}

\tb{Case} \ts{E-CaseN}: \begin{itemize}
\item We have $m \semi \ecase{[]}{[]}{e_1}{x_\mf{hd} :: x_\mf{tl}}{e_2} \mapsto e_1 \semi m$ and:
\begin{mathpar}
\inferrule*[right=TE-If] {
  \inferrule*[right=(1)] { } {\Gamma \semi \Pi \vdexp e_1 : \tau \mid E}
  \and \cdots
} {
  \Gamma \semi \Pi \vdexp \eif{v}{e_1}{e_2} : \tau \mid E
}
\end{mathpar}
\item The conclusion follows immediately from (1).
\end{itemize}

\tb{Case} \ts{E-CaseC}: \begin{itemize}
\item We have $m \semi \ecase{v_\mf{h} :: v_\mf{t}}{[]}{e_1}{x_\mf{h} :: x_\mf{t}}{e_2} 
\mapsto e_2[v_\mf{h} / x_\mf{h}, v_\mf{t} / x_\mf{t}] \semi m$ and:
\begin{mathpar}
\inferrule*[right=TE-If] {
  \inferrule*[right=(1)] { } {\Gamma \semi \Pi \vdval v_\mf{h} :: v_\mf{t} : [\tau']}
  \and
  \inferrule*[right=(2)] { } {\Gamma' \semi \Pi \vdexp e_2 : \tau \mid E}
  \ \cdots
} {
  \Gamma \semi \Pi \vdexp \eif{v}{e_1}{e_2} : \tau \mid E
}
\end{mathpar}
where $\Gamma' = \Gamma[x_\mf{h} \mapsto \tau', x_\mf{t} \mapsto [\tau']]$.
\item We can further derive:
\begin{mathpar}
\inferrule*[right=TV-Cons] {
  \inferrule*[right=(3)] { } {\Gamma \semi \Pi \vdval v_\mf{h}: \tau'}
  \and
  \inferrule*[right=(4)] { } {\Gamma \semi \Pi \vdval v_\mf{t}: [\tau']}
} {
  \Gamma \semi \Pi \vdval v_\mf{h} :: v_\mf{t} : [\tau']\ (1)
}
\end{mathpar}
\item We apply the substitution lemma twice to (2) using values (3) and (4)
to get that 
$\Gamma \semi \Pi \vdval e_2[v_\mf{h} / x_\mf{h}, v_\mf{t} / x_\mf{t}]: \tau \mid E$
and we are done.
\end{itemize}
\end{Proof}
\end{lemma}

The above two lemmas are sufficient for type soundness of expressions,
but we also prove a "purity lemma" on expressions for later use:

\begin{lemma}[Purity of expressions.]
\label{thm:pure}
If $\cdot \semi \Pi \vdexp e: \tau \mid \epsilon$
and $m \semi e \mapsto e' \semi m'$,
then $m' = m$.
\begin{Proof}
By induction on the derivation of the typing judgment.

\tb{Case} \ts{E-Ret}: We have \begin{mathpar}
\inferrule*[right=TE-Ret] {
  \cdot \semi \Pi \vdval v: \tau
} {
  \cdot \semi \Pi \vdexp \eret{v}: \tau \mid \epsilon
}
\end{mathpar}
But there is no rule that steps an expression of form $\eret{v}$,
so the theorem holds vacuously.

\tb{Case} \ts{TE-Let}: We have \begin{mathpar}
\inferrule*[right=E-Let] {
  \inferrule*[right=(1)] { } {m \semi e_1 \mapsto e_1' \semi m'}
} {
  m \semi \elet{x}{e_1}{e_2} \mapsto \elet{x}{e_1'}{e_2} \semi m'
}
\and
\inferrule*[right=TE-Let] {
  \inferrule*[right=(2)] { } {\cdot \semi \Pi \vdexp e_1 : \tau_1 \mid \epsilon}
  \and
  \inferrule*[right=(3)] { } {
    [x \mapsto \tau_1] \semi \Pi \vdexp e_2 : \tau_2 \mid \epsilon
  }
} {
  \cdot \semi \Pi \vdexp \elet{x}{e_1}{e_2}: \tau_2 \mid \epsilon
}
\end{mathpar}
The inductive hypothesis on (1) and (2) gives that $m' = m$, and we are done.

\tb{Case} \ts{TE-AppL}: We have
\begin{mathpar}
\inferrule*[right=TE-AppL] {
  \inferrule*[right=(1)] { } {
    \cdot \semi \Pi \vdval v_1 : \tau_1 \xrightarrow{\epsilon} \tau_2
  }
  \and
  \cdot \semi \Pi \vdval v_2 : \tau_1
} {
  \cdot \semi \Pi \vdexp v_1\ v_2: \tau_2 \mid \epsilon
}
\end{mathpar}
\begin{itemize}
\item By \thref{thm:canon} (canonical forms) on (1), $v_1 = \lambda x.e$.
\item The only step rule which applies is \ts{E-AppL}, which says
that $m \semi (\lambda x.e) v \mapsto e[v/x] \semi m$,
i.e. $m' = m$ and we are done.
\end{itemize}

\tb{Case} \ts{TE-AppC}: We have
\begin{mathpar}
\inferrule*[right=TE-AppC] {
  \inferrule*[right=(1)] { } {
    \cdot \semi \Pi \vdval v_1 : \mf{Component}\ \tau
  }
  \and
  \cdot \semi \Pi \vdval v_2 : \tau
} {
  \cdot \semi \Pi \vdexp v_1\ v_2: \mf{View} \mid \epsilon
}
\end{mathpar}
\begin{itemize}
\item By \thref{thm:canon} (canonical forms) on (1), $v_1 = \mc{C}$.
\item The only step rule that applies is \ts{E-AppC}, which says
that $m \semi (\lambda x.e) v \mapsto \eret{\pair{\mc{C}}{v}} \semi m$,
i.e. $m' = m$ and we are done.
\end{itemize}

\tb{Case} \ts{TE-Case}: We have \begin{mathpar}
\inferrule*[right=TE-Case] {
  \cdot \semi \Pi \vdval v : [\tau_1]
  \and
  \cdot \semi \Pi \vdexp e_1 : \tau_2 \mid E
  \\
  [x_\mf{h} \mapsto \tau_1, x_\mf{t} \mapsto [\tau_1]]
  \semi \Pi \vdexp e_2 : \tau_2 \mid E
} {
  \cdot \semi \Pi \vdexp \ecase{v}{[]}{e_1}{x_\mf{h} :: x_\mf{t}}{e_2} : \tau \mid E
}
\end{mathpar}
\begin{itemize}
\item By \thref{thm:canon} (canonical forms) on (1), 
$v_1 = []$ or $v_1 = v_\mf{hd} :: v_\mf{tl}$.
\item If $v_1=[]$, then \ts{E-CaseN} says that\\
$m \semi \ecase{[]}{[]}{e_1}{x_\mf{h} :: x_\mf{t}}{e_2} \mapsto e_1 \semi m$.
\item If $v_1=\mf{false}$, then \ts{E-CaseC} says that\\
$m \semi \ecase{v_\mf{h} :: v_\mf{t}}{[]}{e_1}{x_\mf{h} :: x_\mf{t}}{e_2} 
\mapsto e_2[v_\mf{h} / x_\mf{h}, v_\mf{t} / x_\mf{t}] \semi m$.
\item In either case, $m' = m$ and we are done.
\end{itemize}

\tb{Cases} \ts{TE-AppS} and \ts{TE-Sub}
are both impossible, since they type their conclusions
with the $\mf{Mut}$ effect.
\end{Proof}
\end{lemma}

\subsubsection{Top-Level Expressions}

\begin{lemma}[Progress of top-level expressions, initial phase.]
\label{thm:progtlei}
If\\[0.2mm]
$\cdot \semi \Pi \semi \Pi[p] \vdtop \hat{e}$ and $\cdot \semi \Pi \vdm m$,
then either:
\begin{enumerate}[label=(\roman*)]
\item $\hat{e} = \eret{v}$, or
\item there exist some $\hat{e}'$ and $m'$ such that
$m \semi \hat{e} \topstep{\mf{Init}}{p} \hat{e}' \semi m'$.
\end{enumerate}

\begin{Proof}
By induction on the derivation of the expression typing judgment.

\tb{Case} \ts{TH-Lift}: \begin{mathpar}
\inferrule*[right=TH-Lift] {
  \inferrule*[right=(1)] { } {\cdot \semi \Pi \vdexp e: \mf{View} \mid \epsilon}
} {
  \cdot \semi \Pi \semi \mc{C} \vdtop e
}
\end{mathpar}\begin{itemize}
\item By \thref{thm:progexp} (progress of expressions) on (1), we have two subcases:
either $e=\eret{v}$ (and we are done by conclusion \emph{(i)}), 
or we have some $e'$ and $m'$ such that
(2) $m \semi e \mapsto e' \semi m'$.
\item By \thref{thm:pure} (purity) on (1) and (2), $m' = m$.
\item We derive the conclusion as follows:
\begin{mathpar}
\inferrule*[right=EH-Lift] {
  m \semi e \mapsto e' \semi m
} {
  m \semi e \topstep{\mf{Init}}{p} e' \semi m
}
\end{mathpar}
\end{itemize}

\tb{Case} \ts{TH-Let}: \begin{mathpar}
\inferrule*[right=TH-Let] {
  \inferrule*[right=(1)] { } {\cdot \semi \Pi \vdexp e: \tau_1 \mid \epsilon}
  \and
  [x \mapsto \tau_1] \semi \Pi \semi \mc{C} \vdtop \hat{e}
} {
  \cdot \semi \Pi \semi \mc{C} \vdtop \elet{x}{e}{\hat{e}}
}
\end{mathpar}\begin{itemize}
\item By \thref{thm:progexp} (progress of expressions) on (1), we have two subcases:
either (a) $e = \eret{v}$,
or (b) we have $m \semi e \mapsto e' \semi m'$.
\begin{itemize}
\item Subcase (a): \ts{EH-Ret} gives us that
$m \semi \elet{x}{\eret{v}}{\hat{e}} \topstep{\mf{Init}}{p} \hat{e}[v/x] \semi m$.
\item Subcase (b): by \thref{thm:pure} (purity) on (1), $m' = m$,
so we derive a step as follows:
\begin{mathpar}
\inferrule*[right=EH-Let] {
  m \semi e \mapsto e' \semi m
} {
  m \semi \elet{x}{e}{\hat{e}} \topstep{\mf{Init}}{p} \elet{x}{e'}{\hat{e}} \semi m
}
\end{mathpar}
\end{itemize}
\item Thus, in either subcase, we take a step and satisfy conclusion \emph{(ii)}.
\end{itemize}
\tb{Case} \ts{TH-State}: \begin{mathpar}
\inferrule*[right=TH-State] {
  \Delta[\mc{C}][\ell] = \tau_1
  \and
  \cdot \semi \Pi \vdval v: \tau_1
  \\
  [x \mapsto \tau_1, x_\mf{set} \mapsto \mf{Setter}\ \tau_1]
  \semi \Pi \semi \mc{C} \vdtop \hat{e}
} {
  \cdot \semi \Pi \semi \mc{C} \vdtop \estate{x}{\ell}{v}{\hat{e}}
}
\end{mathpar}
\begin{itemize}
\item Rule \ts{EH-StI} immediately applies, giving us
$$m \semi \estate{x}{\ell}{v}{\hat{e}}
\topstep{\mf{Init}}{p}
\hat{e}[v/x, \vset{\ell}{p}/x_\mf{set}] \semi m'$$
for $m' = m \mid m[p].\mf{st}[\ell] = v$.
\end{itemize}
\end{Proof}
\end{lemma}

\begin{lemma}[Preservation of top-level expressions, initial phase.]
\label{thm:preztlei}
If
\begin{enumerate}[label=(\roman*)]
\item $\Gamma \semi \Pi \semi \Pi[p] \vdtop \hat{e}$,
\item $\Gamma \semi \Pi \vdm m$ (with $p \in \dom(m)$), and
\item $m \semi \hat{e} \topstep{\mf{Init}}{p} \hat{e}' \semi m'$,
\end{enumerate}
then
\begin{enumerate}[label=(\roman*)]
\item $\Gamma \semi \Pi \semi \Pi[p] \vdtop \hat{e}'$,
\item $\Gamma \semi \Pi \vdm m'$, and
\item $\mc{L}(\hat{e}) \cup \dom(m[p].\mf{st}) 
= \mc{L}(\hat{e}') \cup \dom(m'[p].\mf{st})$.
\end{enumerate}

\begin{Proof}
By induction on derivation of the step judgment \emph{(iii)}.

\tb{Case} \ts{EH-Lift}: We have \begin{mathpar}
\inferrule*[right=EH-Lift] {
  \inferrule*[right=(1)] { } {m \semi e \mapsto e' \semi m}
} {
  m \semi e \topstep{\mf{Init}}{p} e' \semi m
}
\and
\inferrule*[right=TH-Lift] {
  \inferrule*[right=(2)] { } {\Gamma \semi \Pi \vdexp e: \mf{View} \mid \epsilon}
} {
  \Gamma \semi \Pi \semi \mc{C} \vdtop e
}
\end{mathpar}\begin{itemize}
\item By \thref{thm:prezexp} (preservation of expressions)
on (1), (2), and premise \emph{(ii)},
we get (3) $\Gamma \semi \Pi \vdexp e': \mf{View} \mid \epsilon$.
\item We satisfy conclusion \emph{(i)} via the following derivation:
\begin{mathpar}
\inferrule*[right=TH-Lift] {
  \inferrule*[right=(3)] { } {\Gamma \semi \Pi \vdexp e': \mf{View} \mid \epsilon}
} {
  \Gamma \semi \Pi \semi \mc{C} \vdtop e'
}
\end{mathpar}
\item Since $\mc{L}(e)=\mc{L}(e')=\emptyset$ and $m' = m$,
we also get $\mc{L}(e) \cup \dom(m[p].\mf{st}) 
= \mc{L}(e') \cup \dom(m[p].\mf{st})$, i.e. conclusion \emph{(iii)} holds.
\item Conclusion \emph{(ii)} is trivial, and we are done.
\end{itemize}

\tb{Case} \ts{EH-Let}: We have \begin{mathpar}
\inferrule*[right=EH-Let] {
  \inferrule*[right=(1)] { } {m \semi e \mapsto e' \semi m}
} {
  m \semi \elet{x}{e}{\hat{e}} \topstep{\mf{Init}}{p} \elet{x}{e'}{\hat{e}} \semi m
}
\and
\inferrule*[right=TH-Let] {
  \inferrule*[right=(2)] { } {\Gamma \semi \Pi \vdexp e: \tau \mid \epsilon}
  \and
  \inferrule*[right=(3)] { } {
    \Gamma[x \mapsto \tau] \semi \Pi \semi \mc{C} \vdtop \hat{e}
  }
} {
  \Gamma \semi \Pi \semi \mc{C} \vdtop \elet{x}{e}{\hat{e}}
}
\end{mathpar}\begin{itemize}
\item By \thref{thm:prezexp} (preservation of expressions)
on (1), (2), and premise \emph{(ii)},
we get that (4) $\Gamma \semi \Pi \vdexp e': \tau \mid \epsilon$.
\item We derive conclusion \emph{(i)} as follows:
\begin{mathpar}
\inferrule*[right=TH-Let] {
  \inferrule*[right=(4)] { } {\Gamma \semi \Pi \vdexp e': \tau \mid \epsilon}
  \and
  \inferrule*[right=(3)] { } {
    \Gamma[x \mapsto \tau] \semi \Pi \semi \mc{C} \vdtop \hat{e}
  }
} {
  \Gamma \semi \Pi \semi \mc{C} \vdtop \elet{x}{e}{\hat{e}}
}
\end{mathpar}
\item By definition, 
$\mc{L}(\elet{x}{e}{\hat{e}})=\mc{L}(\hat{e})=\mc{L}(\elet{x}{e'}{\hat{e}})$.
Since also $m'=m$, conclusion \emph{(iii)} holds.
\item Conclusion \emph{(ii)} is again trivial, and we are done.
\end{itemize}

\tb{Case} \ts{EH-Ret}:\begin{itemize}
\item We have
(1) $m \semi \elet{x}{\eret{v}}{\hat{e}} 
\topstep{\mf{Init}}{p} \hat{e}[v/x] \semi m$ 
and:
\begin{mathpar}
\inferrule*[right=TH-Let] {
  \inferrule*[right=TE-Ret] { 
    \inferrule*[right=(2)] { } {\Gamma \semi \Pi \vdexp v: \tau}
  } {
    \Gamma \semi \Pi \vdexp \eret{v}: \tau \mid \epsilon
  }
  \ 
  \inferrule*[right=(3)] { } {
    \Gamma[x \mapsto \tau] \semi \Pi \semi \mc{C} \vdtop \hat{e}
  }
} {
  \Gamma \semi \Pi \semi \mc{C} \vdtop \elet{x}{\eret{v}}{\hat{e}}
}
\end{mathpar}
\item By \thref{thm:subtle} (substitution for top-level expressions) on (2) and (3),
we get that $\Gamma \semi \Pi \semi \Pi[p] \vdtop \hat{e}[v/x]$,
satisfying conclusion \emph{(i)}.
\item By \thref{thm:sublab} (substitution for label-sets) and definition of $\mc{L}$,
we get $\mc{L}(\elet{x}{\eret{v}}{\hat{e}})=\mc{L}(\hat{e})=\mc{L}(\hat{e}[v/x])$.
Since also $m' = m$, conclusion \emph{(iii)} holds.
\item Conclusion \emph{(ii)} is once again trivial, and we are done.
\end{itemize}

\tb{Case} \ts{EH-StI}: \begin{itemize}
\item We have
$$m \semi \estate{x}{\ell}{v}{\hat{e}} \topstep{\mf{Init}}{p}
\hat{e}[v/x, \vset{\ell}{p}/x_\mf{set}] \semi m'$$
where $m' = m \mid m[p].\mf{st}[\ell] = v$, and:
\begin{mathpar}
\inferrule*[right=TH-State] {
  \Delta[\Pi[p]][\ell] = \tau_1\ (1)
  \and
  \Gamma \semi \Pi \vdval v: \tau_1\ (2)
  \\
  \Gamma[x \mapsto \tau_1, x_\mf{set} \mapsto \mf{Setter}\ \tau_1]
  \semi \Pi \semi \mc{C} \vdtop \hat{e}\ (3)
} {
  \Gamma \semi \Pi \semi \Pi[p] \vdtop \estate{x}{\ell}{v}{\hat{e}}
}
\end{mathpar}
\item We derive the setter type as follows:
\begin{mathpar}
\inferrule*[right=TV-Set] {
  \inferrule*[right=(1)] { } {\Delta[\Pi[p]][\ell] = \tau_1}
} {
  \Gamma \semi \Pi \vdval \vset{\ell}{p} : \mf{Setter}\ \tau_1\ (4)
}
\end{mathpar}
\item \thref{thm:subtle} (substitution for top-level expressions),
applied twice to (3) using (2) and (4),
yields $\Gamma \semi \Pi \semi \Pi[p] \vdtop \hat{e}[v/x, \vset{\ell}{p}/x_\mf{set}]$,
satisfying conclusion \emph{(i)}.
\item To show conclusion \emph{(ii)} given premise \emph{(ii)}, 
it suffices to show that
$\Gamma \semi \Pi \vdval m'[p].\mf{st}[\ell] : \Delta[\Pi[p]][\ell]$.
By construction of $m'$, we need to show
$\Gamma \semi \Pi \vdval v : \Delta[\Pi[p]][\ell]$.
This is given by (1) and (2).
\item Finally, we need to show conclusion \emph{(iii)}.
By definition of $\mc{L}$, 
$$\mc{L}(\estate{x}{\ell}{v}{\hat{e}}) = \{\ell\} \cup \mc{L}(\hat{e})$$
and by \thref{thm:sublab} (substitution for label-sets), 
$\mc{L}(\hat{e}[v/x, \vset{\ell}{p}/x_\mf{set}])=\mc{L}(\hat{e})$.
By construction, $\dom(m'[p].\mf{st}) = \dom(m[p].\mf{st}) \cup \{\ell\}$.
Conclusion \emph{(iii)} now follows, and we are done:
$$
(\{\ell\} \cup \mc{L}(\hat{e})) \cup \dom(m[p].\mf{st}) 
= \mc{L}(\hat{e}) \cup (\dom(m[p].\mf{st}) \cup \{\ell\})
$$
\end{itemize}

\tb{Case} \ts{EH-StS}: Impossible,
since the judgment requires that $\phi=\mf{Succ}$.
\end{Proof}
\end{lemma}

\begin{lemma}[Progress of top-level expressions, successor phase.]
\label{thm:progtles}
If \begin{enumerate}[label=(\roman*)]
\item $\cdot \semi \Pi \semi \Pi[p] \vdtop \hat{e}$,
\item $\cdot \semi \Pi \vdm m$, and
\item $\mc{L}(\hat{e}) \subseteq \dom(m[p].\mf{st})$,
\end{enumerate}
then either:
\begin{enumerate}[label=(\roman*)]
\item $\hat{e} = \eret{v}$, or
\item there exist some $\hat{e}'$ and $m'$ such that
$m \semi \hat{e} \topstep{\mf{Succ}}{p} \hat{e}' \semi m'$.
\end{enumerate}
\begin{Proof}
By induction over the derivation of the expression typing judgment.

\tb{Cases} \ts{TH-Lift}, \ts{TH-Let}:
Analogous to \thref{thm:progtlei}
(progress at initial phase).

\tb{Case} \ts{TH-State}: \begin{mathpar}
\inferrule*[right=TH-State] {
  \Delta[\mc{C}][\ell] = \tau_1
  \and
  \cdot \semi \Pi \vdval v: \tau_1
  \\
  [x \mapsto \tau_1, x_\mf{set} \mapsto \mf{Setter}\ \tau_1]
  \semi \Pi \semi \mc{C} \vdtop \hat{e}
} {
  \cdot \semi \Pi \semi \mc{C} \vdtop \estate{x}{\ell}{v}{\hat{e}}
}
\end{mathpar}
We want to apply step rule \ts{EH-StS}, which says:
$$m \semi \estate{x}{\ell}{v}{\hat{e}}
\topstep{\mf{Succ}}{p}
\hat{e}[v'/x, \vset{\ell}{p}/x_\mf{set}] \semi m$$
where $v' = m[p].\mf{st}[\ell]$.
It suffices to show that $v'$ is well-defined.

By premise \emph{(iii)} and definition of $\mc{L}$,
we must have that $\ell \in \dom(m[p].\mf{st})$.
\end{Proof}
\end{lemma}

\begin{lemma}[Preservation of top-level expressions, successor phase.]
\label{thm:preztles}
If
\begin{enumerate}[label=(\roman*)]
\item $\Gamma \semi \Pi \semi \Pi[p] \vdtop \hat{e}$,
\item $\Gamma \semi \Pi \vdm m$ (with $p \in \dom(m)$), and
\item $m \semi \hat{e} \topstep{\mf{Succ}}{p} \hat{e}' \semi m'$,
\end{enumerate}
then
\begin{enumerate}[label=(\roman*)]
\item $\Gamma \semi \Pi \semi \Pi[p] \vdtop \hat{e}'$,
\item $m = m'$, and
\item $\mc{L}(\hat{e}') \subseteq \mc{L}(\hat{e})$.
\end{enumerate}

\begin{Proof}
By induction on derivation of the step judgment \emph{(iii)}.

\tb{Case} \ts{EH-Lift}: We have \begin{mathpar}
\inferrule*[right=EH-Lift] {
  m \semi e \mapsto e' \semi m
} {
  m \semi e \topstep{\mf{Succ}}{p} e' \semi m
}
\end{mathpar}\begin{itemize}
\item Conclusion \emph{(i)} is shown analogously to \thref{thm:preztlei}
(preservation at initial phase),
and conclusion \emph{(ii)} is trivial.
\item By definition, $\mc{L}(e)=\emptyset=\mc{L}(e')$,
so conclusion \emph{(iii)} holds and we are done.
\end{itemize}

\tb{Case} \ts{EH-Let}: We have \begin{mathpar}
\inferrule*[right=EH-Let] {
  m \semi e \mapsto e' \semi m
} {
  m \semi \elet{x}{e}{\hat{e}} \topstep{\phi}{p} \elet{x}{e'}{\hat{e}} \semi m
}
\end{mathpar}\begin{itemize}
\item Again, conclusion \emph{(i)} is shown analogously to \thref{thm:preztlei}
and conclusion \emph{(ii)} is trivial.
\item By definition, 
$\mc{L}(\elet{x}{e}{\hat{e}})=\mc{L}(\hat{e})=\mc{L}(\elet{x}{e'}{\hat{e}})$,
so conclusion \emph{(iii)} holds and we are done.
\end{itemize}

\tb{Case} \ts{EH-Ret}: \begin{itemize}
\item We have $m \semi \elet{x}{\eret{v}}{\hat{e}} 
\topstep{\mf{Init}}{p} \hat{e}[v/x] \semi m$.
\item Once again, conclusion \emph{(i)} is shown analogously to \thref{thm:preztlei}
and conclusion \emph{(ii)} is trivial.
\item By definition,
$\mc{L}(\elet{x}{e}{\hat{e}})=\mc{L}(\hat{e})$,
and by \thref{thm:sublab} (substitution for label-sets),
$\mc{L}(\hat{e}[v/x])=\mc{L}(\hat{e})$.
Therefore, conclusion \emph{(iii)} holds and we are done.
\end{itemize}

\tb{Case} \ts{EH-StI}: Impossible,
since the judgment requires that $\phi=\mf{Init}$.

\tb{Case} \ts{EH-StS}: 
\begin{itemize}
\item We have
$$m \semi \estate{x}{\ell}{v}{\hat{e}}
\topstep{\mf{Succ}}{p}
\hat{e}[v'/x, \vset{\ell}{p}/x_\mf{set}] \semi m$$
where $v' = m[p].\mf{st}[\ell]$, and:
\begin{mathpar}
\inferrule*[right=TH-State] {
  \Delta[\Pi[p]][\ell] = \tau_1\ (1)
  \and
  \Gamma \semi \Pi \vdval v: \tau_1
  \\
  \Gamma[x \mapsto \tau_1, x_\mf{set} \mapsto \mf{Setter}\ \tau_1]
  \semi \Pi \semi \mc{C} \vdtop \hat{e}\ (2)
} {
  \Gamma \semi \Pi \semi \Pi[p] \vdtop \estate{x}{\ell}{v}{\hat{e}}
}
\end{mathpar}
\item As in \thref{thm:preztlei}, we derive the setter type as follows:
\begin{mathpar}
\inferrule*[right=TV-Set] {
  \inferrule*[right=(1)] { } {\Delta[\Pi[p]][\ell] = \tau_1}
} {
  \Gamma \semi \Pi \vdval \vset{\ell}{p} : \mf{Setter}\ \tau_1\ (3)
}
\end{mathpar}
\item By the memory typing premise \emph{(ii)},
we get that $\Gamma \semi \Pi \vdval v' : \Delta[\Pi[p]][\ell]$.
By (1), we get (4) that $\Gamma \semi \Pi \vdval v' : \tau_1$.
\item Now, we apply \thref{thm:subtle} (substitution for top-level expressions)
twice to (2) using (3) and (4) to get that 
$\Gamma \semi \Pi \semi \Pi[p] \vdtop \hat{e}[v'/x, \vset{\ell}{p}/x_\mf{set}]$,
satisfying conclusion \emph{(i)}.
\item Conclusion \emph{(iii)} follows from definition of $\mc{L}$:
clearly $\mc{L}(\hat{e}) \subseteq \{\ell\} \cup \mc{L}(\hat{e})$.
\item Conclusion \emph{(ii)} is trivial, and we are done.
\end{itemize}

\end{Proof}

\end{lemma}

Later, we will also require the following (easy) lemma:
\begin{lemma}[Top-level preservation of specifications.]
\label{thm:prezspec}
If $m \semi \hat{e} \topstep{\phi}{p} \hat{e}' \semi m'$,
then $\forall p \in \dom(m). m[p].\mf{spec} = m'[p].\mf{spec}$.

\begin{Proof}
By induction over the step judgment. Cases \ts{EH-Lift},
\ts{EH-Let}, \ts{EH-Ret}, and \ts{EH-StS} are all trivial,
as they all have $m'=m$.

In case \ts{EH-StI}, we have $m' = m \mid m[p].\mf{st}[\ell] = v$.
But this does not change any specs between $m'$ and $m'$,
so the statement still holds and we are done.
\end{Proof}
\end{lemma}

\subsubsection{Programs}

\begin{lemma}[Progress of programs.]
\label{thm:progprog}
If $\cdot \vdprog P$ and $\cdot \semi \Delta \vdd \mc{D}$,
then either:
\begin{enumerate}[label=(\roman*)]
\item $P = \eret{v}$, or
\item there exist some $P'$ and $\mc{D}'$ such that
$\mc{D} \semi P \pstep P' \semi \mc{D}'$.
\end{enumerate}
\begin{Proof}
By induction on the program typing derivation.

\tb{Case} \ts{TP-Lift}: \begin{mathpar}
\inferrule*[right=TP-Lift] {
  \inferrule*[right=(1)] { } {\Gamma \semi \cdot \vdexp e : \mf{View} \mid \epsilon}
} {
  \Gamma \vdprog e
}
\end{mathpar}
\begin{itemize}
\item By definition, the empty memory is well-typed 
by the empty component context, i.e. (2) $\cdot \semi \cdot \vdm \cdot$.
\item We apply \thref{thm:progexp} (progress of expressions)
on (1) and the empty memory
to get that either (a) $e = \eret{v}$ or (b) $\cdot \semi e \mapsto e' \semi m$.
\item In subcase (a), we are done by conclusion \emph{(i)}.
\item In subcase (b), we note by \thref{thm:pure} (purity) that $m=\cdot$.
We derive a step as follows, satisfying conclusion \emph{(ii)}:
\end{itemize}
\begin{mathpar}
\inferrule*[right=EP-Lift] {
  \cdot \semi e \mapsto e' \semi \cdot
} {
  \mc{D} \semi e \pstep e' \semi \mc{D}
}
\end{mathpar}

\tb{Case} \ts{TP-Let}: \begin{mathpar}
\inferrule*[right=TP-Let] {
  \inferrule*[right=(1)] { } {
    \cdot \semi \cdot \vdexp e: \tau \mid \epsilon
  }
  \and 
  [x \mapsto \tau] \vdprog P
} {
  \cdot \vdprog \elet{x}{e}{P}
}
\end{mathpar}
\begin{itemize}
\item We again apply \thref{thm:progexp}
on (1) and the empty memory
to get that either (a) $e = \eret{v}$ or (b) $\cdot \semi e \mapsto e' \semi m$.
\begin{itemize}
\item Subcase (a): \ts{EP-Ret} gives us that
$\mc{D} \semi \elet{x}{\eret{v}}{P} \pstep P[v/x] \semi \mc{D}$.
\item Subcase (b): by \thref{thm:pure} (purity) on (1), $m' = \cdot$,
so we derive a step as follows:
\begin{mathpar}
\inferrule*[right=EP-Let] {
  \cdot \semi e \mapsto e' \semi \cdot
} {
  \mc{D} \semi \elet{x}{e}{P}
  \pstep
  \elet{x}{e'}{P} \semi \mc{D}
}
\end{mathpar}
\end{itemize}
\end{itemize}

\tb{Case} \ts{TP-Def}: \begin{mathpar}
\inferrule*[right=TP-Def] {
  [x \mapsto \Delta[\mc{C}]] \semi \cdot \semi \mc{C} \vdtop \hat{e}
  \and
  \cdot \vdprog P
} {
  \cdot \vdprog \elet{\mc{C}(x)}{\hat{e}}{P}
}
\end{mathpar}
Rule \ts{EP-Def} immediately applies,
giving us $$\mc{D} \semi \elet{\mc{C}(x)}{\hat{e}}{P}
\pstep
P \semi \mc{D}[\mc{C} \mapsto x.\hat{e}]$$
and we are done.
\end{Proof}
\end{lemma}

\begin{lemma}[Preservation of programs.]
\label{thm:prezprog}
If
\begin{enumerate}[label=(\roman*)]
\item $\Gamma \vdprog P$,
\item $\Gamma \semi \Delta \vdd \mc{D}$, and
\item $\mc{D} \semi P \pstep P' \semi \mc{D}'$,
\end{enumerate}
then
\begin{enumerate}[label=(\roman*)]
\item $\Gamma \vdprog P'$,
\item $\Gamma \semi \Delta \vdd \mc{D}'$, and
\item $\mc{L}(P') \cup \dom(\mc{D}') = \mc{L}(P) \cup \dom(\mc{D})$.
\end{enumerate}

\begin{Proof}
By induction on the derivation of the step judgment \emph{(iii)}.

\tb{Case} \ts{EP-Lift}: We have \begin{mathpar}
\inferrule*[right=EP-Lift] {
  \inferrule*[right=(1)] { } {\cdot \semi e \mapsto e' \semi \cdot}
} {
  \mc{D} \semi e \pstep e' \semi \mc{D}
}
\and
\inferrule*[right=TP-Lift] {
  \inferrule*[right=(2)] { } {\Gamma \semi \cdot \vdexp e : \mf{View} \mid \epsilon}
} {
  \Gamma \vdprog e
}
\end{mathpar}
\begin{itemize}
\item \thref{thm:prezexp} (preservation of expressions)
on (1), (2), and the empty memory gives us (3) that 
$\Gamma \semi \cdot \vdexp e': \mf{View} \mid \epsilon$.
\item We derive conclusion \emph{(i)} as follows:
\begin{mathpar}
\inferrule*[right=TP-Lift] {
  \inferrule*[right=(3)] { } {\Gamma \semi \cdot \vdexp e': \mf{View} \mid \epsilon}
} {
  \Gamma \vdprog e'
}
\end{mathpar}
\item Conclusion \emph{(ii)} is trivial. Conclusion \emph{(iii)}
follows from definition of $\mc{L}$: we have $\mc{L}(e)=\emptyset=\mc{L}(e')$.
\end{itemize}

\tb{Case} \ts{EP-Let}: We have \begin{mathpar}
\inferrule*[right=EP-Let] {
  \inferrule*[right=(1)] { } {\cdot \semi e \mapsto e' \semi \cdot}
} {
  \mc{D} \semi \elet{x}{e}{P}
  \pstep
  \elet{x}{e'}{P} \semi \mc{D}
}
\and
\inferrule*[right=TP-Let] {
  \inferrule*[right=(2)] { } {\Gamma \semi \cdot \vdexp e: \tau \mid \epsilon}
  \and 
  \inferrule*[right=(3)] { } {\Gamma[x \mapsto \tau] \vdprog P}
} {
  \Gamma \vdprog \elet{x}{e}{P}
}
\end{mathpar}
\begin{itemize}
\item \thref{thm:prezexp}
on (1), (2), and the empty memory gives us (4) that 
$\Gamma \semi \cdot \vdexp e': \tau \mid \epsilon$.
\item We derive conclusion \emph{(i)} as follows:
\begin{mathpar}
\inferrule*[right=TP-Let] {
  \inferrule*[right=(4)] { } {\Gamma \semi \cdot \vdexp e': \tau \mid \epsilon}
  \and 
  \inferrule*[right=(3)] { } {\Gamma[x \mapsto \tau] \vdprog P}
} {
  \Gamma \vdprog \elet{x}{e'}{P}
}
\end{mathpar}
\item Conclusion \emph{(ii)} is again trivial. Conclusion \emph{(iii)}
follows from definition of $\mc{L}$: we have 
$\mc{L}(\elet{x}{e}{P})=\mc{L}(P)=\mc{L}(\elet{x}{e'}{P})$.
\end{itemize}

\tb{Case} \ts{EP-Ret}: We have
$\mc{D} \semi \elet{x}{\eret{v}}{P} \pstep P[v/x] \semi \mc{D}$ and 
\begin{mathpar}
\inferrule*[right=TP-Let] {
  \inferrule*[right=TE-Ret] {
    \inferrule*[right=(1)] { } {
      \Gamma \semi \cdot \vdval v: \tau
    }
  } {
    \Gamma \semi \cdot \vdexp \eret{v}: \tau \mid \epsilon
  }
  \and 
  \inferrule*[right=(2)] { } {\Gamma[x \mapsto \tau] \vdprog P}
} {
  \Gamma \vdprog \elet{x}{\eret{v}}{P}
}
\end{mathpar}
\begin{itemize}
\item \thref{thm:subprog} (substitution for programs) on (1) and (2)
gives conclusion \emph{(i)}, that $\Gamma \vdprog P[v/x]$.
\item For conclusion \emph{(iii)}, we have
$\mc{L}(\elet{x}{\eret{v}}{P})=\mc{L}(P)$ by definition.
We get that $\mc{L}(P)=\mc{L}(P[v/x])$
from \thref{thm:subname} (component-name substitution).
\item Conclusion \emph{(ii)} is once again trivial, and we are done.
\end{itemize}

\tb{Case} \ts{EP-Def}: We have
$\mc{D} \semi \elet{\mc{C}(x)}{\hat{e}}{P}
\pstep
P \semi \mc{D}[\mc{C} \mapsto x.\hat{e}]$ and
\begin{mathpar}
\inferrule*[right=TP-Def] {
  \inferrule*[right=(1)] { } {
    \Gamma[x \mapsto \Delta[\mc{C}]] \semi \cdot \semi \mc{C} \vdtop \hat{e}
  }
  \and
  \inferrule*[right=(2)] { } {\Gamma \vdprog P}
} {
  \Gamma \vdprog \elet{\mc{C}(x)}{\hat{e}}{P}
}
\end{mathpar}
\begin{itemize}
\item Conclusion \emph{(i)} is immediate from (2).
\item To show conclusion \emph{(ii)} given premise \emph{(ii)},
it suffices to show that 
$\Gamma[x \mapsto \Delta[\mc{C}]] \semi \cdot \semi \mc{C} \vdtop \hat{e}$.
This is immediate from (1).
\item For conclusion \emph{(iii)},
we have that $\mc{L}(\elet{\mc{C}(x)}{\hat{e}}{P}) = \{\mc{C}\} \cup \mc{L}(P)$
by definition, and that
$\dom(\mc{D}') = \{\mc{C}\} \cup \dom(\mc{D})$ by construction of $\mc{D}'$.
Conclusion \emph{(iii)} follows, and we are done:
$$\mc{L}(P) \cup (\{\mc{C}\} \cup \dom(\mc{D})) 
= 
(\{\mc{C}\} \cup \mc{L}(P)) \cup \dom(\mc{D})$$
\end{itemize}
\end{Proof}
\end{lemma}

\pagebreak

\section{Small-Step Runtime}

\subsection{Runtime Syntax}

\begin{align*}
\text{Runtime values } u
&::= \mf{done} \mid []_\rho \mid u ::_\rho u
\mid \mf{lit}\ v \mid \mf{node}(\mf{id}, v, u)
\mid p\\
\text{Runtime expressions } \rho
&::= \mf{done} \mid []_\rho \mid \rho ::_\rho \rho
\mid \mf{lit}\ v \mid \mf{node}(\mf{id}, v, \rho)
\mid p\\
&\quad\mid \mf{init}\ v \mid \mf{inits}\ v \mid \mf{scan}\ v \mid \mf{scans}\ v\\
&\quad\mid \mf{check}\ \rho \mid \mf{checks}\ \rho 
\mid \mf{search}_\mf{id}\ \rho \mid \mf{searches}_\mf{id}\ \rho\\
&\quad\mid \mf{eval}^\phi_p\ \hat{e} \mid \mf{new}^\phi_p\ \rho
\mid \mf{handle}\ e \mid \mf{prog}\ P\\
&\quad\mid \rho \semi \rho\\
\text{Evaluation contexts } \mc{E}
&::= [\cdot] \mid \mf{node}(\mf{id}, v, \mc{E}) \mid \mc{E} ::_\rho \rho
\mid u ::_\rho \mc{E} \mid \mf{new}^\phi_p\ \mc{E}\\
&\quad\mid \mf{check}\ \mc{E} \mid \mf{checks}\ \mc{E} 
\mid \mf{search}_\mf{id}\ \mc{E} \mid \mf{searches}_\mf{id}\ \mc{E}\\
&\quad\mid \mc{E} \semi \rho\\
\text{Runtime type } R &::= \one_\rho \mid \mf{Tree}
\mid [R]_\rho
\end{align*}

We associate each phase $\phi$ with a runtime type as follows:
\[
R(\phi) = \begin{cases}
    \mf{Tree} & \phi=\mf{Init}\\
    \one_\rho & \phi=\mf{Succ}
\end{cases}
\]

\tb{\emph{Constants.}} We distinguish three constants:
\begin{itemize}
\item $p_0$, the "root path" in memory,
\item $\mc{C}_0$, the (dummy) component name associated with that path, and
\item $x_0$, the variable name bound in that component's body.
\end{itemize}
We assume these constants are \emph{globally fresh},
i.e. they do not conflict with any other names in the program
or that are ``generated'' at runtime, e.g. by the premise ``$p$ fresh''.

\subsection{Runtime Typing Rules}

Judgment: $\Gamma \semi \Pi \vdrt \rho: R$
\begin{mathpar}
\inferrule*[right=TR-Done] { } {
  \Gamma \semi \Pi \vdrt \mf{done}: \one_\rho
}
\and
\inferrule*[right=TR-Nil] { } {
  \Gamma \semi \Pi \vdrt []_\rho: [R]_\rho
}
\and
\inferrule*[right=TR-Cons] {
  \Gamma \semi \Pi \vdrt \rho_\mf{hd} : R
  \and
  \Gamma \semi \Pi \vdrt \rho_\mf{tl} : [R]_\rho
} {
  \Gamma \semi \Pi \vdrt \rho_\mf{hd} :: \rho_\mf{tl} : [R]_\rho
}
\and
\inferrule*[right=TR-Lit] {
  \Gamma \semi \Pi \vdval v: \mf{String}
} {
  \Gamma \semi \Pi \vdrt \mf{lit}\ v: \mf{Tree}
}
\and
\inferrule*[right=TR-Node] {
  \Gamma \semi \Pi \vdval v: [\mf{Attr}]
  \and
  \Gamma \semi \Pi \vdrt \rho: [\mf{Tree}]_\rho
} {
  \Gamma \semi \Pi \vdrt \mf{node}(\mf{id}, v, \rho): \mf{Tree}
}
\and
\inferrule*[right=TR-Path] {
  p \in \dom(\Pi) \and \Pi[p] \in \dom(\Delta)
} {
  \Gamma \semi \Pi \vdrt p: \mf{Tree}
}
\and
\inferrule*[right=TR-Init] {
  \Gamma \semi \Pi \vdval v: \mf{View}
} {
  \Gamma \semi \Pi \vdrt \mf{init}\ v: \mf{Tree}
}
\and
\inferrule*[right=TR-Inits] {
  \Gamma \semi \Pi \vdval v: [\mf{View}]
} {
  \Gamma \semi \Pi \vdrt \mf{inits}\ v: [\mf{Tree}]_\rho
}
\and
\inferrule*[right=TR-Sc] {
  \Gamma \semi \Pi \vdval v: \mf{Attr}
} {
  \Gamma \semi \Pi \vdrt \mf{scan}\ v: \one_\rho
}
\and
\inferrule*[right=TR-Scs] {
  \Gamma \semi \Pi \vdval v: [\mf{Attr}]
} {
  \Gamma \semi \Pi \vdrt \mf{scans}\ v: \one_\rho
}
\and
\inferrule*[right=TR-Chk] {
  \Gamma \semi \Pi \vdrt \rho: \mf{Tree}
} {
  \Gamma \semi \Pi \vdrt \mf{check}\ \rho: \one_\rho
}
\and
\inferrule*[right=TR-Chks] {
  \Gamma \semi \Pi \vdrt \rho: [\mf{Tree}]_\rho
} {
  \Gamma \semi \Pi \vdrt \mf{checks}\ \rho: \one_\rho
}
\and
\inferrule*[right=TR-Srch] {
  \Gamma \semi \Pi \vdrt \rho: \mf{Tree}
} {
  \Gamma \semi \Pi \vdrt \mf{search}\ \rho: \one_\rho
}
\and
\inferrule*[right=TR-Srchs] {
  \Gamma \semi \Pi \vdrt \rho: [\mf{Tree}]_\rho
} {
  \Gamma \semi \Pi \vdrt \mf{searches}\ \rho: \one_\rho
}
\and
\inferrule*[right=TR-Eval] {
  \Gamma \semi \Pi \semi \Pi[p] \vdtop \hat{e}
} {
  \Gamma \semi \Pi \vdrt \mf{eval}^\phi_p\ \hat{e}: R(\phi)
}
\and
\inferrule*[right=TR-New] {
  \Gamma \semi \Pi \vdrt \rho: \mf{Tree}
} {
  \Gamma \semi \Pi \vdrt \mf{new}^\phi_p\ \rho: R(\phi)
}
\and
\inferrule*[right=TR-Hdl] {
  \Gamma \semi \Pi \vdexp e: \one \mid E
} {
  \Gamma \semi \Pi \vdrt \mf{handle}\ e: \one_\rho
}
\and
\inferrule*[right=TR-Prog] {
  \Gamma \semi \Pi \vdprog P
  \and
  \Delta[\mc{C}_0] = \one
} {
  \Gamma \semi \Pi \vdrt \mf{prog}\ P: \one_\rho
}
\and
\inferrule*[right=TR-Seq] {
  \Gamma \semi \Pi \vdrt \rho_1: \one_\rho
  \and
  \Gamma \semi \Pi \vdrt \rho_2: R
} {
  \Gamma \semi \Pi \vdrt \rho_1 \semi \rho_2: R
}
\end{mathpar}

Define $\Gamma \semi \Pi \vdcfg \trip{\mc{D}}{m}{\rho}: R$
to abbreviate that
(i) $\Gamma \semi \Delta \vdd \mc{D}$,
(ii) $\Gamma \semi \Pi \vdm m$, and
(iii) $\Gamma \semi \Pi \vdrt \rho: R$.

\subsection{Well-Formedness of Configurations}

In addition to the typing judgment,
we require a \emph{well-formedness} judgment on
configurations, written $\mc{D} \semi m \wf \rho$,
that describes which components have been defined,
and which labels and paths have been initialized.

\subsubsection{Component Definitions}
\begin{mathpar}
\inferrule*[right=WF-Prog] {
  \mc{L}(P) \cup \dom(\mc{D}) \cup \{\mc{C}_0\}= \dom(\Delta)
} {
  \mc{D} \semi m \wf \mf{prog}\ P
}
\end{mathpar}
\emph{All other rules} have the implicit premise that $\dom(\mc{D}) = \dom(\Delta)$.

\subsubsection{Label/Path Initialization}

Define, as abbreviations in the \emph{premises} of all rules
in this section, $m[p].\mf{spec} = \pair{\mc{C}}{v}$ and $\mc{D}[\mc{C}] = x.\hat{e}_0$.
(Thus, as implicit premises, we must have $m[p]$ and $\mc{D}[\mc{C}]$ defined.)
\begin{mathpar}
\inferrule*[right=WF-Path] {
  \mc{L}(\hat{e}_0) = \dom(m[p].\mf{st})
  \and
  \mc{D} \semi m \wf m[p].\mf{tree}
} {
  \mc{D} \semi m \wf p
}
\and
\inferrule*[right=WF-EvalI] {
  \mc{L}(\hat{e}) \cup \dom(m[p].\mf{st}) = \mc{L}(\hat{e}_0)
} {
  \mc{D} \semi m \wf \mf{eval}^\mf{Init}_p\ \hat{e}
}
\and
\inferrule*[right=WF-EvalS] {
  \mc{L}(\hat{e}_0) = \dom(m[p].\mf{st})
  \and
  \mc{L}(\hat{e}) \subseteq \dom(m[p].\mf{st})
} {
  \mc{D} \semi m \wf \mf{eval}^\mf{Succ}_p\ \hat{e}
}
\and
\inferrule*[right=WF-New] {
  \mc{L}(\hat{e}_0) = \dom(m[p].\mf{st})
  \and
  \mc{D} \semi m \wf \rho
} {
  \mc{D} \semi m \wf \mf{new}^\phi_p\ \rho
}
\end{mathpar}

\subsubsection{Boring Rules}

\begin{align*}
\mc{D} \semi m &\wf \mf{done} &&\ts{WF-Done}\\
\mc{D} \semi m &\wf []_\rho &&\ts{WF-Nil}\\
\mc{D} \semi m &\wf \mf{lit}\ v &&\ts{WF-Lit}\\
\mc{D} \semi m &\wf \mf{init}\ v &&\ts{WF-Init}\\
\mc{D} \semi m &\wf \mf{inits}\ v &&\ts{WF-Inits}\\
\mc{D} \semi m &\wf \mf{scan}\ v &&\ts{WF-Scan}\\
\mc{D} \semi m &\wf \mf{scans}\ v &&\ts{WF-Scans}\\
\mc{D} \semi m &\wf \mf{handle}\ e &&\ts{WF-Hdl}\\
\end{align*}

\begin{mathpar}
\inferrule*[right=WF-Cons] {
  \mc{D} \semi m \wf \rho_\mf{hd}
  \and
  \mc{D} \semi m \wf \rho_\mf{tl}
} {
  \mc{D} \semi m \wf \rho_\mf{hd} :: \rho_\mf{tl}
}
\and
\inferrule*[right=WF-Node] {
  \mc{D} \semi m \wf \rho
} {
  \mc{D} \semi m \wf \mf{node}(\mf{id}, v, \rho)
}
\and
\inferrule*[right=WF-Chk] {
  \mc{D} \semi m \wf \rho
} {
  \mc{D} \semi m \wf \mf{check}\ \rho
}
\and
\inferrule*[right=WF-Chks] {
  \mc{D} \semi m \wf \rho
} {
  \mc{D} \semi m \wf \mf{checks}\ \rho
}
\and
\inferrule*[right=WF-Srch] {
  \mc{D} \semi m \wf \rho
} {
  \mc{D} \semi m \wf \mf{search}_\mf{id}\ \rho
}
\and
\inferrule*[right=WF-Srchs] {
  \mc{D} \semi m \wf \rho
} {
  \mc{D} \semi m \wf \mf{searches}_\mf{id}\ \rho
}
\and
\inferrule*[right=WF-Seq] {
  \mc{D} \semi m \wf \rho_1
  \and
  \mc{D} \semi m \wf \rho_2
} {
  \mc{D} \semi m \wf \rho_1 \semi \rho_2
}
\end{mathpar}

\subsection{Runtime-Level Semantics}

Judgment: $\trip{\mc{D}}{m}{\rho} \cstep \trip{\mc{D}}{m}{\rho}$

(The initial configuration is $\trip{\cdot}{\cdot}{\mf{prog}\ P}$.)
\subsubsection{General Rules}

\begin{mathpar}
\inferrule*[right=ER-Ctx] {
  \trip{\mc{D}}{m}{\rho} \cstep \trip{\mc{D}'}{m'}{\rho'}
} {
  \trip{\mc{D}}{m}{\mc{E}[\rho]}
  \cstep 
  \trip{\mc{D}'}{m'}{\mc{E}[\rho']}
}
\end{mathpar}
\begin{align*}
\trip{\mc{D}}{m}{\mf{done} \semi \rho}
&\cstep
\trip{\mc{D}}{m}{\rho}
&&\ts{ER-Seq}
\end{align*}

\subsubsection{Program Evaluation}

\begin{mathpar}
\inferrule*[right=ER-Prog] {
  \mc{D} \semi P \pstep P' \semi \mc{D}'
} {
  \trip{\mc{D}}{m}{\mf{prog}\ P}
  \cstep
  \trip{\mc{D}'}{m}{\mf{prog}\ P'}
}
\and
\inferrule*[right=ER-PDone] {
  \mc{D}' = \mc{D}[\mc{C}_0 \mapsto x.\eret{v}]
  \\
  m' = m[p_0 \mapsto \{\mf{spec}: \pair{\mc{C}_0}{()}, \mf{st}: \cdot,
  \mf{tree}: \cdot, \mf{rerender}: \mf{ff}\}]
} {
  \trip{\mc{D}}{m}{\mf{prog\ \eret{v}}} 
  \cstep
  \trip{\mc{D}'}{m'}{\mf{new}^\mf{Succ}_{p_0}\ \mf{init}\ v}
}
\end{mathpar}

\subsubsection{Initialization}
\begin{mathpar}
\inferrule*[right=ER-ISpec] {
  p\ \mf{fresh}
  \and 
  \mc{D}[\mc{C}]=x.\hat{e}
  \\
  \pi_0=\{
    \mf{spec}=\pair{\mc{C}}{v}, 
    \mf{st}: \cdot, 
    \mf{tree}: \cdot, 
    \mf{rerender}: \mf{ff}
  \}
} {
  \trip{\mc{D}}{m}{\mf{init}\ \pair{\mc{C}}{v}}
  \cstep
  \trip{\mc{D}}{m[p \mapsto \pi_0]}{\mf{eval}^\mf{Init}_p\ \hat{e}[v/x]}
}
\end{mathpar}
\begin{align*}
\trip{\mc{D}}{m}{\mf{init\ text}\ v}
&\cstep
\trip{\mc{D}}{m}{\mf{lit}\ v}
&&\ts{ER-IText}\\
\trip{\mc{D}}{m}{\mf{init\ tag}(v_1, v_2)}
&\cstep
\trip{\mc{D}}{m}{\mf{node}(\mf{id}, v_1, \mf{inits}\ v_2)}
&&\ts{ER-ITag}\\
&(\mf{id\ fresh})\\
\trip{\mc{D}}{m}{\mf{inits}\ []}
&\cstep
\trip{\mc{D}}{m}{[]_\rho}
&&\ts{ER-INil}\\
\trip{\mc{D}}{m}{\mf{inits}\ (v :: vs)}
&\cstep
\trip{\mc{D}}{m}{(\mf{init}\ v) ::_\rho (\mf{inits}\ vs)}
&&\ts{ER-ICons}\\
\end{align*}

\subsubsection{Component Evaluation}
\begin{mathpar}
\inferrule*[right=ER-Eval] {
  m \semi \hat{e} \topstep{\phi}{p} \hat{e}' \semi m'
} {
  \trip{\mc{D}}{m}{\mf{eval}^\phi_p\ \hat{e}}
  \cstep
  \trip{\mc{D}}{m'}{\mf{eval}^\phi_p\ \hat{e}'}
}
\end{mathpar}
\begin{align*}
\trip{\mc{D}}{m}{\mf{eval}^\phi_p\ \eret{v}}
&\cstep
\trip{\mc{D}}{m}{\mf{new}^\phi_p\ (\mf{init}\ v)}
&&\ts{ER-New}\\
\trip{\mc{D}}{m}{\mf{new}^\mf{Init}_p\ u}
&\cstep
\trip{\mc{D}}{m \mid m[p].\mf{tree} = u}{p}
&&\ts{ER-NewI}\\
\trip{\mc{D}}{m}{\mf{new}^\mf{Succ}_p\ u}
&\cstep
\trip{\mc{D}}{m \mid m[p].\mf{tree} = u}{\mf{done}}
&&\ts{ER-NewS}\\
\end{align*}

\subsubsection{Handler Search/Evaluation}
\begin{align*}
\trip{\mc{D}}{m}{\mf{search}_\mf{id}\ \mf{lit}\ v}
&\cstep
\trip{\mc{D}}{m}{\mf{done}}
&&\ts{ER-SLit}\\
\trip{\mc{D}}{m}{\mf{search}_\mf{id}\ p}
&\cstep
\trip{\mc{D}}{m}{\mf{search}_\mf{id}\ m[p].\mf{tree}}
&&\ts{ER-SPath}\\
\trip{\mc{D}}{m}{\mf{search}_\mf{id}\ \mf{node}(\mf{id}', v, \rho)}
&\cstep
\trip{\mc{D}}{m}{\mf{searches}_\mf{id}\ \rho}
&&\ts{ER-SNodeF}\\
\trip{\mc{D}}{m}{\mf{search}_\mf{id}\ \mf{node}(\mf{id}, v, \rho)}
&\cstep
\trip{\mc{D}}{m}{\mf{scans}\ v \semi \mf{searches}_\mf{id}\ \rho}
&&\ts{ER-SNodeT}\\
\trip{\mc{D}}{m}{\mf{searches}_\mf{id}\ []_\rho}
&\cstep
\trip{\mc{D}}{m}{\mf{done}}
&&\ts{ER-SsNil}\\
\trip{\mc{D}}{m}{\mf{searches}_\mf{id}\ (\rho_\mf{1} ::_\rho \rho_\mf{2})}
\cstep&
\trip{\mc{D}}{m}{\mf{search}_\mf{id}\ \rho_\mf{1} \semi \mf{searches}_\mf{id}\ \rho_\mf{2}}
&&\ts{ER-SsCons}\\
\end{align*}
\begin{align*}
\trip{\mc{D}}{m}{\mf{scan\ attr}(v, v)}
&\cstep
\trip{\mc{D}}{m}{\mf{done}}
&&\ts{ER-ScAttr}\\
\trip{\mc{D}}{m}{\mf{scan\ onClick}(v)}
&\cstep
\trip{\mc{D}}{m}{\mf{handle}(v\ ())}
&&\ts{ER-ScClick}\\
\trip{\mc{D}}{m}{\mf{scans}\ []}
&\cstep
\trip{\mc{D}}{m}{\mf{done}}
&&\ts{ER-ScsNil}\\
\trip{\mc{D}}{m}{\mf{scans}\ (v :: vs)}
&\cstep
\trip{\mc{D}}{m}{\mf{scan}\ v \semi \mf{scans}\ vs}
&&\ts{ER-ScsCons}\\
\end{align*}
\begin{mathpar}
\inferrule*[right=ER-Hdl] {
  m \semi e \mapsto e' \semi m'
} {
  \trip{\mc{D}}{m}{\mf{handle}\ e}
  \cstep
  \trip{\mc{D}}{m'}{\mf{handle}\ e'}
}
\end{mathpar}
\begin{align*}
\trip{\mc{D}}{m}{\mf{handle}\ \eret{()}}
&\cstep
\trip{\mc{D}}{m}{\mf{done}}
&&\ts{ER-HdlDone}\\
\end{align*}

\subsubsection{Checking}

\begin{mathpar}
\inferrule*[right=ER-ChkPF] {
  m[p].\mf{rerender} = \mf{ff}
} {
  \trip{\mc{D}}{m}{\mf{check}\ p}
  \cstep
  \trip{\mc{D}}{m}{\mf{check}\ m[p].\mf{tree}}
}
\and
\inferrule*[right=ER-ChkPT] {
  m[p].\mf{rerender} = \mf{tt}
} {
  \trip{\mc{D}}{m}{\mf{check}\ p}
  \cstep
  \trip{\mc{D}}{m \mid m[p].\mf{rerender} = \mf{ff}}{\mf{eval}^\mf{Succ}_p\ \hat{e}_0[v/x]}
}
\end{mathpar}
\begin{align*}
\trip{\mc{D}}{m}{\mf{check\ lit}\ v}
&\cstep
\trip{\mc{D}}{m}{\mf{done}}
&&\ts{ER-ChkLit}\\
\trip{\mc{D}}{m}{\mf{check\ node}(\mf{id}, v, \rho)}
&\cstep
\trip{\mc{D}}{m}{\mf{checks}\ \rho}
&&\ts{ER-ChkNode}\\
\trip{\mc{D}}{m}{\mf{checks}\ []_\rho}
&\cstep
\trip{\mc{D}}{m}{\mf{done}}
&&\ts{ER-ChksNil}\\
\trip{\mc{D}}{m}{\mf{checks}\ \rho_1 ::_\rho \rho_2}
&\cstep
\trip{\mc{D}}{m}{\mf{check}\ \rho_1 \semi \mf{checks}\ \rho_2}
&&\ts{ER-ChksCons}
\end{align*}

\subsubsection{Events}

We introduce a second step judgment, $\cstep_0$,
which subsumes $\cstep$ and additionally allows for
nondeterministic event occurrence at idle state.
(We distinguish $\cstep_0$ to prevent events from occuring inside nested terms.)
\begin{mathpar}
\inferrule*[right=ER-Top] {
  \trip{\mc{D}}{m}{\rho} \cstep \trip{\mc{D}'}{m'}{\rho'}
} {
  \trip{\mc{D}}{m}{\rho} \cstep_0 \trip{\mc{D}'}{m'}{\rho'}
}
\end{mathpar}
\begin{align*}
\trip{\mc{D}}{m}{\mf{done}}
&\cstep_0
\trip{\mc{D}}{m}{\mf{search}_\mf{id}\ p_0 \semi \mf{check}\ p_0}
&&\ts{ER-Event}\\
\end{align*}

\subsection{Configuration Soundness}

% \begin{lemma}[Progress of configurations.]
% \label{thm:progcfg}
% If $\cdot \semi \Pi \vdcfg \trip{\mc{D}}{m}{\rho}: R$
% and also $\mc{D} \semi m \wf \rho$,
% then either:
% \begin{enumerate}[(i)]
% \item $\rho = u$ (i.e. $\rho$ is a value), or
% \item there exist some $\mc{D}'$, $m'$, and $\rho'$ such that
% $\trip{\mc{D}}{m}{\rho}\cstep \trip{\mc{D}'}{m'}{\rho'}$.
% \end{enumerate}
% \begin{Proof}
% By definition of $\vdcfg$, we have sub-premises:
% \begin{enumerate}[(i)]
% \item $\Delta \vdd \mc{D}$,
% \item $\cdot \semi \Pi \vdm m$, and
% \item $\cdot \semi \Pi \vdrt \rho: R$.
% \end{enumerate}
% We proceed by induction over the derivation of \emph{(iii)}.

% \tb{Case} 
% \ts{TR-Done}, 
% \ts{TR-Nil},
% \ts{TR-Lit},
% \ts{TR-Path}: 
% Trivial by conclusion \emph{(i)}. Respectively,
% $\mf{done}$, $[]_\rho$, $\mf{lit}\ v$ (for any $v$), and $p$
% are values.

% \tb{Case} \ts{TR-Cons}: \begin{mathpar}
% \inferrule*[right=TR-Cons] {
%   \inferrule*[right=(1)] { } {\cdot \semi \Pi \vdrt \rho_\mf{hd} : R}
%   \and
%   \inferrule*[right=(2)] { } {\cdot \semi \Pi \vdrt \rho_\mf{tl} : [R]}
% } {
%   \cdot \semi \Pi \vdrt \rho_\mf{hd} :: \rho_\mf{tl} : [R]
% }
% \and
% \inferrule*[right=WF-Cons] {
%   \inferrule*[right=(3)] { } {\mc{D} \semi m \wf \rho_\mf{hd}}
%   \and
%   \inferrule*[right=(4)] { } {\mc{D} \semi m \wf \rho_\mf{tl}}
% } {
%   \mc{D} \semi m \wf \rho_\mf{hd} :: \rho_\mf{tl}
% }
% \end{mathpar}
% \begin{itemize}
% \item By \emph{(i)}, \emph{(ii)}, and (1), we have 
% (5) $\cdot \semi \Pi \vdcfg \trip{\mc{D}}{m}{\rho_\mf{hd}}$.
% \item Inductive hypothesis on (3) and (5) gives either
% (6a) that $\trip{\mc{D}}{m}{\rho_\mf{hd}}\cstep \trip{\mc{D}'}{m'}{\rho'_\mf{hd}}$
% or (6b) that $\rho_\mf{hd} = u_\mf{hd}$.
% \begin{itemize}
% \item In subcase (6a), we apply \ts{ER-Ctx} with $\mc{E}=[\cdot] ::_\rho \rho_\mf{tl}$,
% satisfying conclusion \emph{(ii)}:
% \begin{mathpar}
% \inferrule*[right=ER-Ctx] {
%   \inferrule*[right=\tn{(6a)}] { } {
%     \trip{\mc{D}}{m}{\rho_\mf{hd}} \cstep \trip{\mc{D}'}{m'}{\rho_\mf{hd}'}
%   }
% } {
%   \trip{\mc{D}}{m}{\rho_\mf{hd} ::_\rho \rho_\mf{tl}}
%   \cstep 
%   \trip{\mc{D}'}{m'}{\rho_\mf{hd}' ::_\rho \rho_\mf{tl}}
% }
% \end{mathpar}
% \end{itemize}
% \item Otherwise, we get (7) that 
% $\cdot \semi \Pi \vdcfg \trip{\mc{D}}{m}{\rho_\mf{tl}}$ via 
% \emph{(i)}, \emph{(ii)}, and (2),
% and we apply the inductive hypothesis to (3) and (7)
% to get two more subcases: either 
% (8a) $\trip{\mc{D}}{m}{\rho_\mf{tl}}\cstep \trip{\mc{D}'}{m'}{\rho'_\mf{tl}}$
% or (8b) $\rho_\mf{tl} = u_\mf{tl}$.
% \item In subcase (8a), we apply \ts{ER-Ctx} with $\mc{E}= u_\mf{hd} ::_\rho [\cdot]$,
% satisfying conclusion \emph{(ii)}:
% \begin{mathpar}
% \inferrule*[right=ER-Ctx] {
%   \inferrule*[right=\tn{(8a)}] { } {
%     \trip{\mc{D}}{m}{\rho_\mf{tl}} \cstep \trip{\mc{D}'}{m'}{\rho_\mf{tl}'}
%   }
% } {
%   \trip{\mc{D}}{m}{u_\mf{hd} ::_\rho \rho_\mf{tl}}
%   \cstep 
%   \trip{\mc{D}'}{m'}{u_\mf{hd} ::_\rho \rho_\mf{tl}'}
% }
% \end{mathpar}
% \item Otherwise, by (6b) and (8b), we have that $\rho = u_\mf{hd} ::_\rho u_\mf{tl}$,
% i.e. $\rho$ is a value, satisfying conclusion \emph{(i)}.
% \end{itemize}

% \tb{Case} \ts{TR-Node}: \begin{mathpar}
% \inferrule*[right=TR-Node] {
%   \cdot \semi \Pi \vdval v: [\mf{Attr}]
%   \and
%   \inferrule*[right=(1)] { } {\cdot \semi \Pi \vdrt \rho: [\mf{Tree}]_\rho}
% } {
%   \cdot \semi \Pi \vdrt \mf{node}(\mf{id}, v, \rho): \mf{Tree}
% }
% \and
% \inferrule*[right=WF-Node] {
%   \inferrule*[right=(2)] { } {\mc{D} \semi m \wf \rho}
% } {
%   \mc{D} \semi m \wf \mf{node}(\mf{id}, v, \rho)
% }
% \end{mathpar}
% \begin{itemize}
% \item By \emph{(i)}, \emph{(ii)}, and (1), we have 
% (3) $\cdot \semi \Pi \vdcfg \trip{\mc{D}}{m}{\rho}$.
% We apply the inductive hypothesis to (2) and (3) to get that either
% (4a) $\rho = u$ or (4b) 
% $\trip{\mc{D}}{m}{\rho}
% \cstep 
% \trip{\mc{D}'}{m'}{\rho'}$.
% \item In subcase (4a), we have that $\mf{node}(\mf{id}, v, u)$ is a value,
% satisfying conclusion \emph{(i)}.
% \item In subcase (4b), we apply \ts{ER-Ctx} with $\mc{E}=\mf{node}(\mf{id}, v, [\cdot])$
% to satisfy conclusion \emph{(ii)} as follows:
% \begin{mathpar}
% \inferrule*[right=ER-Ctx] {
%   \inferrule*[right=\tn{(4b)}] { } {
%     \trip{\mc{D}}{m}{\rho} \cstep \trip{\mc{D}'}{m'}{\rho'}
%   }
% } {
%   \trip{\mc{D}}{m}{\mf{node}(\mf{id}, v, \rho)}
%   \cstep 
%   \trip{\mc{D}'}{m'}{\mf{node}(\mf{id}, v, \rho')}
% }
% \end{mathpar}
% \end{itemize}

% \tb{Case} \ts{TR-Init}: \begin{mathpar}
% \inferrule*[right=TR-Init] {
%   \inferrule*[right=(1)] { } {\cdot \semi \Pi \vdval v: \mf{View}}
% } {
%   \cdot \semi \Pi \vdrt \mf{init}\ v: \mf{Tree}
% }
% \end{mathpar}
% \begin{itemize}
% \item By \thref{thm:canon} (canonical forms) on (1),
% we have three subcases: either
% (a) $v=\mf{text}\ v'$,
% (b) $v=\mf{tag}(v_1, v_2)$, or
% (c) $v=\pair{\mc{C}}{v_0}$.
% \item In subcase (a), \ts{ER-IText} simply says
% $\trip{\mc{D}}{m}{\mf{init\ text}\ v'} \cstep
% \trip{\mc{D}}{m}{\mf{lit}\ v'}$.
% \item In subcase (b), we apply \ts{ER-ITag} (with some fresh $\mf{id}$)
% to get that $\trip{\mc{D}}{m}{\mf{init\ tag}(v_1, v_2)} \cstep
% \trip{\mc{D}}{m}{\mf{node}(\mf{id}, v_1, \mf{inits}\ v_2)}$.
% \item (HERE DOWN NEEDS REFACTORING)
% \item In subcase (c), we appeal to our well-formedness premise:
% \begin{mathpar}
% \inferrule*[right=WF-Init] {
%   \inferrule*[right=WF-Spec] {
%     \inferrule*[right=(2)] { } {\mc{C} \in \mc{D}}
%   } {
%     \mc{D} \wf \pair{\mc{C}}{v_0}
%   }
% } {
%   \mc{D} \semi m \wf \mf{init}\ \pair{\mc{C}}{v_0}
% }
% \end{mathpar}
% \item By (2), we can now apply \ts{ER-ISpec} (with some fresh $p$:)
% \begin{mathpar}
% \inferrule*[right=ER-ISpec] {
%   p\ \mf{fresh}
%   \and 
%   \inferrule*[right=(2)] { } {\mc{D}[\mc{C}]=x.\hat{e}}
%   \\
%   \pi_0=\{
%     \mf{spec}=\pair{\mc{C}}{v_0}, 
%     \mf{st}: \cdot, 
%     \mf{tree}: \cdot, 
%     \mf{rerender}: \mf{ff}
%   \}
% } {
%   \trip{\mc{D}}{m}{\mf{init}\ \pair{\mc{C}}{v}}
%   \cstep
%   \trip{\mc{D}}{m[p \mapsto \pi_0]}{\mf{eval}^\mf{Init}_p\ \hat{e}[v/x]}
% }
% \end{mathpar}
% \item Thus, in any subcase, we satisfy conclusion \emph{(ii)}.
% \end{itemize}

% \tb{Case} \ts{TR-Inits}: IH/canonical lists, inil/icons

% \tb{Case} \ts{TR-Scan}: \begin{mathpar}
% \inferrule*[right=TR-Scan] {
%   \cdot \semi \Pi \vdval v: \mf{Attr}
% } {
%   \cdot \semi \Pi \vdrt \mf{scan}\ v: \one_\rho
% }
% \end{mathpar}
% \begin{itemize}
% \item By \thref{thm:canon} (canonical forms) on $v$,
% either (a) $v=\mf{attr}(v_1, v_2)$ or (b) $v=\mf{onClick}\ v'$.
% \item In subcase (a), \ts{ER-ScAttr} says
% $\trip{\mc{D}}{m}{\mf{scan}\ \mf{attr}(v_1, v_2)} \cstep \trip{\mc{D}}{m}{\mf{done}}$.
% \item In subcase (b), we apply \ts{ER-ScClick} to get
% $\trip{\mc{D}}{m}{\mf{scan}\ \mf{onClick}(v')} \cstep \trip{\mc{D}}{m}{\mf{handle}\ (v\ ())}$.
% \item Thus, in either subcase we satisfy conclusion \emph{{ii}}.
% \end{itemize}

% \tb{Case} \ts{TR-Scans}: \begin{mathpar}
% \inferrule*[right=TR-Scans] {
%   \Gamma \semi \Pi \vdval v: [\mf{Attr}]
% } {
%   \Gamma \semi \Pi \vdrt \mf{scans}\ v: \one_\rho
% }
% \end{mathpar}
% \begin{itemize}
% \item By \thref{thm:canon} (canonical forms) on $v$,
% we have two subcases: (a) $v=[]$, or else (b) $v=v_\mf{hd} :: v_\mf{tl}$.
% \item In subcase (a), \ts{ER-ScsNil} gives
% $\trip{\mc{D}}{m}{\mf{scans}\ []} \cstep \trip{\mc{D}}{m}{\mf{done}}$.
% \item In subcase (b), we apply \ts{ER-ScsCons} to get
% $\trip{\mc{D}}{m}{\mf{scans}\ v_\mf{hd} :: v_\mf{tl}} 
% \cstep \trip{\mc{D}}{m}{\mf{scan}\ v_\mf{hd} \semi \mf{scans}\ v_\mf{tl}}$.
% \item Thus, in either subcase we satisfy conclusion \emph{{ii}}.
% \end{itemize}

% \tb{Case} \ts{TR-Chk}: IH on $\rho$, if not val then e-ctx,
% if val then canonical on rho: either lit, node, or p; rules for each
% (wf-path tells us m[p], d[c] are defined)

% \tb{Case} \ts{TR-Chks}: IH/canonical on lists, chksnil/chkscons

% \tb{Case} \ts{TR-Srch}: IH/e-ctx or canonically tree => step,
% m[p].tree defined by wf p

% \tb{Case} \ts{TR-Srchs}: standard for lists

% \tb{Case} \ts{TR-Eval $(\phi=\mf{Init})$}: progtle-i,
% if e steps then er-eval, else ret v then er-new

% \tb{Case} \ts{TR-Eval $(\phi=\mf{Succ})$}: analogous,
% but uses wf propery that $\mc{L}(\hat{e}) \subseteq \dom(m[p].\mf{st})$
% to apply progtle-s

% \tb{Case} \ts{TR-New}: IH/e-ctx, canon tree -> step newi/news

% \tb{Case} \ts{TR-Hdl}: progexp, er-hdl or canon unit -> er-hdldone

% \tb{Case} \ts{TR-Prog}: progprog, er-prog or canon end -> er-pdone

% \tb{Case} \ts{TR-Seq}: IH/e-ctx, canon done -> er-seq
% \end{Proof}
% \end{lemma}

\begin{lemma}[Preservation of configurations.]
\label{thm:prezcfg}
If \begin{enumerate}[label=(\roman*)]
\item $\trip{\mc{D}}{m}{\rho} \cstep \trip{\mc{D}'}{m'}{\rho'}$,
\item $\Gamma \semi \Pi \vdcfg \trip{\mc{D}}{m}{\rho}: R$, and
\item $\mc{D} \semi m \wf \rho$,
\end{enumerate}
then there exists some $\Pi'$ such that
\begin{enumerate}[label=(\roman*)]
\item $\Pi \po \Pi'$,
\item $\Gamma \semi \Pi' \vdcfg \trip{\mc{D}'}{m'}{\rho'}: R$, and
\item $\mc{D}' \semi m' \wf \rho'$.
\end{enumerate}

\begin{Proof}
To show conclusion \emph{(ii)}, we need the following three sub-conclusions:
\begin{enumerate}[label=(\emph{ii}-\alph*)]
\item $\Gamma \semi \Delta \vdd \mc{D}'$,
\item $\Gamma \semi \Pi' \vdm m'$, and
\item $\Gamma \semi \Pi' \vdrt \rho': R$.
\end{enumerate}

(We refer to premise \emph{(ii)} in similar fashion, 
e.g. \emph{(ii)}-a says $\Gamma \semi \Delta \vdd \mc{D}$.)

Unless otherwise mentioned, we let $\Pi'=\Pi$, 
and conclusion \emph{(i)} holds trivially.

We proceed by induction on the small-step derivation \emph{(iii)}.

\tb{Case} \ts{ER-Ctx}: \hl{TODO}; need eval ctx typing, but should be standard

\tb{Case} \ts{ER-Seq}: We have
$\trip{\mc{D}}{m}{\mf{done} \semi \rho}
\cstep
\trip{\mc{D}}{m}{\rho}$ and, by \emph{(ii)}, \begin{mathpar}
\inferrule*[right=TR-Seq] {
  \Gamma \semi \Pi \vdrt \mf{done}: \one_\rho
  \and
  \inferrule*[right=(1)] { } {\Gamma \semi \Pi \vdrt \rho: R}
} {
  \Gamma \semi \Pi \vdrt \mf{done} \semi \rho: R
}
\end{mathpar}
\begin{itemize}
\item Conclusion \emph{(ii)} follows from (1).
\item Conclusion \emph{(iii)} is given by premise \emph{(iii)}
as follows: \begin{mathpar}
\inferrule*[right=WF-Seq] {
  \inferrule*[right=WF-Done] { } {\mc{D} \semi m \wf \mf{done}}
  \and
  \inferrule*[right=\textnormal{\emph{(iii)}}] { } {\mc{D} \semi m \wf \rho}
} {
  \mc{D} \semi m \wf \mf{done} \semi \rho
}
\end{mathpar}
\end{itemize}

\tb{Case} \ts{ER-Prog}: We have \begin{mathpar}
\inferrule*[right=ER-Prog] {
  \inferrule*[right=(1)] { } {\mc{D} \semi P \pstep P' \semi \mc{D}'}
} {
  \trip{\mc{D}}{m}{\mf{prog}\ P}
  \cstep
  \trip{\mc{D}'}{m}{\mf{prog}\ P'}
}
\and
\inferrule*[right=TR-Prog] {
  \inferrule*[right=(2)] { } {\Gamma \semi \Pi \vdprog P}
  \and
  \inferrule*[right=(3)] { } {\Delta[\mc{C}_0] = \one}
} {
  \Gamma \semi \Pi \vdrt \mf{prog}\ P: \one_\rho
}
\and
\inferrule*[right=WF-Prog] {
  \inferrule*[right=(4)] { } {
    \mc{L}(P) \cup \dom(\mc{D}) \cup \{\mc{C}_0\}= \dom(\Delta)
  }
} {
  \mc{D} \semi m \wf \mf{prog}\ P
}
\end{mathpar}
\begin{itemize}
\item \thref{thm:prezprog} (preservation of programs) on (1) and (2)
gives us that
(5) $\Gamma \semi \Pi \vdprog P'$,
(6) $\Gamma \semi \Delta \vdd \mc{D}'$, and
(7) $\mc{L}(P') \cup \dom(\mc{D}') = \mc{L}(P) \cup \dom(\mc{D})$.
\item To show conclusion \emph{(ii)}, it suffices to show that
$\Gamma \semi \Pi \vdrt \mf{prog}\ P': \one_\rho$ and $\Gamma \semi \Delta \vdd \mc{D}'$.
The latter is precisely given by (6), and
the former is given by the following derivation:
\begin{mathpar}
\inferrule*[right=TR-Prog] {
  \inferrule*[right=(5)] { } {\Gamma \semi \Pi \vdprog P'}
  \and
  \inferrule*[right=(3)] { } {\Delta[\mc{C}_0] = \one}
} {
  \Gamma \semi \Pi \vdrt \mf{prog}\ P': \one_\rho
}
\end{mathpar}
\item By (4) and (7), we get (8) that 
$\mc{L}(P') \cup \dom(\mc{D'}) \cup \{\mc{C}_0\}= \dom(\Delta)$,
and then we can derive conclusion \emph{(iii)} as follows:
\begin{mathpar}
\inferrule*[right=WF-Prog] {
  \inferrule*[right=(8)] { } {
    \mc{L}(P') \cup \dom(\mc{D}') \cup \{\mc{C}_0\}= \dom(\Delta)
  }
} {
  \mc{D'} \semi m \wf \mf{prog}\ P'
}
\end{mathpar}
\end{itemize}

\tb{Case} \ts{ER-PDone}: We have
$\trip{\mc{D}}{m}{\mf{prog\ \eret{v}}} 
\cstep
\trip{\mc{D}'}{m'}{\mf{new}^\mf{Succ}_{p_0}\ \mf{init}\ v}$ and
\begin{mathpar}
\inferrule*[right=TR-Prog] {
  \inferrule*[right=TP-Lift] {
    \inferrule*[right=TE-Ret] {
      \inferrule*[right=(1)] { } {
        \Gamma \semi \cdot \vdval v: \mf{View}
      }
    } {
      \Gamma \semi \cdot \vdval \eret{v}: \mf{View} \mid \epsilon
    }
  } {
    \Gamma \vdprog \eret{v}
  }
  \and
  \inferrule*[right=(2)] { } {\Delta[\mc{C}_0] = \one}
} {
  \Gamma \semi \Pi \vdrt \mf{prog}\ \eret{v}: \one_\rho
}
\and
\inferrule*[right=WF-Prog] {
  \inferrule*[right=(3)] { } {
    \mc{L}(\eret{v}) \cup \dom(\mc{D}) \cup \{\mc{C}_0\}= \dom(\Delta)
  }
} {
  \mc{D} \semi m \wf \mf{prog}\ \eret{v}
}
\end{mathpar}
with abbreviations
\begin{align*}
m' &= m[p_0 \mapsto \{\mf{spec}: \pair{\mc{C}_0}{()}, \mf{st}: \cdot,
  \mf{tree}: \cdot, \mf{rerender}: \mf{ff}\}]\\
\mc{D}' &= \mc{D}[\mc{C}_0 \mapsto x_0.\eret{v}]
\end{align*}

\begin{itemize}
\item Choose $\Pi' = \Pi[p_0 \mapsto \mc{C}_0]$.
By global freshness of $p_0$, we have $p_0 \notin \dom(\Pi)$, 
so $\Pi \po \Pi'$,
i.e. conclusion \emph{(i)} holds.
\item To show conclusion (\emph{ii}-a) (that $\Gamma \semi \Delta \vdd \mc{D}'$),
by premise (\emph{ii}-a) and global freshness of $\mc{C}_0$, it suffices to show that 
$\Gamma[x_0 \mapsto \Delta[\mc{C}_0]] \semi \cdot \semi \mc{C}_0 \vdtop \eret{v}$.
(Say $\Gamma' = \Gamma[x_0 \mapsto \Delta[\mc{C}_0]]$.)
\item By global freshness of $x_0$, we have $x_0 \notin \dom(\Gamma)$,
and so $\Gamma \po \Gamma'$. By \thref{thm:wktpv} (weakening) on (1),
we get (4) that $\Gamma' \semi \cdot \vdval v: \mf{View}$,
and we derive the conclusion as follows:
\begin{mathpar}
\inferrule*[right=TH-Lift] {
  \inferrule*[right=TE-Ret] {
    \inferrule*[right=(4)] { } {
      \Gamma' \semi \cdot \vdval v: \mf{View}
    }
  } {
    \Gamma' \semi \cdot \vdexp \eret{v} : \mf{View} \mid \epsilon
  }
} {
  \Gamma' \semi \cdot \semi \mc{C}_0 \vdtop \eret{v}
}
\end{mathpar}
\item To show conclusion (\emph{ii}-b) (that $\Gamma \semi \Pi' \vdm m'$),
we first apply \thref{thm:wkpim} (weakening) to premise (\emph{ii}-b) to get that
$\Gamma \semi \Pi' \vdm m$. Now, by definition 
(and since both the $\mf{st}$ and $\mf{tree}$ fields of $m'[p_0]$ are empty),
it suffices to show that $\Pi'[p_0]=\mc{C}_0$ and 
$\Gamma \semi \Pi' \vdval (): \Delta[\mc{C}_0]$.
\item The former is true by definition of $\Pi'$, and the latter holds by (2)
and rule \ts{TVal-Unit}: $\Gamma \semi \Pi' \vdval (): \one$. Thus,
conclusion (\emph{ii}-b) holds.
\item By \thref{thm:wkpiv} (weakening) on (1),
we get (5) that $\Gamma \semi \Pi' \vdval v: \mf{View}$. 
We derive conclusion (\emph{ii}-c) as follows:
\begin{mathpar}
\inferrule*[right=TR-New] {
  \inferrule*[right=TR-Init] {
    \inferrule*[right=(5)] { } {\Gamma \semi \Pi' \vdval v: \mf{View}}
  } {
    \Gamma \semi \Pi' \vdrt \mf{init}\ v: \mf{Tree}
  }
} {
  \Gamma \semi \Pi' \vdrt \mf{new}^\mf{Succ}_{p_0}\ \mf{init}\ v: \one_\rho
}
\end{mathpar}
\item Finally, to show conclusion \emph{(iii)}, we first observe that
$\mc{L}(\eret{v})=\emptyset$, so from (3) we get (6) that
$\dom(\mc{D}) \cup \{\mc{C}_0\} = \dom(\Delta)$.
By definition of $\mc{D}'$, $\dom(\mc{D}')=\dom(\mc{D}) \cup \{\mc{C}_0\}$,
so therefore $\dom(\mc{D}') = \dom(\Delta)$, satisfying the implicit premise.
\item We construct the following derivation, where (7) is given by
definition of $\mc{L}$ and construction of $m'$:
\begin{mathpar}
\inferrule*[right=WF-New] {
  \inferrule*[right=(7)] { } {\mc{L}(\eret{v}) = \dom(m'[p].\mf{st})}
  \quad
  \inferrule*[right=WF-Init] { } {\mc{D}' \semi m' \wf \mf{init}\ v}
} {
  \mc{D}' \semi m' \wf \mf{new}^\mf{Succ}_{p_0}\ \mf{init}\ v
}
\end{mathpar}
\item We have thus shown all desired conclusions, and the case holds.
\end{itemize}

\tb{Case} \ts{ER-ISpec}: We have
$\trip{\mc{D}}{m}{\mf{init}\ \pair{\mc{C}}{v}}
\cstep
\trip{\mc{D}}{m'}{\mf{eval}^\mf{Init}_p\ \hat{e}[v/x]}$ and
\begin{mathpar}
\inferrule*[right=TR-Init] {
  \inferrule*[right=TV-Spec] {
    \inferrule*[right=TV-Comp] {
      \inferrule*[right=(1)] { } {\Delta[\mc{C}] = \tau}
    } {
      \Gamma \semi \Pi \vdval \mc{C}: \mf{Component}\ \tau
    }
    \ 
    \inferrule*[right=(2)] { } {\Gamma \semi \Pi \vdval v: \tau}
  } {
    \Gamma \semi \Pi \vdval \pair{\mc{C}}{v}: \mf{View}
  }
} {
  \Gamma \semi \Pi \vdrt \mf{init}\ \pair{\mc{C}}{v}: \mf{Tree}
}
\end{mathpar}
with $p$ fresh and abbreviations
\begin{align*}
m' &= m[p \mapsto \{
    \mf{spec}=\pair{\mc{C}}{v}, 
    \mf{st}: \cdot, 
    \mf{tree}: \cdot, 
    \mf{rerender}: \mf{ff}
  \}]\\
\mc{D}[\mc{C}]&=x.\hat{e}
\end{align*}
\begin{itemize}
\item Choose $\Pi' = \Pi[p \mapsto \mc{C}]$.
Since $p$ is fresh, $\Pi \po \Pi'$, i.e. conclusion \emph{(i)} holds.
Furthermore, conclusion (\emph{ii}-a) is trivial.
\item We weaken premise (\emph{ii}-b) by \thref{thm:wkpim}
to get that $\Gamma \semi \Pi' \vdm m$. Thus, to show $\Gamma \semi \Pi' \vdm m'$,
by freshness of $p$, to show conclusion (\emph{ii}-b),
it suffices to show that $\Pi'[p]=\mc{C}$ (trivial)
and that $\Gamma \semi \Pi' \vdval v: \Delta[\mc{C}]$.
\item By (1) and (2), $\Gamma \semi \Pi \vdval v: \Delta[\mc{C}]$.
To this, we apply \thref{thm:wkpiv} (weakening)
to show conclusion (\emph{ii}-b) holds.
\item To show conclusion (\emph{ii}-c), we appeal to premise (\emph{ii}-a),
which gives us that 
$\Gamma[x \mapsto \Delta[\mc{C}]] \semi \cdot \semi \mc{C} \vdtop \hat{e}$.
We first weaken this (by \thref{thm:wkpih}) to get 
$\Gamma[x \mapsto \Delta[\mc{C}]] \semi \Pi \semi \mc{C} \vdtop \hat{e}$,
where we can then apply \thref{thm:subtle} (substitution)
with value from (2) and using the equality from (1) to get
$\Gamma \semi \Pi \semi \mc{C} \vdtop \hat{e}[v/x]$.
\item We apply \thref{thm:wkpih} once more to say (4) that
$\Gamma \semi \Pi' \semi \mc{C} \vdtop \hat{e}[v/x]$.
We note that $\mc{C}=\Pi'[p]$, and we conclude with the derivation:
\begin{mathpar}
\inferrule*[right=TR-Eval] {
  \inferrule*[right=(4)] { } {
    \Gamma \semi \Pi' \semi \Pi'[p] \vdtop \hat{e}[v/x]
  }
} {
  \Gamma \semi \Pi' \vdrt \mf{eval}^\mf{Init}_p\ \hat{e}[v/x]: \mf{Tree}
}
\end{mathpar}
\item Only conclusion \emph{(iii)} remains to be shown. Via the following
derivation, it suffices to show (5):
\begin{mathpar}
\inferrule*[right=WF-EvalI] {
  \inferrule*[right=(5)] { } {
    \mc{L}(\hat{e}[v/x]) \cup \dom(m'[p].\mf{st}) = \mc{L}(\hat{e})
  }
} {
  \mc{D} \semi m' \wf \mf{eval}^\mf{Init}_p\ \hat{e}[v/x]
}
\end{mathpar}
\item By construction, $\dom(m'[p].\mf{st})=\emptyset$, so
it suffices to show $\mc{L}(\hat{e}[v/x]) = \mc{L}(\hat{e})$.
This is given directly by \thref{thm:sublab} (substitution for label-sets),
and we are done.
\end{itemize}

\tb{Case} \ts{ER-IText}: We have 
$\trip{\mc{D}}{m}{\mf{init\ text}\ v}
\cstep
\trip{\mc{D}}{m}{\mf{lit}\ v}$ and: \begin{mathpar}
\inferrule*[right=TR-Init] {
  \inferrule*[right=TV-Text] {
    \inferrule*[right=(1)] { } {\Gamma \semi \Pi \vdval v: \mf{String}}
  } {
    \Gamma \semi \Pi \vdval \mf{text}\ v: \mf{View}
  }
} {
  \Gamma \semi \Pi \vdrt \mf{init\ text}\ v: \mf{Tree}
}
\end{mathpar}
\begin{itemize}
\item Conclusion (\emph{ii}-c) is given by the following derivation:
\begin{mathpar}
\inferrule*[right=TR-Lit] {
  \inferrule*[right=(1)] { } {\Gamma \semi \Pi \vdval v: \mf{String}}
} {
  \Gamma \semi \Pi \vdrt \mf{lit}\ v: \mf{Tree}
}
\end{mathpar}
\item Conclusions (\emph{ii}-a), (\emph{ii}-b), and \emph{(iii)}
are all trivial, and we are done.
\end{itemize}

\tb{Case} \ts{ER-ITag}: We have 
$\trip{\mc{D}}{m}{\mf{init\ tag}(v_1, v_2)}
\cstep
\trip{\mc{D}}{m}{\mf{node}(\mf{id}, v_1, \mf{inits}\ v_2)}$ and:
\begin{mathpar}
\inferrule*[right=TR-Init] {
  \inferrule*[right=TV-Tag] {
    \inferrule*[right=(1)] { } {\Gamma \semi \Pi \vdval v_1: [\mf{Attr}]}
    \and
    \inferrule*[right=(2)] { } {\Gamma \semi \Pi \vdval v_2: [\mf{View}]}
  } {
    \Gamma \semi \Pi \vdval \mf{tag}(v_1, v_2): \mf{View}
  }
} {
  \Gamma \semi \Pi \vdrt \mf{init\ tag}(v_1, v_2): \mf{Tree}
}
\end{mathpar}
\begin{itemize}
\item Conclusion (\emph{ii}-c) is given by the following derivation:
\begin{mathpar}
\inferrule*[right=TR-Node] {
  \inferrule*[right=(1)] { } {\Gamma \semi \Pi \vdval v_1: [\mf{Attr}]}
  \
  \inferrule*[right=TR-Inits] {
    \inferrule*[right=(2)] { } {\Gamma \semi \Pi \vdval v_2: [\mf{View}]}
  } {
    \Gamma \semi \Pi \vdrt \mf{inits}\ v_2: [\mf{Tree}]_\rho
  }
} {
  \Gamma \semi \Pi \vdrt \mf{node}(\mf{id}, v_1, \mf{inits}\ v_2): \mf{Tree}
}
\end{mathpar}
\item Conclusions (\emph{ii}-a) and (\emph{ii}-b) are trivial. \emph{(iii)}
is shown as follows:
\begin{mathpar}
\inferrule*[right=WF-Node] {
  \inferrule*[right=WF-Inits] { } {\mc{D} \semi m \wf \mf{inits}\ v_2}
} {
  \mc{D} \semi m \wf \mf{node}(\mf{id}, v, \mf{inits}\ v_2)
}
\end{mathpar}
\end{itemize}

\tb{Case} \ts{ER-INil}: We have
$\trip{\mc{D}}{m}{\mf{inits}\ []}
\cstep
\trip{\mc{D}}{m}{[]_\rho}$ and:
\begin{mathpar}
\inferrule*[right=TR-Inits] {
  \Gamma \semi \Pi \vdval []: [\mf{View}]
} {
  \Gamma \semi \Pi \vdrt \mf{inits}\ []: [\mf{Tree}]_\rho
}
\end{mathpar}
\begin{itemize}
\item The typing conclusion (\emph{ii}-c) is immediate: 
\begin{mathpar}
\inferrule*[right=TR-Nil] { } {
  \Gamma \semi \Pi \vdrt []_\rho: [\mf{Tree}]_\rho
}
\end{mathpar}
\item Conclusions (\emph{ii}-a), (\emph{ii}-b), and \emph{(iii)} are all trivial.
\end{itemize}

\tb{Case} \ts{ER-ICons}: We have
$\trip{\mc{D}}{m}{\mf{inits}\ (v :: vs)}
\cstep
\trip{\mc{D}}{m}{(\mf{init}\ v) ::_\rho (\mf{inits}\ vs)}$ and:
\begin{mathpar}
\inferrule*[right=TR-Inits] {
  \inferrule*[right=TV-Cons] {
    \inferrule*[right=(1)] { } {\Gamma \semi \Pi \vdval v: \mf{View}}
    \and
    \inferrule*[right=(2)] { } {\Gamma \semi \Pi \vdval vs: [\mf{View}]}
  } {
    \Gamma \semi \Pi \vdval v :: vs: [\mf{View}]
  }
} {
  \Gamma \semi \Pi \vdrt \mf{inits}\ (v :: vs): [\mf{Tree}]_\rho
}
\end{mathpar}
\begin{itemize}
\item We derive conclusion (\emph{ii}-c) from (1) and (2) as follows:
\end{itemize}
\begin{mathpar}
\inferrule*[right=TR-Cons] {
  \inferrule*[right=TR-Init] {
    \inferrule*[right=(1)] { } {\Gamma \semi \Pi \vdval v: \mf{View}}
  } {
    \Gamma \semi \Pi \vdrt \mf{init}\: v : \mf{Tree}
  }
  \,
  \inferrule*[right=TR-Inits] {
    \inferrule*[right=(2)] { } {\Gamma \semi \Pi \vdval vs: [\mf{View}]}
  } {
    \Gamma \semi \Pi \vdrt \mf{inits}\: vs : [\mf{Tree}]_\rho
  }
} {
  \Gamma \semi \Pi \vdrt (\mf{init}\ v) ::_\rho (\mf{inits}\ vs) : [\mf{Tree}]_\rho
}
\end{mathpar}
\begin{itemize}
\item Conclusions (\emph{ii}-a), (\emph{ii}-b), and \emph{(iii)} are all trivial.
\end{itemize}

\tb{Case} \ts{ER-Eval ($\phi=\mf{Init}$)}: We have
\begin{mathpar}
\inferrule*[right=ER-Eval] {
  \inferrule*[right=(1)] { } {
    m \semi \hat{e} \topstep{\mf{Init}}{p} \hat{e}' \semi m'
  }
} {
  \trip{\mc{D}}{m}{\mf{eval}^\mf{Init}_p\ \hat{e}}
  \cstep
  \trip{\mc{D}}{m'}{\mf{eval}^\mf{Init}_p\ \hat{e}'}
}
\and
\inferrule*[right=TR-Eval] {
  \inferrule*[right=(2)] { } {\Gamma \semi \Pi \semi \Pi[p] \vdtop \hat{e}}
} {
  \Gamma \semi \Pi \vdrt \mf{eval}^\mf{Init}_p\ \hat{e}: \mf{Tree}
}
\and
\inferrule*[right=WF-EvalI] {
  \inferrule*[right=(3)] { } {\mc{L}(\hat{e}) \cup \dom(m[p].\mf{st}) = \mc{L}(\hat{e}_0)}
} {
  \mc{D} \semi m \wf \mf{eval}^\mf{Init}_p\ \hat{e}
}
\end{mathpar}
\begin{itemize}
\item It follows from the well-definedness 
of $\hat{e}_0$ in (3) that $p \in \dom(m)$. 
Therefore, we can apply \thref{thm:preztlei} 
(top-level preservation at $\phi=\mf{Init}$) 
using (2), premise (\emph{ii}-b),and (1) to get the following:
\begin{enumerate}[label=(\arabic*),start=(4)]
\item $\Gamma \semi \Pi \semi \Pi[p] \vdtop \hat{e}'$,
\item $\Gamma \semi \Pi \vdm m'$, and
\item $\mc{L}(\hat{e}) \cup \dom(m[p].\mf{st}) 
= \mc{L}(\hat{e}') \cup \dom(m'[p].\mf{st})$.
\end{enumerate}
\item Conclusion (\emph{ii}-a) is trivial, and (\emph{ii}-b) is given by (5).
(\emph{ii}-c) is shown as follows:
\begin{mathpar}
\inferrule*[right=TR-Eval] {
  \inferrule*[right=(4)] { } {\Gamma \semi \Pi \semi \Pi[p] \vdtop \hat{e}'}
} {
  \Gamma \semi \Pi \vdrt \mf{eval}^\mf{Init}_p\ \hat{e}': \mf{Tree}
}
\end{mathpar}
\item By \thref{thm:prezspec} (spec preservation) and (3),
we have that $m'[p].\mf{spec} = m[p].\mf{spec}$,
and therefore $\hat{e}_0$ remains the same between $m$ and $m'$.
By (4) and (6), we get (8) that
$\mc{L}(\hat{e}') \cup \dom(m'[p].\mf{st}) = \mc{L}(\hat{e}_0)$.
We can now derive conclusion \emph{(iii)} as follows:
\end{itemize}
\begin{mathpar}
\inferrule*[right=WF-EvalI] {
  \inferrule*[right=(8)] { } {
    \mc{L}(\hat{e}') \cup \dom(m'[p].\mf{st}) = \mc{L}(\hat{e}_0)
  }
} {
  \mc{D} \semi m' \wf \mf{eval}^\mf{Init}_p\ \hat{e}'
}
\end{mathpar}

\tb{Case} \ts{ER-Eval ($\phi=\mf{Succ}$)}: We have
\begin{mathpar}
\inferrule*[right=ER-Eval] {
  \inferrule*[right=(1)] { } {
    m \semi \hat{e} \topstep{\mf{Succ}}{p} \hat{e}' \semi m'
  }
} {
  \trip{\mc{D}}{m}{\mf{eval}^\mf{Succ}_p\ \hat{e}}
  \cstep
  \trip{\mc{D}}{m'}{\mf{eval}^\mf{Succ}_p\ \hat{e}'}
}
\and
\inferrule*[right=TR-Eval] {
  \inferrule*[right=(2)] { } {\Gamma \semi \Pi \semi \Pi[p] \vdtop \hat{e}}
} {
  \Gamma \semi \Pi \vdrt \mf{eval}^\mf{Succ}_p\ \hat{e}: \one_\rho
}
\and
\inferrule*[right=WF-EvalS] {
  \inferrule*[right=(3)] { } {\mc{L}(\hat{e}_0) = \dom(m[p].\mf{st})}
  \and
  \inferrule*[right=(4)] { } {\mc{L}(\hat{e}) \subseteq \dom(m[p].\mf{st})}
} {
  \mc{D} \semi m \wf \mf{eval}^\mf{Succ}_p\ \hat{e}
}
\end{mathpar}
\begin{itemize}
\item By (4), we again have $p \in \dom(m)$, so we can apply \thref{thm:preztles}
(top-level preservation at $\phi=\mf{Succ}$)
using (2), premise (\emph{ii}-b), and (1) to get the following:
\begin{enumerate}[label=(\arabic*),start=(5)]
\item $\Gamma \semi \Pi \semi \Pi[p] \vdtop \hat{e}'$,
\item $m = m'$, and
\item $\mc{L}(\hat{e}') \subseteq \mc{L}(\hat{e})$.
\end{enumerate}
\item Conclusion (\emph{ii}-a) is trivial, 
and (\emph{ii}-b) is immediate from (6) and premise (\emph{ii}-b).
Conclusion (\emph{ii}-c) is shown as follows:
\begin{mathpar}
\inferrule*[right=TR-Eval] {
  \inferrule*[right=(5)] { } {\Gamma \semi \Pi \semi \Pi[p] \vdtop \hat{e}'}
} {
  \Gamma \semi \Pi \vdrt \mf{eval}^\mf{Succ}_p\ \hat{e}': \one_\rho
}
\end{mathpar}
\item Finally, (4) and (7) give us (8) 
that $\mc{L}(\hat{e}') \subseteq \dom(m[p].\mf{st})$,
and conclusion \emph{(iii)} follows:
\begin{mathpar}
\inferrule*[right=WF-EvalS] {
    \inferrule*[right=(3)] { } {\mc{L}(\hat{e}_0) = \dom(m[p].\mf{st})}
  \and
  \inferrule*[right=(8)] { } {\mc{L}(\hat{e}') \subseteq \dom(m[p].\mf{st})}
} {
  \mc{D} \semi m \wf \mf{eval}^\mf{Succ}_p\ \hat{e}'
}
\end{mathpar}
\end{itemize}

\tb{Case} \ts{ER-New ($\phi=\mf{Init}$)}: We have
$$\trip{\mc{D}}{m}{\mf{eval}^\mf{Init}_p\ \eret{v}}
\cstep
\trip{\mc{D}}{m}{\mf{new}^\mf{Init}_p\ \mf{init}\ v}$$
\begin{mathpar}
\inferrule*[right=TR-Eval] {
  \inferrule*[right=TH-Lift] {
    \inferrule*[right=TE-Ret] {
      \inferrule*[right=(1)] { } {
        \Gamma \semi \Pi \vdval v: \mf{View}
      }
    } {
      \Gamma \semi \Pi \vdexp \eret{v} : \mf{View} \mid \epsilon
    }
  } {
    \Gamma \semi \Pi \semi \Pi[p] \vdtop \eret{v}
  }
} {
  \Gamma \semi \Pi \vdrt \mf{eval}^\mf{Init}_p\ \eret{v}: \mf{Tree}
}
\and
\inferrule*[right=WF-EvalI] {
  \inferrule*[right=(2)] { } {
    \mc{L}(\eret{v}) \cup \dom(m[p].\mf{st}) = \mc{L}(\hat{e}_0)
  }
} {
  \mc{D} \semi m \wf \mf{eval}^\mf{Init}_p\ \eret{v}
}
\end{mathpar}
\begin{itemize}
\item Conclusions (\emph{ii}-a) and (\emph{ii}-b) are trivial.
The following derivation gives conclusion (\emph{ii}-c):
\begin{mathpar}
\inferrule*[right=TR-New] {
  \inferrule*[right=TR-Init] {
    \inferrule*[right=(1)] { } {
      \Gamma \semi \Pi \vdval v: \mf{View}
    }
  } {
    \Gamma \semi \Pi \vdrt \mf{init}\ v: \mf{Tree}
  }
} {
  \Gamma \semi \Pi \vdrt \mf{new}^\mf{Init}_p\ \mf{init}\ v: \mf{Tree}
}
\end{mathpar}
\item To show conclusion \emph{(iii)}, we note that $\mc{L}(\eret{v}) = \emptyset$
by definition, so (2) reduces to (3) $\dom(m[p].\mf{st}) = \mc{L}(\hat{e}_0)$.
Well-formedness then follows by derivation:
\begin{mathpar}
\inferrule*[right=WF-New] {
  \inferrule*[right=(3)] { } {\mc{L}(\hat{e}_0) = \dom(m[p].\mf{st})}
  \and
  \inferrule*[right=WF-Init] { } {\mc{D} \semi m \wf \mf{init}\ v}
} {
  \mc{D} \semi m \wf \mf{new}^\mf{Init}_p\ \mf{init}\ v
}
\end{mathpar}
\end{itemize}

\tb{Case} \ts{ER-New ($\phi=\mf{Succ}$)}: We have
$$\trip{\mc{D}}{m}{\mf{eval}^\mf{Succ}_p\ \eret{v}}
\cstep
\trip{\mc{D}}{m}{\mf{new}^\mf{Succ}_p\ \mf{init}\ v}$$
\begin{mathpar}
\inferrule*[right=TR-Eval] {
  \inferrule*[right=TH-Lift] {
    \inferrule*[right=TE-Ret] {
      \inferrule*[right=(1)] { } {
        \Gamma \semi \Pi \vdval v: \mf{View}
      }
    } {
      \Gamma \semi \Pi \vdexp \eret{v} : \mf{View} \mid \epsilon
    }
  } {
    \Gamma \semi \Pi \semi \Pi[p] \vdtop \eret{v}
  }
} {
  \Gamma \semi \Pi \vdrt \mf{eval}^\mf{Succ}_p\ \eret{v}: \one_\rho
}
\and
\inferrule*[right=WF-EvalS] {
  \inferrule*[right=(2)] { } {\mc{L}(\hat{e}_0) = \dom(m[p].\mf{st})}
  \and
  \mc{L}(\eret{v}) \subseteq \dom(m[p].\mf{st})
} {
  \mc{D} \semi m \wf \mf{eval}^\mf{Succ}_p\ \eret{v}
}
\end{mathpar}
\begin{itemize}
\item Conclusions (\emph{ii}-a) and (\emph{ii}-b) are trivial.
The following derivation gives conclusion (\emph{ii}-c):
\begin{mathpar}
\inferrule*[right=TR-New] {
  \inferrule*[right=TR-Init] {
    \inferrule*[right=(1)] { } {
      \Gamma \semi \Pi \vdval v: \mf{View}
    }
  } {
    \Gamma \semi \Pi \vdrt \mf{init}\ v: \mf{Tree}
  }
} {
  \Gamma \semi \Pi \vdrt \mf{new}^\mf{Succ}_p\ \mf{init}\ v: \one_\rho
}
\end{mathpar}
\item Conclusion \emph{(iii)} also follows quickly, and we are done:
\begin{mathpar}
\inferrule*[right=WF-New] {
  \inferrule*[right=(2)] { } {\mc{L}(\hat{e}_0) = \dom(m[p].\mf{st})}
  \and
  \inferrule*[right=WF-Init] { } {\mc{D} \semi m \wf \mf{init}\ v}
} {
  \mc{D} \semi m \wf \mf{new}^\mf{Succ}_p\ \mf{init}\ v
}
\end{mathpar}
\end{itemize}

\tb{Case} \ts{ER-NewI}: We have
$\trip{\mc{D}}{m}{\mf{new}^\mf{Init}_p\ u}
\cstep
\trip{\mc{D}}{m'}{p}$ and:
\begin{mathpar}
\inferrule*[right=TR-New] {
  \inferrule*[right=(1)] { } {\Gamma \semi \Pi \vdrt u: \mf{Tree}}
} {
  \Gamma \semi \Pi \vdrt \mf{new}^\mf{Init}_p\ u: \mf{Tree}
}
\and
\inferrule*[right=WF-New] {
  \inferrule*[right=(2)] { } {\mc{L}(\hat{e}_0) = \dom(m[p].\mf{st})}
  \and
  \inferrule*[right=(3)] { } {\mc{D} \semi m \wf u}
} {
  \mc{D} \semi m \wf \mf{new}^\mf{Init}_p\ u
}
\end{mathpar}
where $m' = m \mid m[p].\mf{tree} = u$.
\begin{itemize}
\item Conclusion (\emph{ii}-a) is trivial.
\item To show conclusion (\emph{ii}-b) given premise (\emph{ii}-b),
it suffices to show that $\Gamma \semi \Pi \vdrt u: \mf{Tree}$.
This is given directly by (1).
\item To show conclusion (\emph{ii}-c), we note that,
since $p \in \dom(m)$ (e.g. by (2))
and $\Gamma \semi \Pi \vdm m$ (by premise (\emph{ii}-b)),
it must be that (4) $p \in \Pi$. Also
from the memory typing judgment, $\Pi[p] =\mc{C}$
for some $\mc{C} \in \dom(\Delta)$, i.e. (5) $\Pi[p] \in \dom(\Delta)$.
Conclusion (\emph{ii}-c) now follows from the derivation:
\begin{mathpar}
\inferrule*[right=TR-Path] {
  \inferrule*[right=(4)] { } {p \in \dom(\Pi)}
  \and
  \inferrule*[right=(5)] { } {\Pi[p] \in \dom(\Delta)}
} {
  \Gamma \semi \Pi \vdrt p: \mf{Tree}
}
\end{mathpar}
\item Finally, we need to show conclusion \emph{(iii)}.
Clearly $m'[p].\mf{spec} = m[p].\mf{spec}$ and $m'[p].\mf{st} = m[p].\mf{st}$,
which lets us conclude from (2) that (6) $\mc{L}(\hat{e}_0) = \dom(m'[p].\mf{st})$.
We construct the following derivation:
\begin{mathpar}
\inferrule*[right=WF-Path] {
  \inferrule*[right=(6)] { } {\mc{L}(\hat{e}_0) = \dom(m'[p].\mf{st})}
  \and
  \inferrule*[right=(7)] { } {\mc{D} \semi m' \wf m'[p].\mf{tree}}
} {
  \mc{D} \semi m' \wf p
}
\end{mathpar}
\item Therefore, it suffices to show (7). \hl{This is... not given by (3)}
\end{itemize}

\tb{Case} \ts{ER-NewS}: We have
$\trip{\mc{D}}{m}{\mf{new}^\mf{Succ}_p\ u}
\cstep
\trip{\mc{D}}{m'}{\mf{done}}$ and:
\begin{mathpar}
\inferrule*[right=TR-New] {
  \inferrule*[right=(1)] { } {\Gamma \semi \Pi \vdrt u: \mf{Tree}}
} {
  \Gamma \semi \Pi \vdrt \mf{new}^\mf{Succ}_p\ u: \one_\rho
}
\and
\inferrule*[right=WF-New] {
  \mc{L}(\hat{e}_0) = \dom(m[p].\mf{st})
  \and
  \mc{D} \semi m \wf u
} {
  \mc{D} \semi m \wf \mf{new}^\mf{Succ}_p\ u
}
\end{mathpar}
where $m' = m \mid m[p].\mf{tree} = u$.
\begin{itemize}
\item Conclusions (\emph{ii}-a), (\emph{ii}-c), and \emph{(iii)} are all trivial.
\item To show conclusion (\emph{ii}-b), given premise (\emph{ii}-b),
it suffices to show that $\Gamma \semi \Pi \vdrt u: \mf{Tree}$.
This is given by (1) and we are done.\\
\hl{WE LOSE INFO HERE: work wf-judgment into memory?}
\end{itemize}

\tb{Case} \ts{ER-ScAttr}: We have
$\trip{\mc{D}}{m}{\mf{scan\ attr}(v_1, v_2)}
\cstep
\trip{\mc{D}}{m}{\mf{done}}$
and:
\begin{mathpar}
\inferrule*[right=TR-Sc] {
  \Gamma \semi \Pi \vdval v: \mf{Attr}
} {
  \Gamma \semi \Pi \vdrt \mf{scan\ attr}(v_1, v_2): \one_\rho
}
\end{mathpar}
\begin{itemize}
\item All conclusions are trivial.
\end{itemize}

\tb{Case} \ts{ER-ScClick}: We have
$\trip{\mc{D}}{m}{\mf{scan\ onClick}(v)}
\cstep
\trip{\mc{D}}{m}{\mf{handle}(v\ ())}$
and:
\begin{mathpar}
\inferrule*[right=TR-Sc] {
  \inferrule*[right=TV-OnClick] {
    \inferrule*[right=(1)] { } {
      \Gamma \semi \Pi \vdval v: \one \xrightarrow{E} \one
    }
  } {
    \Gamma \semi \Pi \vdval \mf{onClick}\ v: \mf{Attr}
  }
} {
  \Gamma \semi \Pi \vdrt \mf{scan\ onClick}(v): \one_\rho
}
\end{mathpar}
\begin{itemize}
\item Conclusions (\emph{ii}-a), (\emph{ii}-b), and \emph{(iii)} are all trivial.
\item We derive conclusion (\emph{ii}-c) as follows:
\begin{mathpar}
\inferrule*[right=TR-Hdl] {
  \inferrule*[right=TE-AppL] { 
    \inferrule*[right=(1)] { } {\Gamma \semi \Pi \vdval v: \one \xrightarrow{E} \one}
    \
    \inferrule*[right=TV-Unit] { } {\Gamma \semi \Pi \vdval (): \one}
  } {
    \Gamma \semi \Pi \vdexp v\ (): \one \mid E
  }
} {
  \Gamma \semi \Pi \vdrt \mf{handle}(v\ ()): \one_\rho
}
\end{mathpar}
\end{itemize}

\tb{Case} \ts{ER-ScsNil}: We have
$\trip{\mc{D}}{m}{\mf{scans}\ []}
\cstep
\trip{\mc{D}}{m}{\mf{done}}$
and:
\begin{mathpar}
\inferrule*[right=TR-Scs] {
  \Gamma \semi \Pi \vdval []: [\mf{Attr}]
} {
  \Gamma \semi \Pi \vdrt \mf{scans}\ []: \one_\rho
}
\end{mathpar}
\begin{itemize}
\item All conclusions are trivial.
\end{itemize}

\tb{Case} \ts{ER-ScsCons}: We have
$\trip{\mc{D}}{m}{\mf{scans}\ (v :: vs)}
\cstep
\trip{\mc{D}}{m}{\mf{scan}\ v \semi \mf{scans}\ vs}$
and:
\begin{mathpar}
\inferrule*[right=TR-Scs] {
  \inferrule*[right=TV-Cons] {
    \inferrule*[right=(1)] { } {\Gamma \semi \Pi \vdval v: \mf{Attr}}
    \and
    \inferrule*[right=(2)] { } {\Gamma \semi \Pi \vdval vs: [\mf{Attr}]}
  } {
    \Gamma \semi \Pi \vdval v :: vs: [\mf{Attr}]
  }
} {
  \Gamma \semi \Pi \vdrt \mf{scans}\ (v :: vs): \one_\rho
}
\end{mathpar}
\begin{itemize}
\item Conclusions (\emph{ii}-a), (\emph{ii}-b), and \emph{(iii)} are all trivial.
\item We derive conclusion (\emph{ii}-c) as follows:
\begin{mathpar}
\inferrule*[right=TR-Seq] {
  \inferrule*[right=TR-Sc] {
    \inferrule*[right=(1)] { } {\Gamma \semi \Pi \vdval v: \mf{Attr}}
  } {
    \Gamma \semi \Pi \vdrt \mf{scan}\ v: \one_\rho
  }
  \,
  \inferrule*[right=TR-Scs] {
    \inferrule*[right=(2)] { } {\Gamma \semi \Pi \vdval vs: [\mf{Attr}]}
  } {
    \Gamma \semi \Pi \vdrt \mf{scans}\ vs: \one_\rho
  }
} {
  \Gamma \semi \Pi \vdrt \mf{scan}\ v \semi \mf{scans}\ vs: \one_\rho
}
\end{mathpar}
\end{itemize}

\tb{Case} \ts{ER-Hdl}: We have
\begin{mathpar}
\inferrule*[right=ER-Hdl] {
  \inferrule*[right=(1)] { } {m \semi e \mapsto e' \semi m'}
} {
  \trip{\mc{D}}{m}{\mf{handle}\ e}
  \cstep
  \trip{\mc{D}}{m'}{\mf{handle}\ e'}
}
\and
\inferrule*[right=TR-Hdl] {
  \inferrule*[right=(2)] { } {\Gamma \semi \Pi \vdexp e: \one \mid E}
} {
  \Gamma \semi \Pi \vdrt \mf{handle}\ e: \one_\rho
}
\end{mathpar}
\begin{itemize}
\item Conclusions (\emph{ii}-a) and (\emph{iii}) are trivial.
\item We apply \thref{thm:prezexp} (expression preservation)
to (1), (2), and premise (\emph{ii}-b) to get:
\begin{enumerate}[start=3,label=(\arabic*)]
\item $\Gamma \semi \Pi \vdexp e': \one \mid E$ and
\item $\Gamma \semi \Pi \vdm m'$.
\end{enumerate}
\item (4) gives us conclusion (\emph{ii}-b) directly.
Conclusion (\emph{ii}-c) follows from the derivation:
\begin{mathpar}
\inferrule*[right=TR-Hdl] {
  \inferrule*[right=(3)] { } {\Gamma \semi \Pi \vdexp e': \one \mid E}
} {
  \Gamma \semi \Pi \vdrt \mf{handle}\ e': \one_\rho
}
\end{mathpar}
\end{itemize}

\tb{Case} \ts{ER-HdlDone}: We have
$\trip{\mc{D}}{m}{\mf{handle}\ \eret{()}}
\cstep
\trip{\mc{D}}{m}{\mf{done}}$ and:
\begin{mathpar}
\inferrule*[right=TR-Hdl] {
  \Gamma \semi \Pi \vdexp \eret{()}: \one \mid E
} {
  \Gamma \semi \Pi \vdrt \mf{handle}\ \eret{()}: \one_\rho
}
\end{mathpar}
\begin{itemize}
\item All conclusions are trivial.
\end{itemize}

\tb{Case} \ts{ER-ChkPF}: We have \begin{mathpar}
\inferrule*[right=ER-ChkPF] {
  m[p].\mf{rerender} = \mf{ff}
} {
  \trip{\mc{D}}{m}{\mf{check}\ p}
  \cstep
  \trip{\mc{D}}{m}{\mf{check}\ m[p].\mf{tree}}
}
\and
\inferrule*[right=TR-Chk] {
  \Gamma \semi \Pi \vdrt p: \mf{Tree}
} {
  \Gamma \semi \Pi \vdrt \mf{check}\ p: \one_\rho
}
\and
\inferrule*[right=WF-Chk] {
  \inferrule*[right=(WF-Path)] {
    \mc{L}(\hat{e}_0) = \dom(m[p].\mf{st})
    \and
    \inferrule*[right=(1)] { } {\mc{D} \semi m \wf m[p].\mf{tree}}
  } {
    \mc{D} \semi m \wf p
  }
} {
  \mc{D} \semi m \wf \mf{check}\ p
}
\end{mathpar}
\begin{itemize}
\item Conclusions (\emph{ii}-a) and (\emph{ii}-b) are trivial.
\item By premise (\emph{ii}-b), we get (2) that 
$\Gamma \semi \Pi \vdrt m[p].\mf{tree}: \mf{Tree}$.
Therefore, we can show conclusion (\emph{ii}-c):
\begin{mathpar}
\inferrule*[right=TR-Chk] {
  \inferrule*[right=(2)] { } {\Gamma \semi \Pi \vdrt m[p].\mf{tree}: \mf{Tree}}
} {
  \Gamma \semi \Pi \vdrt \mf{check}\ m[p].\mf{tree}: \one_\rho
}
\end{mathpar}
\item Finally, conclusion \emph{(iii)} follows from (1):
\begin{mathpar}
\inferrule*[right=WF-Chk] {
  \inferrule*[right=(1)] { } {\mc{D} \semi m \wf m[p].\mf{tree}}
} {
  \mc{D} \semi m \wf \mf{check}\ m[p].\mf{tree}
}
\end{mathpar}
\end{itemize}

\tb{Case} \ts{ER-ChkPT}: We have \begin{mathpar}
\inferrule*[right=ER-ChkPT] {
  m[p].\mf{rerender} = \mf{tt}
} {
  \trip{\mc{D}}{m}{\mf{check}\ p}
  \cstep
  \trip{\mc{D}}{m \mid m[p].\mf{rerender} = \mf{ff}}{\mf{eval}^\mf{Succ}_p\ \hat{e}_0[v_0/x]}
}
\and
\inferrule*[right=TR-Chk] {
  \inferrule*[right=(1)] { } {\Gamma \semi \Pi \vdrt m[p].\mf{tree}: \mf{Tree}}
} {
  \Gamma \semi \Pi \vdrt \mf{check}\ m[p].\mf{tree}: \one_\rho
}
\and
\inferrule*[right=WF-Chk] {
  \inferrule*[right=WF-Path] {
    \inferrule*[right=(2)] { } {\mc{L}(\hat{e}_0) = \dom(m[p].\mf{st})}
    \and
    \mc{D} \semi m \wf m[p].\mf{tree}
  } {
    \mc{D} \semi m \wf p
  }
} {
  \mc{D} \semi m \wf \mf{check}\ p
}
\end{mathpar}
\begin{itemize}
\item Conclusions (\emph{ii}-a) and (\emph{ii}-b) are trivial.
\item Recall that we defined $\hat{e}_0$ by $m[p].\mf{spec}=\pair{\mc{C}}{v_0}$
and $\mc{D}[\mc{C}] = x.\hat{e}_0$.
By premise (\emph{ii}-a), we have
(3) $\Gamma[x \mapsto \Delta[\mc{C}]] \semi \cdot \semi \mc{C} \vdtop \hat{e}_0$.
By premise (\emph{ii}-b), we also have
(4) $\Gamma \semi \Pi \vdval v_0: \Delta[\mc{C}]$.
\item By \thref{thm:wkpih} and \thref{thm:subtle} on (3) and (4),
we get $\Gamma \semi \cdot \semi \mc{C} \vdtop \hat{e}_0[v_0/x]$.
Noting that $\Pi[p]=\mc{C}$ by premise (\emph{ii}-b), we can now derive
conclusion (\emph{ii}-c):
\begin{mathpar}
\inferrule*[right=TR-Eval] {
  \Gamma \semi \Pi \semi \Pi[p] \vdtop \hat{e}_0[v_0/x]
} {
  \Gamma \semi \Pi \vdrt \mf{eval}^\mf{Succ}_p\ \hat{e}_0[v_0/x]: \one_\rho
}
\end{mathpar}
\item Conclusion \emph{(iii)} remains to be shown. By \thref{thm:sublab}
(label-set substitution), we have $\mc{L}(\hat{e}_0) = \mc{L}(\hat{e}_0[v_0/x])$,
and by (2), this means (5) that $\mc{L}(\hat{e}_0[v_0/x]) \subseteq \dom(m[p].\mf{st})$.
We conclude via derivation:
\begin{mathpar}
\and
\inferrule*[right=WF-EvalS] {
  \inferrule*[right=(2)] { } {\mc{L}(\hat{e}_0) = \dom(m[p].\mf{st})}
  \and
  \mc{L}(\hat{e}_0[v_0/x]) \subseteq \dom(m[p].\mf{st})
} {
  \mc{D} \semi m \wf \mf{eval}^\mf{Succ}_p\ \hat{e}_0[v_0/x]
}
\end{mathpar}
\end{itemize}

\tb{Case} \ts{ER-ChkLit}: Trivial.

\tb{Case} \ts{ER-ChkNode}: We have
$\trip{\mc{D}}{m}{\mf{check\ node}(\mf{id}, v, \rho)}
\cstep
\trip{\mc{D}}{m}{\mf{checks}\ \rho}$ and:
\begin{mathpar}
\inferrule*[right=TR-Chk] {
  \inferrule*[right=TV-Node] {
    \Gamma \semi \Pi \vdval v: [\mf{Attr}]
    \and
    \inferrule*[right=(1)] { } {\Gamma \semi \Pi \vdrt \rho: [\mf{Tree}]_\rho}
  } {
    \Gamma \semi \Pi \vdrt \mf{node}(\mf{id}, v, \rho): \mf{Tree}
  }
} {
  \Gamma \semi \Pi \vdrt \mf{check\ node}(\mf{id}, v, \rho): \one_\rho
}
\and
\inferrule*[right=WF-Chk] {
  \inferrule*[right=WF-Node] {
    \inferrule*[right=(2)] { } {\mc{D} \semi m \wf \rho}
  } {
    \mc{D} \semi m \wf \mf{node}(\mf{id}, v, \rho)
  }
} {
  \mc{D} \semi m \wf \mf{check\ node}(\mf{id}, v, \rho)
}
\end{mathpar}
\begin{itemize}
\item Conclusions (\emph{ii}-a) and (\emph{ii}-b) are trivial.
\item Conclusion (\emph{ii}-c) is derived as follows:
\begin{mathpar}
\inferrule*[right=TR-Chks] {
  \inferrule*[right=(1)] { } {\Gamma \semi \Pi \vdrt \rho: [\mf{Tree}]_\rho}
} {
  \Gamma \semi \Pi \vdrt \mf{checks}\ \rho: [\one_\rho]_\rho
}
\end{mathpar}
\item Conclusion \emph{(iii)} is derived as follows:
\begin{mathpar}
\inferrule*[right=WF-Chks] {
  \inferrule*[right=(2)] { } {\mc{D} \semi m \wf \rho}
} {
  \mc{D} \semi m \wf \mf{checks}\ \rho
}
\end{mathpar}
\end{itemize}

\tb{Case} \ts{ER-ChksNil}: Trivial.

\tb{Case} \ts{ER-ChksCons}: We have
$$\trip{\mc{D}}{m}{\mf{checks}\ \rho_1 ::_\rho \rho_2}
\cstep
\trip{\mc{D}}{m}{\mf{check}\ \rho_1 \semi \mf{checks}\ \rho_2}$$
\begin{mathpar}
\inferrule*[right=TR-Chks] {
  \inferrule*[right=TR-Cons] {
    \inferrule*[right=(1)] { } {\Gamma \semi \Pi \vdrt \rho_1: \mf{Tree}}
    \and
    \inferrule*[right=(2)] { } {\Gamma \semi \Pi \vdrt \rho_2: [\mf{Tree}]_\rho}
  } {
    \Gamma \semi \Pi \vdrt \rho_1 ::_\rho \rho_2: [\mf{Tree}]_\rho
  }
} {
  \Gamma \semi \Pi \vdrt \mf{checks}\ \rho_1 ::_\rho \rho_2 : \one_\rho
}
\and
\inferrule*[right=WF-Chks] {
  \inferrule*[right=WF-Cons] {
    \inferrule*[right=(3)] { } {\mc{D} \semi m \wf \rho_1}
    \and
    \inferrule*[right=(4)] { } {\mc{D} \semi m \wf \rho_2}
  } {
    \mc{D} \semi m \wf \rho_1 ::_\rho \rho_2
  }
} {
  \mc{D} \semi m \wf \mf{checks}\ \rho_1 ::_\rho \rho_2
}
\end{mathpar}
\begin{itemize}
\item Conclusions (\emph{ii}-a) and (\emph{ii}-b) are trivial.
\item Conclusion (\emph{ii}-c) is as follows:
\begin{mathpar}
\inferrule* {
  \inferrule*[right=TR-Chk] {
    \inferrule*[right=(1)] { } {\Gamma \semi \Pi \vdrt \rho_1: \mf{Tree}}
  } {
    \Gamma \semi \Pi \vdrt \mf{check}\ \rho_1: \one_\rho
  }
  \,
  \inferrule*[right=TR-Chks] {
    \inferrule*[right=(2)] { } {\Gamma \semi \Pi \vdrt \rho_2: [\mf{Tree}]_\rho}
  } {
    \Gamma \semi \Pi \vdrt \mf{checks}\ \rho_2: \one_\rho
  }
} {
  \Gamma \semi \Pi \vdrt \mf{check}\ \rho_1 \semi \mf{checks}\ \rho_2: \one_\rho
  \and (\ts{TR-Seq})
}
\end{mathpar}
\item Conclusion \emph{(iii)} is as follows:
\begin{mathpar}
\inferrule*[right=WF-Seq] {
  \inferrule*[right=WF-Chk] {
    \inferrule*[right=(3)] { } {\mc{D} \semi m \wf \rho_1}
  } {
    \mc{D} \semi m \wf \mf{check}\ \rho_1
  }
  \and
  \inferrule*[right=WF-Chk] {
    \inferrule*[right=(4)] { } {\mc{D} \semi m \wf \rho_2}
  } {
    \mc{D} \semi m \wf \mf{check}\ \rho_1
  }
} {
  \mc{D} \semi m \wf \mf{check}\ \rho_1 \semi \mf{checks}\ \rho_2
}
\end{mathpar}
\end{itemize}

\tb{Cases} to do with searching: \hl{unfinished}, but should be easy.

\end{Proof}
\end{lemma}

\pagebreak

\subsection{27 Jan}

My current "well-formedness" judgment is inadequate:
I need some condition on the \emph{memory}, not just the current term

Bigger question: what \emph{should} it mean for a memory or
a tree to be well-formed?

Intuitively: $\trip{\mc{D}}{m}{\mf{done}}$ is well-formed if, for all $p \in \dom(m)$:
\begin{itemize}
\item (define $\pair{\mc{C}}{v}=m[p].\mf{spec}$, $x.\hat{e}_0 = \mc{D}[\mc{C}]$)
\item $\mc{L}(\hat{e}_0) = \dom(m[p].\mf{spec})$ and
\item $m[p].\mf{tree}$ is well-formed
\end{itemize}
where a "well-formed tree" just means any paths in the tree are in memory?

Is this kind of thing sufficient?

How do I extend it for other terms? e.g. $\rho=\mf{new}^\mf{Init}_p \rho'$
means we'll \emph{later} assign a tree to $m[p]$, but right now it's uninitialized

What is my goal?

We can "realize" a tree (and a memory) into a concrete data structure as follows:
\begin{align*}
\mf{realize}(\mf{lit}\ v) &= \mf{lit}\ v\\
\mf{realize}(\mf{node}(\mf{id}, v, u)) &= \mf{node}(\mf{id}, v, \mf{realize}\ u)\\
\mf{realize}(p) &= \mf{realize}(m[p].\mf{tree})\\
\mf{realize}([]_\rho) &= []_\rho\\
\mf{realize}(u ::_\rho us) &= \mf{realize}(u) ::_\rho \mf{realize}(us)
\end{align*}
(In my current setup, to realize a memory, we would run $\mf{realize}(p_0)$.)

I want to show that $\mf{realize}(p_0)$ terminates 
at configuration $\trip{\mc{D}}{m}{\mf{done}}$.

Two issues:
\begin{itemize}
\item We want \tf{realize} to terminate (result in a finite tree)
\item We want \tf{realize} to never "get stuck", 
i.e. $m[p].\mf{tree}$ should always be defined.
\end{itemize}

\tf{realize} should terminate if the paths form a tree
(i.e. there are no reference cycles).

This is \emph{intuitively} true based on how the semantics unfolds:
we only create new paths at $\mf{init} \pair{\mc{C}}{v}$, and that path can't
"leave" its subterm (i.e. refer to things outside its subterm)

But it is \emph{syntactically} possible to write a term like:
$$\mf{new}^\mf{Succ}_{p_1} p_2 \semi \mf{new}^\mf{Succ}_{p_2} p_1$$
which assigns $m[p_1].\mf{tree} = p_2$ and $m[p_2].\mf{tree} = p_1$,
creating a cycle.

\begin{itemize}
\item How can I prove that such a situation will never arise
from some "valid" starting state?
\item Smaller point:
How do I actually go from "the paths form a tree" to "\tf{realize} terminates"?
\end{itemize}

\section{Big-Step Runtime (In Complete Disarray)}

\subsection{(Old) Runtime Typing}

\subsubsection{Trees}

\emph{\tb{Typing}}. Judgment: $\Pi \vdtree t$
\begin{mathpar}
\inferrule*[right=TTree-Lit] {
  \cdot \semi \Pi \vdval v: \mf{String}
} {
  \Pi \vdtree \mf{lit}\ v
}
\and
\inferrule*[right=TTree-Path] {
  p \in \dom(\Pi)
} {
  \Pi \vdtree p
}
\and
\inferrule*[right=TTree-Node] {
  \mf{id\ fresh}
  \and
  \cdot \semi \Pi \vdval v: \mf{String}
  \\
  (\forall i)\ \cdot \semi \Pi \vdval v_i: \mf{Attr}
  \and
  (\forall j)\ \Pi \vdtree t_j
} {
  \Pi \vdtree \mf{node}(\mf{id}, v, \num{v_i}, \num{t_j})
}
\end{mathpar}

\emph{\tb{Well-Formedness}}.
In addition to our typing judgment on trees,
we have a \emph{well-formedness} judgment on trees
with respect to a concrete memory $m$ and definition set $\mc{D}$.
Informally, this says that all paths in a tree have been
"properly initialized".

Judgment: $\mc{D} \semi m \wf t$
\begin{mathpar}
\inferrule*[right=WF-Lit] { } {\mc{D} \semi m \wf \mf{lit}\ v}
\and
\inferrule*[right=WF-Node] {
  (\forall j)
  \quad
  \mc{D} \semi m \wf t_j
} {
  \mc{D} \semi m \wf \mf{node}(\mf{id}, v, \num{v_i}, \num{t_j})
}
\and
\inferrule*[right=WF-Path] {
  m[p].\mf{spec} = \pair{\mc{C}}{v}
  \and
  \mc{D}[\mc{C}] = x.\hat{e}
  \\
  \mc{L}(\hat{e}) \subseteq \dom(m[p].\mf{st})
  \and
  \mc{D} \semi m \wf m[p].\mf{child}
} {
  \mc{D} \semi m \wf p
}
\end{mathpar}

\subsubsection{Component Context}
Define $\Delta \vDash \mc{D}$ to say that,
for all $\mc{C} \in \dom(\mc{D})$
(such that $\mc{D}[\mc{C}] = x.\hat{e}$),
we have $[x \mapsto \Delta[\mc{C}]] \semi \cdot \semi \mc{C} \vdtop \hat{e}$.
(Implicitly, $\dom(\mc{D}) \subseteq \dom(\Delta)$).

\subsubsection{Memory}
We say $\Gamma \semi \Pi \vdm m$ if, for all $p \in \dom(m)$,
\begin{itemize}
\item for $m[p].\mf{spec}=\pair{\mc{C}}{v}$,
$\Pi[p] = \mc{C}$ and $\Gamma \semi \Pi \vdval v: \Delta[\mc{C}]$, and
\item for all $\ell$ in $m[p].\mf{st}$,
$\Gamma \semi \Pi \vdval m[p].\mf{st}[\ell] : \Delta[\mc{C}][\ell]$.
\item if $m[p].\mf{child}$ is defined, then $\Pi \vdtree m[p].\mf{child}$.
\end{itemize}

\subsubsection{Modes}

Judgment: $\Pi \vdmu \mu$
\begin{mathpar}
\inferrule*[right=TMode-Idle] { } {\Pi \vdmu \mf{idle}}
\and
\inferrule*[right=TMode-Search] { } {\Pi \vdmu \mf{search}(id)}
\and
\inferrule*[right=TMode-Start] {
  \cdot \semi \Pi \vdprog P
} {
  \Pi \vdmu \mf{start}(P)
}
\and
\inferrule*[right=TMode-HandleNil] { } {\Pi \vdmu \mf{handle}([])}
\and
\inferrule*[right=TMode-HandleCons] {
  \cdot \semi \Pi \vdexp e : \one \mid E
  \and
  \Pi \vdmu \mf{handle}(es)
} {
  \Pi \vdmu \mf{handle}(e :: es)
}
\end{mathpar}

\subsubsection{Configurations}

Define a configuration $\four{\mc{D}}{m}{t}{\mu}$
to be well-typed via $\Pi \vdcfg \four{\mc{D}}{m}{t}{\mu}$ if:
\begin{itemize}
\item $\Delta \vDash \mc{D}$;
\item $\Pi \vdmu \mu$;
\item if $m$ is defined, then $\cdot \semi \Pi \vdm m$; and
\item if $t$ is defined, then both $\Pi \vdtree t$ and $\mc{D} \semi m \wf t$.
\end{itemize}

\subsection{Program Semantics}

Throughout the semantics, we use the notation $m \mid x = v$
to concisely denote nested map/record updates.
The above should be read as ``$m$ with field $x$ set to $v$''.
For instance, $m \mid m[p].\mf{st}[\ell] = v$ updates the
state field of $m$ at path $p$ and location $\ell$ to value $v$.

Again, we \hl{highlight} interesting rules.

\subsubsection{Expression Semantics}

Judgment: $m \semi e \mapsto e' \semi m'$
\begin{mathpar}
\inferrule*[right=E-Let] {
  m \semi e_1 \mapsto e_1' \semi m'
} {
  m \semi \elet{x}{e_1}{e_2} \mapsto \elet{x}{e_1'}{e_2} \semi m'
}
\end{mathpar}
\begin{align*}
m \semi \elet{x}{\eret{v}}{e} &\mapsto e[v/x] \semi m
&&\ts{E-Ret}\\
m \semi (\lambda x.e) v &\mapsto e[v/x] \semi m
&&\ts{E-AppFun}\\
m \semi \mc{C}\ v &\mapsto \eret{\pair{\mc{C}}{v}} \semi m
&&\ts{\hl{E-AppC}}\\
m \semi \vset{\ell}{p}\ v &\mapsto \eret{()} \semi m'
&&\ts{\hl{E-AppSet}}\\
(m' = m \mid m[p].\mf{st}[\ell] &= v\ 
    \mf{and}\ m[p].\mf{rerender} = \mf{tt})\\
m \semi \eif{\mf{true}}{e_1}{e_2} &\mapsto e_1 \semi m
&&\ts{E-IfTrue}\\
m \semi \eif{\mf{false}}{e_1}{e_2} &\mapsto e_2 \semi m
&&\ts{E-IfFalse}
\end{align*}

\subsubsection{Top-Level Expression Semantics}

Judgment: $m \semi \hat{e} \topstep{\phi}{p} \hat{e}' \semi m'$
\begin{mathpar}
\inferrule*[right=E-Lift] {
  \cdot \semi e \mapsto e' \semi \cdot
} {
  m \semi e \topstep{\phi}{p} e' \semi m
}
\and
\inferrule*[right=E-LetTop] {
  \cdot \semi e \mapsto e' \semi \cdot
} {
  m \semi \elet{x}{e}{\hat{e}} \topstep{\phi}{p} \elet{x}{e'}{\hat{e}} \semi m
}
\end{mathpar}
\begin{align*}
m \semi \elet{x}{\eret{v}}{\hat{e}} &\topstep{\phi}{p} \hat{e}[v/x] \semi m
&&\ts{E-RetTop}\\
m \semi \estate{x}{\ell}{v}{\hat{e}}
&\topstep{\mf{Init}}{p}
\hat{e}[v/x, \vset{\ell}{p}/x_\mf{set}] \semi m'
&&\ts{\hl{E-StInit}}\\
(m' &= m \mid m[p].\mf{st}[\ell] = v)\\
m \semi \estate{x}{\ell}{v}{\hat{e}}
&\topstep{\mf{Succ}}{p}
\hat{e}[v'/x, \vset{\ell}{p}/x_\mf{set}] \semi m
&&\ts{\hl{E-StSucc}}\\
(v' &= m[p].\mf{st}[\ell])
\end{align*}

\subsubsection{Program Expression Semantics}

Judgment: $\mc{D} \semi P \pstep P' \semi \mc{D}'$
\begin{mathpar}
\inferrule*[right=P-Lift] {
  \cdot \semi e \mapsto e' \semi \cdot
} {
  \mc{D} \semi e \pstep e \semi \mc{D}
}
\and
\inferrule*[right=P-Let] {
  \cdot \semi e \mapsto e' \semi \cdot
} {
  \mc{D} \semi \elet{x}{e}{P}
  \pstep
  \elet{x}{e'}{P} \semi \mc{D}
}
\end{mathpar}
\begin{align*}
\mc{D} \semi \elet{x}{\eret{v}}{P}
&\pstep
P[v/x] \semi \mc{D}
&&\ts{E-RetProg}\\
\mc{D} \semi \elet{\mc{C}(x)}{\hat{e}}{P}
&\pstep
P \semi \mc{D}[\mc{C} \mapsto x.\hat{e}]
&&\ts{\hl{E-CompDef}}
\end{align*}

\subsection{Event-Loop Semantics}

Say $m \semi \hat{e} \Topstep{\phi}{p} v \semi m'$ 
if $m \semi \hat{e} \xmapsto{*}^{\phi}_{p} \eret{v} \semi m'$.
\pagebreak
\subsubsection{Initialization}

Judgment: $\mc{D} \vdash v \bsinit{n} \pair{t}{m}$.

Define $\pi_0(\mc{C}, v)$ as $
\{\mf{spec}: \pair{\mc{C}}{v}, \mf{st}: \emptyset, \mf{child}: \emptyset,
\mf{rerender}: \mf{ff}\}$.

\begin{mathpar}
\inferrule*[right=\ts{Init-Lit}] {
  n > 0
} {
  \mc{D} \vdash \mf{text}\ v 
  \bsinit{n}
  \pair{\mf{lit}\ v}{\cdot}
}
\and
\inferrule*[right=\ts{Init-Node}] {
  n > 0
  \and
  \mf{id\ fresh}
  \and
  (\forall j)
  \quad
  \mc{D} \vdash v_j
  \bsinit{n-1}
  \pair{t_j}{m_j}
} {
  \mc{D} \vdash \mf{tag}(v, \num{v_i}, \num{v_j})
  \bsinit{n}
  \pair{\mf{node}(\mf{id}, v, \num{v_i}, \num{t_j})}{\bigcup_j m_j}
}
\and
\inferrule*[right=\ts{Init-Comp}] {
  n > 0
  \and
  p\ \mf{fresh}
  \and
  \mc{D}[\mc{C}] = x.\hat{e}
  \\
  [p \mapsto \pi_0(\mc{C}, v)] 
  \semi \hat{e}[v/x] \Topstep{\mf{Init}}{p} v' \semi m 
  \\
  \mc{D} \vdash v' \bsinit{n-1} \pair{t}{m'}
} {
  \mc{D} \vdash \pair{\mc{C}}{v}
  \bsinit{n}
  \pair{p}{m' \uplus (m \mid m[p].\mf{child} = t)}
}
\and
\inferrule*[right=\ts{Init-Div}] { } {
  \mc{D} \vdash \bsinit{0} \bot
}
\and
\inferrule*[right=\ts{Init-Node}$_\bot$] {
  n > 0
  \and
  (\exists j)
  \quad
  \mc{D} \vdash v_j \bsinit{n-1} \bot
} {
  \mc{D} \vdash \mf{node}(\mf{id}, v, \num{v_i}, \num{v_j}) \bsinit{n} \bot
}
\and
\inferrule*[right=\ts{Init-Comp}$_\bot$] {
  n > 0
  \and
  p\ \mf{fresh}
  \and
  \mc{D}[\mc{C}] = x.\hat{e}
  \\
  [p \mapsto \pi_0(\mc{C}, v)] 
  \semi \hat{e}[v/x] \Topstep{\mf{Init}}{p} v' \semi m 
  \\
  \mc{D} \vdash v' \bsinit{n-1} \bot
} {
  \mc{D} \vdash \pair{\mc{C}}{v}
  \bsinit{n}
  \pair{p}{m' \uplus (m \mid m[p].\mf{child} = t)}
}
\end{mathpar}

\subsubsection{Rerendering}

Judgment: $\mc{D} \semi m \vdash t \bssucc{n} m'$

\begin{mathpar}
\inferrule*[right=Check-Lit] { } {
  \mc{D} \semi m \vdash \mf{lit}\ v \bssucc{n} \cdot
}
\and
\inferrule*[right=Check-Node] {
  (\forall j) \quad
  \mc{D} \semi m \vdash t_j \bssucc{n} m_j
} {
  \mc{D} \semi m \vdash \mf{node}(\mf{id}, v, \num{v_i}, \num{t_j}) 
  \bssucc{n} \bigcup_{\forall j} m_j
}
\and
\inferrule*[right=\hl{Check-PathF}] {
  m[p].\mf{rerender} = \mf{ff}
  \\
  \mc{D} \semi m \vdash m[p].\mf{child} \bssucc{n} m'
} {
  \mc{D} \semi m \vdash p \bssucc{n} m'
}
\and
\inferrule*[right=\hl{Check-PathT}] { 
  m[p].\mf{rerender} = \mf{tt}
  \\
  m'[p].\mf{spec} = \pair{\mc{C}}{v}
  \and
  \mc{D}[\mc{C}] = x.\hat{e}
  \\
  m \semi \hat{e}[v/x] \Topstep{\mf{Succ}}{p} \eret{v'} \semi \_
  \\
  \mc{D} \semi m \vdash v' \bsinit{n} \pair{t}{m'}
} {
  \mc{D} \semi m \vdash p \bssucc{n} m' \mid 
  m'[p].\mf{child}=t\ \mf{and}\ m'[p].\mf{rerender}=\mf{ff}
}
\and
\inferrule*[right=\hl{Check-PathDiv}] { 
  m[p].\mf{rerender} = \mf{tt}
  \\
  m'[p].\mf{spec} = \pair{\mc{C}}{v}
  \and
  \mc{D}[\mc{C}] = x.\hat{e}
  \\
  m \semi \hat{e}[v/x] \Topstep{\mf{Succ}}{p} \eret{v'} \semi \_
  \\
  \mc{D} \semi m \vdash v' \bsinit{n} \bot
} {
  \mc{D} \semi m \vdash p \bssucc{n} \bot
}
\and
\inferrule*[right=Check-Node$_\bot$] {
  (\exists j)
  \quad
  \mc{D} \semi m \vdash t_j \bssucc{n} \bot
} {
  \mc{D} \semi m \vdash \mf{node}(\mf{id}, v, \num{v_i}, \num{t_j}) \bssucc{n} \bot
}
\and
\inferrule*[right=Check-PathF$_\bot$] {
  m[p].\mf{rerender} = \mf{ff}
  \\
  \mc{D} \semi m \vdash m[p].\mf{child} \bssucc{n} \bot
} {
  \mc{D} \semi m \vdash p \bssucc{n} \bot
}
\end{mathpar}

\subsubsection{Handler Search}

Judgment: $m \semi \mf{id} \vdash \mf{handlers}(t) = \num{v}$

\begin{mathpar}
\inferrule*[right=Handlers-Lit] { } {
  m \semi \mf{id} \vdash \mf{handlers}(\mf{lit}\ v) = \{\}
}
\and
\inferrule*[right=Handlers-Comp] {
  m \semi \mf{id} \vdash \mf{handlers}(m[p].\mf{child}) = \num{v}
} {
  m \semi \mf{id} \vdash \mf{handlers}(p) = \num{v}
}
\and
\inferrule*[right=Handlers-Tgt] { } {
  m \semi \mf{id} \vdash \mf{handlers}(\mf{node}(id, v, \num{v_i}, \num{v_j}))
  = \{v_h \mid \mf{onClick}\ v_h \in \num{v_i}\}
}
\and
\inferrule*[right=Handlers-Node] {
  \mf{id} \neq \mf{id}'
  \and
  (\forall j)\ 
  m \semi \mf{id}\vdash \mf{handlers}(v_j) = h_j
} {
  m \semi \mf{id} \vdash \mf{handlers}(\mf{node}(id', v, \num{v_i}, \num{v_j}))
  = \bigcup_{\forall j} h_j
}
\end{mathpar}

\subsubsection{Event Loop}

Judgment: $\four{\mc{D}}{m}{t}{\mu} \cstep \four{\mc{D}'}{m'}{t}{\mu'}$

\begin{mathpar}
\inferrule*[right=\hl{Loop-Start}] {
  \mc{D} \semi P \pstep P' \semi \mc{D}'
} {
  \four{\mc{D}}{m}{t}{\mf{start}(P)} 
  \cstep
  \four{\mc{D}'}{m}{t}{\mf{start}(P')}
}
\and
\inferrule*[right=\hl{Loop-Init}] {
  (\forall n \in \mathbb{N})
  \quad
  \mc{D} \vdash v \bsinit{n} \pair{t}{m}
} {
  \four{\mc{D}}{\_}{\_}{\mf{start}(\eret{v})} 
  \cstep 
  \four{\mc{D}}{m}{t}{\mf{idle}}
}
\and
\inferrule*[right=\hl{Loop-Event}] { } {
  \four{\mc{D}}{m}{t}{\mf{idle}} 
  \cstep 
  \four{\mc{D}}{m}{t}{\mf{search}(id)}
}
\and
\inferrule*[right=\hl{Loop-Search}] {
  m \semi \mf{id} \vdash \mf{handlers}(t) = \num{v}
} {
  \four{\mc{D}}{m}{t}{\mf{search}(id)}
  \cstep
  \four{\mc{D}}{m}{t}{\mf{handle}(\num{v\ ()})}
}
\and
\inferrule*[right=\hl{Loop-HandleCong}] {
  m \semi e \mapsto e' \semi m'
} {
  \four{\mc{D}}{m}{t}{\mf{handle}(e :: es)}
  \cstep
  \four{\mc{D}}{m'}{t}{\mf{handle}(e' :: es)}
}
\and
\inferrule*[right=\hl{Loop-HandleNext}] { } {
  \four{\mc{D}}{m}{t}{\mf{handle}(\eret{()} :: es)}
  \cstep
  \four{\mc{D}}{m}{t}{\mf{handle}(es)}
}
\and
\inferrule*[right=\hl{Loop-Rerender}] {
  (\forall n \in \mathbb{N})
  \quad
  \mc{D} \semi m \vdash t \bssucc{n} m'
} {
  \four{\mc{D}}{m}{t}{\mf{handle}([])}
  \cstep
  \four{\mc{D}}{m'}{t}{idle}
}
\end{mathpar}

\subsection{(Old) Runtime Type Soundness}

\subsubsection{Initialization}

\begin{lemma}[Progress of initialization]
If $\Delta \vDash \mc{D}$ and $\cdot \semi \Pi \vdval v: \mf{View}$,
then, for all $n \in \mathbb{N}$,
either $\mc{D} \vdash v \divinit{n}$ or $\mc{D} \vdash v \bsinit{n} \pair{t}{m}$.
\end{lemma}

\begin{proof}
By induction on $n$.

In the base case, when $n=0$, we can always apply rule \ts{Init-Div}
to conclude that $\mc{D} \vdash v \divinit{0}$ and we are done.

In the inductive case, we have $n>0$ by assumption.
We apply the canonical forms lemma to the typing premise
and case-analyze the result.

\tb{Case} $v = \mf{text}\ v'$: 
We simply apply $\ts{Init-Lit}$ to get 
$\mc{D} \vdash \mf{text}\ v' \bsinit{n} \pair{\mf{lit}\ v'}{\cdot}$
and we are done.

\tb{Case} $v = \mf{tag}(v', \num{v_i}, \num{v_j})$: unfinished

\tb{Case} $v = \pair{\mc{C}}{v}$: unfinished

\end{proof}

\begin{lemma}[Preservation of initialization]
If \begin{enumerate}[label=(\roman*)]
\item $\Delta \vDash \mc{D}$, 
\item $\cdot \semi \Pi \vdval v: \mf{View}$, and
\item $\mc{D} \vdash v \bsinit{n} \pair{t}{m}$,
\end{enumerate}
then there exists some $\Pi'$ with \begin{enumerate}[label=(\roman*)]
\item $\Pi \po \Pi'$, 
\item $\cdot \semi \Pi' \vdm m$, and
\item $\Pi' \vdtree t$. We also have
\item $\mc{D} \semi m \wf t$.
\end{enumerate}
\end{lemma}

\begin{proof}
By induction on the derivation of the big-step judgment.
\end{proof}

\begin{theorem}[Progress of configurations]
If $\Pi \vdcfg \four{\mc{D}}{m}{t}{\mu}$, either:
\begin{enumerate}[label=(\roman*)]
\item $\mu = \mf{idle}$,
\item $\mu = \mf{start}(\eret{v})$ and $\mc{D} \vdash v \divinit{ }$,
\item $\mu = \mf{handle}([])$ and $\mc{D} \semi m \vdash t \divsucc{ }$, or
\item $\four{\mc{D}}{m}{t}{\mu} \cstep \four{\mc{D}'}{m'}{t'}{\mu'}$.
\end{enumerate}
\end{theorem}

\begin{theorem}[Preservation of configurations]
If $\Pi \vdcfg \four{\mc{D}}{m}{t}{\mu}$\\
and $\four{\mc{D}}{m}{t}{\mu} \cstep \four{\mc{D}'}{m'}{t'}{\mu'}$,
there exists some $\Pi'$ with $\Pi \po \Pi'$
and $\Pi' \vdcfg \four{\mc{D}'}{m'}{t'}{\mu'}$.
    
\end{theorem}

\subsubsection{Initialization (Even Older)}

We first prove a number of lemmas which are necessary in our critical property.

\begin{lemma}[Path Preservation]
If $m \semi \hat{e} \topstep{\phi}{p} \hat{e}' \semi m'$
and $p \in \dom(m)$,
then $\dom(m) = \dom(m')$.
\end{lemma}

\begin{proof}
By induction on the derivation of the small-step judgment.

\tb{Case} \ts{E-StInit}: \begin{itemize}
\item We have
$m \semi \estate{x}{\ell}{v}{\hat{e}} \topstep{\mf{Init}}{p}
\hat{e}[v/x, \vset{\ell}{p}/x_\mf{set}] \semi m'$
where $m' = m \mid m[p].\mf{st}[\ell] = v$.
\item Since we assumed $p \in \dom(m)$, the addition of $p$
to $m'$ does not change its domain, and therefore we have $\dom(m)=\dom(m')$.
\end{itemize}

\tb{All other cases} do not alter $m$, and the lemma holds trivially.

\end{proof}

\begin{corollary}
If $m \semi \hat{e} \Topstep{\phi}{p} v \semi m'$
and $p \in \dom(m)$,
then $\dom(m) = \dom(m')$.
\end{corollary}

\begin{proof}
Immediate from path preservation and definition of $\Topstep{\phi}{p}$.
\end{proof}

\begin{lemma}[Memory Freshness]
If $\mc{D} \vdash v \bsinit{n} \pair{t}{m}$,
then for all $p \in \dom(m)$, we have $p\ \mf{fresh}$.
\end{lemma}

\begin{proof}
By induction over the derivation of the big-step judgment.

\tb{Case} \ts{Init-Lit}: $m=\cdot$, so the conclusion holds trivially.

\tb{Case} \ts{Init-Node}: \begin{itemize}
\item We have
$\mc{D} \vdash \mf{tag}(v, \num{v_i}, \num{v_j}) 
\bsinit{n}
\pair{\mf{node}(\mf{id}, v, \num{v_i}, \num{t_j})}{\bigcup_j m_j}$.
\item For all $j$, we also have $\mc{D} \vdash v_j \bsinit{n} \pair{t_j}{m_j}$.
\item We apply the inductive hypothesis to each, getting that
all $m_j$ only contain fresh $p$.
\item Therefore, their union $\bigcup_j m_j$ also only contains fresh $p$,
and we are done.
\end{itemize}

\tb{Case} \ts{Init-Comp}: \begin{itemize}
\item We start with $\mc{D} \vdash \pair{\mc{C}}{v} 
\bsinit{n} \pair{p}{m' \cup (m \mid m[p].\mf{child} = t)}$.
\item Let $\mc{D}[\mc{C}] = x.\hat{e}$ and
$m_0 = \{\mf{spec}: \pair{\mc{C}}{v_0}, \mf{st}: \emptyset, \mf{child}: \emptyset\}$.
\item We have the following premises:
\begin{enumerate}[label=(\alph*)]
  \item $p$ is fresh,
  \item $m_0 \semi \hat{e}[v_0/x] \Topstep{\mf{Init}}{p} v \semi m$, and
  \item $\mc{D} \vdash v \bsinit{n} \pair{t}{m'}$.
\end{enumerate}
\item By (a), clearly $m_0$ has all $p$ fresh.
\item By the path preservation corollary on (b), $m$ also has all $p$ fresh.
\item By inductive hypothesis on $c$, $m'$ also has all $p$ fresh.
\item Therefore, $m' \cup (m \mid m[p].\mf{child} = t)$
has all $p$ fresh and we are done.
\end{itemize}
\end{proof}

\begin{lemma}[Label Initialization]
If $m \semi \hat{e} \topstep{\mf{Init}}{p} \hat{e}' \semi m'$,
then $\mc{L}(\hat{e}) \setminus \mc{L}(\hat{e}') \subseteq \dom(m'[p].\mf{st})$.
\end{lemma}

\begin{proof}
By case-analysis of the small-step judgment. Note that if 
$\mc{L}(\hat{e}) \setminus \mc{L}(\hat{e}') = \emptyset$,
the conclusion holds trivially.

\tb{Case} \ts{E-Lift}: \begin{itemize}
\item We have $\cdot \semi e \mapsto e' \semi \cdot$.
\item By definition, $\mc{L}(e)=\emptyset=\mc{L}(e')$.
\item Clearly then $\mc{L}(e) \setminus \mc{L}(e')=\emptyset$ and we are done.
\end{itemize}

\tb{Case} \ts{E-LetTop}: \begin{itemize}
\item We have $m \semi \elet{x}{e}{\hat{e}} \topstep{\mf{Init}}{p} \elet{x}{e'}{\hat{e}} \semi m$,
with $\cdot \semi e \mapsto e' \semi \cdot$.
\item By definition, $\mc{L}(\elet{x}{e}{\hat{e}})=\mc{L}(\elet{x}{e'}{\hat{e}})=\mc{L}(\hat{e})$.
\item We get $\mc{L}(\elet{x}{e}{\hat{e}}) \setminus \mc{L}(\elet{x}{e'}{\hat{e}})=\emptyset$ and we are done.
\end{itemize}

\tb{Case} \ts{E-RetTop}: \begin{itemize}
\item We have $m \semi \elet{x}{\eret{v}}{\hat{e}} 
\topstep{\mf{Init}}{p} \hat{e}[v/x] \semi m$.
\item By definition, $\mc{L}(\elet{x}{\eret{v}}{\hat{e}})=\mc{L}(\hat{e})$.
\item By label preservation over substitution, 
$\mc{L}(\hat{e}[v/x])=\mc{L}(\hat{e})$.
\item We get $\mc{L}(\elet{x}{\eret{v}}{\hat{e}}) 
\setminus \mc{L}(\hat{e}[v/x])=\emptyset$ and we are done.
\end{itemize}

\tb{Case} \ts{E-StInit}: \begin{itemize}
\item We have $m \semi \estate{x}{\ell}{v}{\hat{e}} 
\topstep{\mf{Init}}{p} \hat{e}[v/x, \vset{\ell}{p}/x_\mf{set}]
\semi m'$, where $m' = m \mid m[p].\mf{st}[\ell] = v$.
\item By definition, $\mc{L}(\estate{x}{\ell}{v}{\hat{e}}) 
= \{\ell\} \cup \mc{L}(\hat{e})$.
\item By label preservation over substitution,
$\mc{L}(\hat{e}[v/x, \vset{\ell}{p}/x_\mf{set}]) = \mc{L}(\hat{e})$.
\item We subcase on whether $\ell \in \mc{L}(\hat{e})$.
\item Subcase 1: $\ell \in \mc{L}(\hat{e})$.
\begin{itemize}
  \item It follows that $\mc{L}(\estate{x}{\ell}{v}{\hat{e}}) = \mc{L}(\hat{e})$.
  \item Thus, $\mc{L}(\estate{x}{\ell}{v}{\hat{e}}) \setminus
\mc{L}(\hat{e}[v/x, \vset{\ell}{p}/x_\mf{set}]) = \emptyset$ and we are done.
\end{itemize}
\item Subcase 2: $\ell \notin \mc{L}(\hat{e})$.
\begin{itemize}
  \item Then $\mc{L}(\estate{x}{\ell}{v}{\hat{e}}) \setminus \mc{L}(\hat{e}) = \{\ell\}$.
  \item Thus, we need to show $\ell \in m'[p].\mf{st}$.
  \item This is given by definition of $m'$, and we are done.
\end{itemize}
\end{itemize}

\tb{Case} \ts{E-StSucc}: Impossible, 
since this case requires $\phi=\mf{Succ}$,
which contradicts our assumption that $\phi=\mf{Init}$.

\end{proof}

\begin{corollary}
If $m \semi \hat{e} \Topstep{\mf{Init}}{p} v \semi m'$,
then $\mc{L}(\hat{e}) \subseteq \dom(m'[p].\mf{st})$.
\end{corollary}

\begin{proof}
Follows from label initialization and the fact that 
$\mc{L}(\eret{v})=\emptyset$, which means that 
$\mc{L}(\hat{e}) \setminus \mc{L}(\eret{v}) = \mc{L}(\hat{e})$.
\end{proof}

\begin{lemma}[Strong Disjoint Memory]
If \begin{enumerate}[label=(\roman*)]
\item $\Gamma \semi \Pi \semi \mc{D} \vdm m$,
\item $\Gamma \semi \Pi \semi \mc{D} \vdm m'$, and
\item $\dom(m) \cap \dom(m') = \emptyset$,
\end{enumerate}
then $\Gamma \semi \Pi \semi \mc{D} \vdm m \cup m'$.
\end{lemma}

\begin{proof}
Follows easily from definition of $\vdm$.
\end{proof}

\begin{theorem}[Critical Property of Initialization]
If \begin{enumerate}[label=(\roman*)]
\item $\mc{D} \vdash v \bsinit{n} \pair{t}{m}$,
\item $\Delta \vDash \mc{D}$, and
\item $\cdot \semi \Pi \vdval v: \mf{View}$,
\end{enumerate}
then there exists some $\Pi'$ such that: \begin{enumerate}[label=(\roman*)]
\item $\Pi \sqsubseteq \Pi'$,
\item $\Pi' \vdtree t$, and
\item $\cdot \semi \Pi' \semi \mc{D} \vdm m$.
\end{enumerate}
\end{theorem}

\begin{proof}
By induction on derivation of \emph{(i)}.

\tb{Case} \ts{Init-Lit}: \begin{itemize}
\item We have $\mc{D} \vdash \mf{text}\ v \bsinit{n} \pair{\mf{lit}\ v}{\cdot}$.
\item Choose $\Pi'=\Pi$; conclusions \emph{(i)} and \emph{(iii)} are trivial.
\item By premise \emph{(iii)} and \ts{TV-Text},\\
we have $\cdot \semi \Pi \vdval v: \mf{String}$.
\item We conclude via \ts{TTree-Lit} that $\Pi \vdtree \mf{lit}\ v$
and we are done. 
\end{itemize}

\tb{Case} \ts{Init-Node}: \begin{itemize}
\item We have $\mc{D} \vdash \mf{tag}(v, \num{v_i}, \num{v_j})
  \bsinit{n}
  \pair{\mf{node}(\mf{id}, v, \num{v_i}, \num{t_j})}{\bigcup_j m_j}$.
\item Furthermore, for all $j$, we have (a) that $\mc{D} \vdash v_j \bsinit{n} \pair{t_j}{m_j}$.
\item By \emph{(iii)} and \ts{TV-Tag}, we have \begin{enumerate}[label=(\alph*)]
  \stepcounter{enumi}
  \item $\cdot \semi \Pi \vdval v: \mf{String}$,
  \item $\cdot \semi \Pi \vdval v_i: \mf{Attr}$ for all $i$, and
  \item $\cdot \semi \Pi \vdval v_j: \mf{View}$ for all $j$.
\end{enumerate}
\item We \emph{iteratively} apply the inductive hypothesis for each $j$.
\begin{itemize}
  \item Let $j$ range from $0$ to $n$ (exclusive).
  \item Let $\Pi_0=\Pi$.
  \item At step $j$, we weaken (d) to get (e) that $\cdot \semi \Pi_j \vdval v_j: \mf{View}$.
  \begin{itemize}
    \item That $\Pi \po \Pi_j$ is given by reflexivity at $j=0$, 
    or the previous conclusion (f) \& transitivity of $\po$ otherwise.
    (We have a chain $\Pi_0 \po \Pi_1 \po \cdots \po \Pi_j$).
  \end{itemize}
  \item We apply the inductive hypothesis using (a), (e), and premise \emph{(ii)}.
  \item We get some $\Pi_{j+1}$ with:
  \begin{enumerate}[label=(\alph*)]
    \setcounter{enumi}{5}
    \item $\Pi_j \po \Pi_{j+1}$,
    \item $\Pi_{j+1} \vdtree t_j$, and
    \item $\cdot \semi \Pi_{j+1} \semi \mc{D} \vdm m_{j}$.
  \end{enumerate}
\end{itemize}
\item Choose $\Pi'=\Pi_n$.
\item By (f) and transitivity of $\po$, we get $\Pi \po \Pi_n$, i.e. conclusion \emph{(i)}.
\item Since $\Pi_{j+1} \po \Pi_n$ for all $j$, weaken (g) and (h) to say
\begin{enumerate}[label=(\alph*')]
  \setcounter{enumi}{6}
  \item $\Pi_n \vdtree t_j$ and
  \item $\cdot \semi \Pi_{n} \semi \mc{D} \vdm m_j$.
\end{enumerate}
\item By the memory freshness lemma on (a), we get that all $p \in m_j$ are fresh,
i.e. all $m_j$ are disjoint.
\item Therefore, by the strong disjoint memory lemma,
we get by (h$'$) that
$\cdot \semi \Pi_n \vdm \bigcup_j m_j$, i.e. conclusion \emph{(iii)}.
\item Finally, we apply \ts{TTree-Node} using (b) and (c) weakened on $\Pi$,
as well as (g$'$),
to get $\Pi_n \vdtree \mf{node}(\mf{id}, v, \num{v_i}, \num{t_j})$, 
i.e conclusion \emph{(ii)}, and we are done.
\end{itemize}

\tb{Case} \ts{Init-Comp}: \begin{itemize}
\item We start with $\mc{D} \vdash \pair{\mc{C}}{v_0}
  \bsinit{n}
  \pair{p}{m' \cup (m \mid m[p].\mf{child} = t)}$.
\item Let $\mc{D}[\mc{C}] = x.\hat{e}$ and
$m_0 = \{\mf{spec}: \pair{\mc{C}}{v_0}, \mf{st}: \emptyset, \mf{child}: \emptyset\}$.
\item We have the following premises:
\begin{enumerate}[label=(\alph*)]
  \item $p$ is fresh,
  \item $m_0 \semi \hat{e}[v_0/x] \Topstep{\mf{Init}}{p} v \semi m$ (S3), and
  \item $\mc{D} \vdash v \bsinit{n} \pair{t}{m'}$.
\end{enumerate}
\item Let $\Pi_0=\Pi[p \mapsto \mc{C}]$; by freshness of $p$, $\Pi \po \Pi_0$.
Note that $\cdot \semi \Pi_0 \vdm m_0$ (S2) by definition.
\item By premise \emph{(ii)}, we get 
$[x \mapsto \Delta[\mc{C}]] \semi \cdot \semi \mc{C} \vdtop \hat{e}$.
\item Premise \emph{(iii)} gives the following derivation (bottom-up):
\begin{mathpar}
\inferrule*[right=TV-Spec] {
  \inferrule*[right=TV-Comp] {
    \inferrule*{(d)} {\Delta[\mc{C}] = \tau}
  } {
    \cdot \semi \Pi \vdval \mc{C}: \mf{Component}\ \tau
  }
  \and
  \inferrule*{(e)} {\cdot \semi \Pi \vdval v_0: \tau}
} {
  \cdot \semi \Pi \vdval \pair{\mc{C}}{v_0} : \mf{View}
}
\end{mathpar}
\item By (d), $[x \mapsto \tau] \semi \cdot \semi \mc{C} \vdtop \hat{e}$.
\item By (e), weakening, and substitution lemma,
$\cdot \semi \Pi \semi \mc{C} \vdtop \hat{e}[v/x]$.
\item We weaken this again to get $\cdot \semi \Pi_0 \semi \mc{C} \vdtop \hat{e}[v/x]$ (S1).
\item We can therefore apply the soundness of top-level expressions at initial phase
to (S1-3), and get that:
\begin{enumerate}[label=(\alph*)]
  \setcounter{enumi}{5}
  \item $\cdot \semi \Pi_0 \vdval v: \mf{View}$, and
  \item $\cdot \semi \Pi_0 \vdm m$.
\end{enumerate}
\item We then apply the inductive hypothesis to (c) and (f) to get some $\Pi_1$ with:
\begin{enumerate}[label=(\alph*)]
  \setcounter{enumi}{7}
  \item $\Pi_0 \po \Pi_1$,
  \stepcounter{enumi}
  \item $\Pi_1 \vdtree t$, and
  \item $\cdot \semi \Pi_1 \semi \mc{D} \vdm m'$.
\end{enumerate}
\item Choose this $\Pi_1$ as our $\Pi'$. We now need to show each of the conclusions.
\item That $\Pi \po \Pi'$ follows from (h), definition of $\Pi_0$,
and transitivity of $\po$. (Conclusion \emph{(i)} satisfied.)
\item That $\Pi' \vdtree p$ (conclusion \emph{(ii)}) 
is immediate from \ts{TTree-Path} and $\Pi_0 \po \Pi'$.
\item We then need to show conclusion \emph{(iii)}, that 
$\cdot \semi \Pi' \semi \mc{D} \vdm m' \cup (m \mid m[p].\mf{child} = t)$.
\item Let $m_1 = m \mid m[p].\mf{child} = t$.
By (k) and the strong disjoint memory lemma, it suffices to show
that $\cdot \semi \Pi' \semi \mc{D} \vdm m_1$.
\item By definition of $\vdm$, we need the following sub-conclusions:
\begin{enumerate}[label=(iii-\roman*)]
  \item $\cdot \semi \Pi' \vdm m_1$,
  \item $\Delta \vDash \mc{D}$,
  \item for all $p \in \dom(m_1)$, for $\mc{D}[\Pi[p]]=x.\hat{e}$,
  $\mc{L}(\hat{e}) \subseteq \dom(m_1[p].\mf{st})$, and
  \item for all $p \in \dom(m_1)$, $\Pi' \vdtree m_1[p].\mf{child}$.
\end{enumerate}
\item (\emph{iii}-a) follows from (g) and weakening via (h).
\item (\emph{iii}-b) is immediate by premise \emph{(ii)}.
\item For (\emph{iii}-c) and (\emph{iii}-d), we apply the path preservation lemma
to (b) to say that $\dom(m_0) = \dom(m)$.
\item By definition of each of $m_0$ and $m_1$, we get $\dom(m_1)=\{p\}$.
\item Therefore, (\emph{iii}-d) is satisfied by (j).
\item Finally, the label initialization corollary gives us that\\
$\mc{L}(\hat{e}) \subseteq \dom(m[p].\mf{st})$, 
i.e. sub-conclusion (\emph{iii}-c), and we are done.
\end{itemize}
\end{proof}

\section{$\lambda_{CC}$}

\subsection{Syntax}

Additions/changes from React-tRace Core are \hl{highlighted}.
Intermediate-only constructs are \grey{greyed}
\begin{align*}
\text{Value } v &::= x \mid () \mid \mf{true} \mid \mf{false} 
\mid s \mid \lambda x: \tau.e\\
&\mid \mf{text}\ v \mid \mf{tag}(v, \num{v}, \num{v})
\mid \mf{attr}(v, v) \mid \mf{onClick}\ v \mid \mathgrey{\ell}\\
\text{Expression } e &::= \eret{v} \mid \elet{x}{e}{e} 
\mid v\ v \mid \eif{v}{e}{e}\\
&\mid \mhl{\eref{v}} \mid \mhl{\eget{e}} \mid \mhl{\eset{v}{v}}\\
\text{Type } \tau &::= \one \mid \mf{Bool} \mid \mhl{\tau \rightarrow \tau}
\mid \mf{View} \mid \mf{Attr} \mid \mhl{\mf{Ref}\ \tau}\\
\end{align*}
Internal representations:
\begin{align*}
\text{Type context } \Gamma &::= \map{x}{\tau}\\
\text{Store } \sigma &::= \map{\ell}{v}\\
\text{Tree }t &::= \mf{lit}\ v \mid \mf{node}(id, v, \num{v}, \num{t})\\
\text{Mode }\mu &::= \mf{start}(e) \mid \mf{idle} \mid \mf{render}
\end{align*}

\subsection{Typing}

Non-standard rules are \hl{highlighted}, although on the whole,
they are the same as in React-tRace Core. (Most relevant are the omissions.)

\pagebreak

\subsubsection{Value Typing}

Judgment: $\Gamma \vdashv v: \tau$

\begin{mathpar}
\inferrule*[right=TVal-Var] {
  x : \tau \in \Gamma
} {
  \Gamma \vdashv x: \tau
}
\and
\inferrule*[right=TVal-Unit] { } {
  \Gamma \vdashv () : \one
}
\and
\inferrule*[right=TVal-True] { } {
  \Gamma \vdashv \mf{true} : \mf{Bool}
}
\and
\inferrule*[right=TVal-False] { } {
  \Gamma \vdashv \mf{false} : \mf{Bool}
}
\and
\inferrule*[right=TVal-Str] { } {
  \Gamma \vdashv s : \mf{String}
}
\and
\inferrule*[right=TVal-Lam] {
  \Gamma, x: \tau_1 \vdashe e : \tau_2
} {
  \Gamma \vdashv \lambda x: \tau_1.e: \tau_1 \rightarrow \tau_2
}
\and
\inferrule*[right=\hl{TVal-Text}] {
  \Gamma \vdashv v: \tau
  \and
  \tau \in \{\mf{Bool},\mf{String}\}
} {
  \Gamma \vdashv \mf{text}\ v: \mf{View}
}
\and
\inferrule*[right=\hl{TVal-Tag}] {
  \Gamma \vdashv v: \mf{String}
  \and
  (\forall i) \
  \Gamma \vdashv v_i : \mf{Attr}
  \and
  (\forall j) \
  \Gamma \vdashv v_j: \mf{View}
}{
  \Gamma \vdashv \mf{tag}(v, \num{v_i}, \num{v_j}): \mf{View}
}
\and
\inferrule*[right=\hl{TVal-Attr}] {
  \Gamma \vdashv v_1: \mf{String}
  \and
  \Gamma \vdashv v_2: \mf{String}
} {
  \Gamma \vdashv \mf{attr}(v_1, v_2): \mf{Attr}
}
\and
\inferrule*[right=\hl{TVal-OnClick}] {
  \Gamma \vdashv v: \one \rightarrow \one
} {
  \Gamma \vdashv \mf{onClick}\ v: \mf{Attr}
}
\end{mathpar}

\pagebreak

\subsubsection{Expression Typing}

Additions relative to React-tRace core are \hl{highlighted},
although these are just standard rules for references.

A top-level program should have type $\one \rightarrow \mf{View}$.

Judgment: $\Gamma \vdashe e: \tau$

\begin{mathpar}
\inferrule*[right=TExp-Ret] {
  \Gamma \vdashv v: \tau
} {
  \Gamma \vdashe \eret{v}: \tau
}
\and
\inferrule*[right=TExp-Let] {
  \Gamma \vdashe e_1: \tau_1
  \and
  \Gamma, x: \tau_1 \vdashe e_2: \tau_2
} {
  \Gamma \vdashe \elet{x}{e_1}{e_2}: \tau_2
}
\and
\inferrule*[right=TExp-App] {
  \Gamma \vdashv v_1: \tau_1 \rightarrow \tau_2
  \and
  \Gamma \vdashv v_2: \tau_2
} {
  \Gamma \vdashe v_1\ v_2: \tau_2
}
\and
\inferrule*[right=TExp-If] {
  \Gamma \vdashv v : \mf{Bool}
  \and
  \Gamma \vdashe e_1 : \tau
  \and
  \Gamma \vdashe e_2 : \tau
} {
  \Gamma \vdashe \eif{v}{e_1}{e_2} : \tau
}
\and
\inferrule*[right=\hl{TExp-Ref}] {
  \Gamma \vdashv v: \tau
} {
  \Gamma \vdashe \eref{v}: \mf{Ref}\ \tau
}
\and
\inferrule*[right=\hl{TExp-Get}] {
  \Gamma \vdashv v : \mf{Ref}\ \tau
} {
  \Gamma \vdashe \eget{v} : \tau
}
\and
\inferrule*[right=\hl{TExp-Set}] {
  \Gamma \vdashv v_1: \mf{Ref}\ \tau
  \and
  \Gamma \vdashv v_2: \tau
} {
  \Gamma \vdashe \eset{v_1}{v_2}: \one
}
\end{mathpar}

\subsection{Operational Semantics}

\subsubsection{Expression Evaluation}

Judgment: $\sigma, e \mapsto e' \semi \sigma'$
\begin{mathpar}
\inferrule*[right=E-Let] {
  \sigma \semi e_1 \mapsto e_1' \semi \sigma'
} {
  \sigma \semi \elet{x}{e_1}{e_2} \mapsto \elet{x}{e_1'}{e_2} \semi \sigma'
}
\end{mathpar}
\begin{align*}
\sigma \semi \elet{x}{\eret{v}}{e} 
&\mapsto e[v/x] \semi \sigma
&&\ts{E-Ret}\\
\sigma \semi (\lambda x.e) v
&\mapsto e[v/x] \semi \sigma
&&\ts{E-Beta}\\
\sigma \semi \eif{\mf{true}}{e_1}{e_2} 
&\mapsto e_1 \semi \sigma
&&\ts{E-IfTrue}\\
\sigma \semi \eif{\mf{false}}{e_1}{e_2} 
&\mapsto e_2 \semi \sigma
&&\ts{E-IfFalse}\\
\sigma \semi \eref{v} 
&\mapsto \eret{\ell} \semi \sigma[\ell \mapsto v] \quad (\ell \text{ fresh})
&&\ts{E-Ref}\\
\sigma \semi \eget{\ell}
&\mapsto \eret{\sigma[\ell]} \semi \sigma
&&\ts{E-Get}\\
\sigma \semi \eset{\ell}{v}
&\mapsto \eret{()} \semi \sigma[\ell \mapsto v]
&&\ts{E-Set}
\end{align*}

\subsubsection{Tree Annotation}

Judgment: $\mf{annotate}(v) = t$

\begin{mathpar}
\inferrule*[right=Annot-Text] { } {
  \mf{annotate}(\mf{text}\ v) = \mf{lit}\ v
}
\and
\inferrule*[right=Annot-Tag] {
    id \text{ fresh}
    \and
    (\forall j)\
    \mf{annotate}(v_j) = t_j
} {
  \mf{annotate}(\mf{tag}(v, \num{v_i}, \num{v_j})) 
  = \mf{node}(id, v, \num{v_i}, \num{t_j})\\
}
\end{mathpar}

\subsubsection{Handler Search}

Judgment: $\mf{handlers}(t, id) = \num{v}$

\begin{mathpar}
\inferrule*[right=Handlers-Lit] { } {
  m \vdash \mf{handlers}(\mf{lit}\ v, id) = \{\}
}
\and
\inferrule*[right=Handlers-Tgt] { } {
  m \vdash \mf{handlers}(\mf{node}(id, v, \num{v_i}, \num{v_j}), id)
  = \{v_h \mid \mf{onClick}\ v_h \in \num{v_i}\}
}
\and
\inferrule*[right=Handlers-Node] {
  id \neq id'
  \and
  (\forall j)\ 
  m \vdash \mf{handlers}(v_j, id) = h_j
} {
  m \vdash \mf{handlers}(\mf{node}(id', v, \num{v_i}, \num{v_j}), id)
  = \bigcup_{\forall j} h_j
}
\end{mathpar}

\subsubsection{Reactive Semantics}

Judgment: $\four{\sigma}{t}{v}{\mu} \cstep \four{\sigma'}{t'}{v'}{\mu'}$

\begin{mathpar}
\inferrule*[right=R-Start] {
  \cdot \semi e \mapsto^* \eret{v} \semi \sigma
} {
  \four{\cdot}{\cdot}{\cdot}{\mf{start}(e)}
  \cstep
  \four{\sigma}{\cdot}{v}{\mf{render}}
}
\and
\inferrule*[right=R-Render] {
  \sigma \semi v_r\; () \mapsto^* \eret{v} \semi \sigma'
  \and
  \mf{annotate}(v) = t
} {
  \four{\sigma}{\_}{v_r}{\mf{render}}
  \cstep
  \four{\sigma'}{t}{v_r}{\mf{idle}}
}
\and
\inferrule*[right=R-Event] {
  \exists id
  \and
  \mf{handlers}(t, id) = \num{v_i}
  \\
  (0 \leq i < n) \quad
  \sigma_i \semi v_i\; () \mapsto^* \eret{()} \semi \sigma_{i+1}
} {
  \four{\sigma_0}{t}{v_r}{\mf{idle}}
  \cstep
  \four{\sigma_n}{t}{v_r}{\mf{render}}
}
\end{mathpar}

\subsection{Type Soundness (Expressions)}

\subsubsection{Store Context}

We introduce the additional context $\Sigma=\map{\ell}{\tau}$
to type standalone locations.

Define $\Gamma \semi \Sigma \vDash \sigma$ to say that 
$\mf{dom}(\Sigma) = \mf{dom}(\sigma)$ and
$\Gamma \semi \Sigma \vdashv \sigma[\ell]: \Sigma[\ell]$.

Further define $\Sigma \sqsubseteq \Sigma'$ to hold when,
for all $\ell \in \Sigma$, both $\ell \in \Sigma'$ and $\Sigma(\ell) = \Sigma'(\ell)$.
Overload $\sqsubseteq$ analogously for type contexts $\Gamma$.

\subsubsection{Typing Rules}

We extend all typing rules with $\Sigma$, and introduce
a new rule \ts{TVal-Loc} to use it.

Augmented value judgment: $\Gamma \semi \Sigma \vdashv v: \tau$

\begin{mathpar}
\inferrule*[right=\hl{TVal-Loc}] {
  \ell : \tau \in \Sigma
} {
  \Gamma \semi \Sigma \vdashv \ell: \mf{Ref}\ \tau
}
\and
\inferrule*[right=TVal-Var] {
  x : \tau \in \Gamma
} {
  \Gamma \semi \Sigma \vdashv x: \tau
}
\and
\inferrule*[right=TVal-Unit] { } {
  \Gamma \semi \Sigma \vdashv () : \one
}
\and
\inferrule*[right=TVal-True] { } {
  \Gamma \semi \Sigma \vdashv \mf{true} : \mf{Bool}
}
\and
\inferrule*[right=TVal-False] { } {
  \Gamma \semi \Sigma \vdashv \mf{false} : \mf{Bool}
}
\and
\inferrule*[right=TVal-Str] { } {
  \Gamma \semi \Sigma \vdashv s : \mf{String}
}
\and
\inferrule*[right=TVal-Lam] {
  \Gamma, x: \tau_1 \semi \Sigma \vdashe e : \tau_2
} {
  \Gamma \semi \Sigma \vdashv \lambda x: \tau_1.e: \tau_1 \rightarrow \tau_2
}
\and
\inferrule*[right=TVal-Text] {
  \Gamma \semi \Sigma \vdashv v: \tau
  \and
  \tau \in \{\mf{Bool},\mf{String}\}
} {
  \Gamma \semi \Sigma \vdashv \mf{text}\ v: \mf{View}
}
\and
\inferrule*[right=TVal-Tag] {
  \Gamma \semi \Sigma \vdashv v: \mf{String}
  \\
  (\forall i) \
  \Gamma \semi \Sigma \vdashv v_i : \mf{Attr}
  \and
  (\forall j) \
  \Gamma \semi \Sigma \vdashv v_j: \mf{View}
}{
  \Gamma \semi \Sigma \vdashv \mf{tag}(v, \num{v_i}, \num{v_j}): \mf{View}
}
\and
\inferrule*[right=TVal-Attr] {
  \Gamma \semi \Sigma \vdashv v_1: \mf{String}
  \and
  \Gamma \semi \Sigma \vdashv v_2: \mf{String}
} {
  \Gamma \semi \Sigma \vdashv \mf{attr}(v_1, v_2): \mf{Attr}
}
\and
\inferrule*[right=TVal-OnClick] {
  \Gamma \semi \Sigma \vdashv v: \one \rightarrow \one
} {
  \Gamma \semi \Sigma \vdashv \mf{onClick}\ v: \mf{Attr}
}
\end{mathpar}

Augmented expression judgment: $\Gamma \semi \Sigma \vdashe e: \tau$

\begin{mathpar}
\inferrule*[right=TExp-Ret] {
  \Gamma \semi \Sigma \vdashv v: \tau
} {
  \Gamma \semi \Sigma \vdashe \eret{v}: \tau
}
\and
\inferrule*[right=TExp-Let] {
  \Gamma \semi \Sigma \vdashe e_1: \tau_1
  \and
  \Gamma, x: \tau_1 \semi \Sigma \vdashe e_2: \tau_2
} {
  \Gamma \semi \Sigma \vdashe \elet{x}{e_1}{e_2}: \tau_2
}
\and
\inferrule*[right=TExp-App] {
  \Gamma \semi \Sigma \vdashv v_1: \tau_1 \rightarrow \tau_2
  \and
  \Gamma \semi \Sigma \vdashv v_2: \tau_2
} {
  \Gamma \semi \Sigma \vdashe v_1\ v_2: \tau_2
}
\and
\inferrule*[right=TExp-If] {
  \Gamma \semi \Sigma \vdashv v : \mf{Bool}
  \and
  \Gamma \semi \Sigma \vdashe e_1 : \tau
  \and
  \Gamma \semi \Sigma \vdashe e_2 : \tau
} {
  \Gamma \semi \Sigma \vdashe \eif{v}{e_1}{e_2} : \tau
}
\and
\inferrule*[right=TExp-Ref] {
  \Gamma \semi \Sigma \vdashv v: \tau
} {
  \Gamma \semi \Sigma \vdashe \eref{v}: \mf{Ref}\ \tau
}
\and
\inferrule*[right=TExp-Get] {
  \Gamma \semi \Sigma \vdashv v : \mf{Ref}\ \tau
} {
  \Gamma \semi \Sigma \vdashe \eget{v} : \tau
}
\and
\inferrule*[right=TExp-Set] {
  \Gamma \semi \Sigma \vdashv v_1: \mf{Ref}\ \tau
  \and
  \Gamma \semi \Sigma \vdashv v_2: \tau
} {
  \Gamma \semi \Sigma \vdashe \eset{v_1}{v_2}: \one
}
\end{mathpar}

\subsubsection{Unproved Lemmas}

We assume that the typical structural rules for intuitionistic calculi,
i.e. weakening, contraction, and exchange, hold on both $\Gamma$ and $\Sigma$.
For example:

\begin{lemma}[Weakening $\Gamma$ in $\vdashv$]
If $\Gamma \semi \Sigma \vdashv v: \tau$
then $\Gamma' \semi \Sigma \vdashv v: \tau$ 
for any $\Gamma'$ such that $\Gamma \sqsubseteq \Gamma'$.
\end{lemma}

\begin{proof}
By induction on the derivation of the typing judgment.
\end{proof}

We also assume the following canonical forms lemma on values:

\begin{lemma}[Canonical Forms]
If $\cdot \semi \Sigma \vdashv v: \tau$, then
\begin{itemize}
\item If $\tau=\mf{Bool}$, then $v=\mf{true}$ or $v=\mf{false}$.
\item If $\tau=\one$, then $v=()$.
\item If $\tau=\mf{String}$, then $v=s$ for some string $s$.
\item If $\tau=\tau_1 \rightarrow \tau_2$, then $v=\lambda x: \tau_1.e$.
\item If $\tau=\mf{Attr}$, then $v=\mf{attr}(v_1, v_2)$ or $v=\mf{onClick} v$.
\item If $\tau=\mf{View}$, then $v=\mf{text}\ v$ or $v=\mf{tag}(v, \num{v_i}, \num{v_j})$. 
\item If $\tau=\mf{Ref}\ \tau'$, then $v=\ell$.
\end{itemize}
\end{lemma}

\subsubsection{Substitution}

\begin{lemma}[Type preservation under substitution]
If $\Gamma \semi \Sigma \vdashv v: \tau$, then the following two conditionals hold:
\begin{enumerate}[label=(\roman*)]
\item If $\Gamma, x: \tau \semi \Sigma \vdashv v': \tau'$, 
then $\Gamma \semi \Sigma \vdashv v'[v/x]: \tau'$.
\item If $\Gamma, x: \tau \semi \Sigma \vdashe e: \tau'$, 
then $\Gamma \semi \Sigma \vdashe e[v/x]: \tau'$.
\end{enumerate}
\end{lemma}

\begin{proof}
By mutual induction on \emph{(i)} and \emph{(ii)}. For each,
we induct over the derivation of the respective typing judgment
in the antecedent.

\tb{Case} \ts{TVal-Var}:
We have $\Gamma, x: \tau \semi \Sigma \vdashv x: \tau'$.
Applying substitution, we want to show 
$\Gamma \semi \Sigma \vdashv v: \tau'$, which is exactly our first premise
(with $\tau'=\tau$).

\tb{Case} \ts{TVal-Unit}:
We have $\Gamma, x: \tau \semi \Sigma \vdashv (): \one$.
Trivial, since $()[v/x]=()$ and $\Gamma \semi \Sigma \vdashv (): \one$ axiomatically.
We elide analogous cases for boolean and string constants.

\tb{Case} \ts{TVal-Loc}:
We have $\Gamma, x: \tau \semi \Sigma \vdashv \ell: \tau$.
By inversion, $\ell: \tau \in \Sigma$,
and since $\ell[v/x]=\ell$, we get $\Gamma \semi \Sigma \vdashv \ell: \tau$.

\tb{Case} \ts{TVal-Lam}:
We have $\Gamma, x: \tau \semi \Sigma \vdashv \lambda y.e : \tau_1 \rightarrow \tau_2$.
First suppose $y=x$: by definition of substitution, $(\lambda x.e)[v/x] = \lambda x.e$,
and the conclusion follows from inversion \& reapplication of \ts{TVal-Lam},
noting that $x: \tau$ is overridden by $x: \tau_1$ in the premise.

If $y \neq x$, then $(\lambda y.e)[v/x] = \lambda y.(e[v/x])$,
and inversion gives us $\Gamma, x: \tau, y: \tau_1 \semi \Sigma \vdashe e: \tau_2$.
Applying case \emph{(ii)} of the inductive hypothesis, alongside exchange,
gives that $\Gamma, y: \tau_1 \semi \Sigma \vdashe e[v/x]: \tau_2$.
We conclude $\Gamma \semi \Sigma \vdashv (\lambda y.e) [v/x]: \tau_1 \rightarrow \tau_2$
by reapplication of \ts{TVal-Lam}.

\tb{Case} \ts{TVal-Tag}:
We have  $\Gamma, x: \tau \semi \Sigma \vdashv 
\mf{tag}(v', \num{v_i}, \num{v_j}): \mf{View}$.
By inversion, $\Gamma, x: \tau \semi \Sigma \vdashv v': \mf{String}$,
$\Gamma, x: \tau \semi \Sigma \vdashv v_i: \mf{Attr}$,
and $\Gamma, x: \tau \semi \Sigma \vdashv v_j: \mf{View}$
for all $i, j$.
We apply case \emph{(i)} of the inductive hypothesis to each, to obtain
$\Gamma \semi \Sigma \vdashv v'[v/x]: \mf{String}$,
$\Gamma \semi \Sigma \vdashv v_i[v/x]: \mf{Attr}$,
and $\Gamma \semi \Sigma \vdashv v_j[v/x]: \mf{View}$, respectively. 
By \ts{TVal-Tag} and definition of substitution on $\mf{tag}$,
we conclude $\Gamma \semi \Sigma \vdashv 
\mf{tag}(v', \num{v_i}, \num{v_j})[v/x]: \mf{View}$.
We elide the simpler but analogous cases for data constructors 
\textsf{text}, \textsf{attr}, and \textsf{onClick}.

\tb{Case} \ts{TExp-Ret}:
We have $\Gamma, x: \tau \semi \Sigma \vdashe \eret{v'}: \tau'$.
By inversion, $\Gamma, x: \tau \semi \Sigma \vdashv v': \tau'$.
Case \emph{(i)} of the inductive hypothesis gives $\Gamma \semi \Sigma \vdashv v'[v/x]: \tau'$,
and reapplication \& definition of substitution let us conclude
$\Gamma \semi \Sigma \vdashe (\eret{v'})[v/x]: \tau'$.

\tb{Case} \ts{TVal-Let}:
$\Gamma, x: \tau \semi \Sigma \vdashe
\elet{y}{e_1}{e_2}: \tau_2$. Inversion gives 
$\Gamma, x: \tau \semi \Sigma \vdashe e_1: \tau_1$.
We have two subcases, when $y=x$ and when not.
In either case, $\Gamma \semi \Sigma \vdashe e_1[v/x] : \tau_1$
by inductive hypothesis.

If $y=x$, then $(\elet{y}{e_1}{e_2})[v/x]=\elet{y}{e_1[v/x]}{e_2}$.
We reapply \ts{TExp-Let} with the second premise obtained via inversion and
we are done. (Like the $\lambda$ case, we note the overwriting of $x: \tau$ in
the second premise.)

If $y \neq x$, then we apply the inductive hypothesis to the second premise,
giving $\Gamma, x: \tau, y: \tau_1 \semi \Sigma \vdashe e_2: \tau_2$, and
by inductive hypothesis and exchange, $\Gamma, y: \tau_1 \vdashe e_2[v/x]: \tau_2$.
Reapplication of \ts{TExpr-Let} and definition of substitution give
$\Gamma \semi \Sigma \vdashe (\elet{y}{e_1}{e_2})[v/x]: \tau_2$.

We elide the remaining cases for expressions, which are analogous to the above.

\end{proof}

\subsubsection{Progress}

\begin{theorem}[Progress of expressions]
If $\cdot \semi \Sigma \vdashe e: \tau$, then either:
\begin{enumerate}[label=(\roman*)]
\item $e = \eret{v}$, or
\item for all $\sigma$ such that $\cdot \semi \Sigma \vDash \sigma$, 
there exist some $e'$ and $\sigma'$ such that 
$\sigma \semi e \mapsto e' \semi \sigma'$.
\end{enumerate}
\end{theorem}

\begin{proof}
By induction on the derivation of the typing judgment.

\tb{Case} \ts{TExp-Ret}: 
We have $\cdot \semi \Sigma \vdashe \eret{v} \semi \tau$.
The theorem holds trivially by case \emph{(i)}.

\tb{Case} \ts{TExp-Let}:
We have $\cdot \semi \Sigma \vdashe \elet{x}{e_1}{e_2}: \tau_2$.
By inversion, $\cdot \semi \Sigma \vdashe e_1: \tau_1$
and $x: \tau_1 \semi \Sigma \vdashe e_2: \tau_2$.
By inductive hypothesis on the first premise,
we have two subcases on $e_1$: either $e_1=\eret{v}$ 
or $\sigma \semi e_1 \mapsto e_1' \semi \sigma'$
(for $\cdot \semi \Sigma \vDash \sigma$).

If $e_1=\eret{v}$, then by \ts{E-Ret},
we have $\sigma, \elet{x}{\eret{v}}{e_2} \mapsto e_2[v/x] \semi \sigma$
for any $\sigma$.

If instead $\sigma \semi e_1 \mapsto e_1' \semi \sigma'$,
by \ts{E-Let}, we can conclude that 
$\sigma \semi \elet{x}{e_1}{e_2} \mapsto \elet{x}{e_1'}{e_2} \semi \sigma'$.

Thus, in either subcase, $\elet{x}{e_1}{e_2}$ steps given a proper $\sigma$,
so the case holds.

\tb{Case} \ts{TExp-App}:
We have $\cdot \semi \Sigma \vdashe v_1\ v_2: \tau_2$.
By inversion, $\cdot \semi \Sigma \vdashv v_1: \tau_1 \rightarrow \tau_2$
and $\cdot \semi \Sigma \vdashv v_2: \tau_1$.
By the canonical forms lemma, $v_1=\lambda x.e$,
and so we apply rule \ts{E-Beta}, i.e.
$\sigma \semi (\lambda x.e) v_2 \mapsto e[v_2/x] \semi \sigma$ for
any $\sigma$ such that $\cdot \semi \Sigma \vDash \sigma$.

\tb{Case} \ts{TExp-If}:
We have $\cdot \semi \Sigma \vdashe \eif{v}{e_1}{e_2} : \tau$.
By inversion, $\cdot \semi \Sigma \vdashv v: \mf{Bool}$,
and by canonical forms, either $v=\mf{true}$ or $v=\mf{false}$.
If $v=true$, then \ts{E-IfTrue} says 
$\sigma \semi \eif{v}{e_1}{e_2} \mapsto e_1 \semi \sigma$,
and if $v=false$, then $\sigma \semi \eif{v}{e_1}{e_2} \mapsto e_2 \semi \sigma$
by \ts{E-IfFalse}. In either case, we satisfy \emph{(ii)}.

\tb{Case} \ts{TExp-Ref}:
We have $\cdot \semi \Sigma \vdashe \eref{v} : \mf{Ref}\ \tau$.
Immediately apply \ts{E-Ref} to get that 
$\sigma, \eref{v} \mapsto \eret{\ell} \semi \sigma[\ell \mapsto v]$,
which satisfies \emph{(ii)}.

\tb{Case} \ts{TExp-Get}:
We have $\cdot \semi \Sigma \vdashe \eget{v} : \tau$.
By inversion, $\cdot \semi \Sigma \vdashv v: \mf{Ref}\ \tau$,
and by canonical forms, $v=\ell$.
Again by inversion (on \ts{TVal-Loc}), $\ell: \tau \in \Sigma$.
Since $\cdot \semi \Sigma \vDash \sigma$ for our arbitrary $\sigma$,
we have that $\sigma[\ell]$ exists.
Thus, we apply \ts{E-Get} to say
$\sigma \semi \eget{v} \mapsto \eret{\sigma[\ell]} \semi \sigma$.

\tb{Case} \ts{TExp-Set}:
We have $\cdot \semi \Sigma \vdashe \eset{v_1}{v_2}: \one$.
By inversion, $\cdot \semi \Sigma \vdashv v_1: \mf{Ref}\ \tau$,
and by canonical forms, $v=\ell$.
Applying \ts{E-Set} then gives 
$\sigma \semi \eset{v_1}{v_2} \mapsto \eret{()} \semi \sigma[\ell \mapsto v_2]$.

\end{proof}

\subsubsection{Preservation}

\begin{theorem}[Preservation of expressions]
If
\begin{enumerate}[label=(\roman*)]
\item $\Gamma \semi \Sigma \vdashe e: \tau$,
\item $\Gamma \semi \Sigma \vDash \sigma$ for some $\sigma$, and
\item $\sigma \semi e \mapsto e' \semi \sigma'$ for some $e'$ and $\sigma'$,
\end{enumerate}
then there is some $\Sigma'$ such that:
\begin{itemize}
\item $\Sigma \sqsubseteq \Sigma'$,
\item $\Gamma \semi \Sigma' \vDash \sigma'$, and
\item $\Gamma \semi \Sigma' \vdashe e': \tau$.
\end{itemize}
\end{theorem}

\begin{proof}
By induction on derivation of \emph{(iii)}.

\tb{Case} \ts{E-Let}:
We have $\sigma \semi \elet{x}{e_1}{e_2} \mapsto \elet{x}{e_1'}{e_2} \semi \sigma'$.
By inversion on \ts{E-Let}, we also have $\sigma \semi e_1 \mapsto e_1' \semi \sigma'$.
By inversion on \emph{(i)}, $\Gamma \semi \Sigma \vdashe e_1: \tau_1$.
By inductive hypothesis, there exists some $\Sigma'$
with $\Gamma \semi \Sigma' \vDash \sigma'$
and $\Gamma \semi \Sigma' \vdashe e_1: \tau$.
Since $\Sigma \sqsubseteq \Sigma'$, and $\Gamma, x: \tau_1 \semi \Sigma \vdashe e_2: \tau$,
we also get that $\Gamma, x: \tau_1 \semi \Sigma' \vdashe e_2: \tau_2$.
We apply \ts{E-Let} to get $\Gamma \semi \Sigma' \vdashe \elet{x}{e_1'}{e_2}: \tau$.

\tb{Case} \ts{E-Ret}:
We have $\sigma \semi \elet{x}{\eret{v}}{e} \mapsto e[v/x] \semi \sigma$.
Inversion on \emph{(i)} gives that $\Gamma \semi \Sigma \vdashv v: \tau_1$ and
$\Gamma, x: \tau_1 \semi \Sigma \vdashe e: \tau$. We conclude
$\Gamma \semi \Sigma \vdashe e[v/x]: \tau$ via the substitution lemma.

\tb{Case} \ts{E-Beta}:
We have $\sigma \semi (\lambda x: \tau_1.e)v \mapsto e[v/x] \semi \sigma$.
Inversion on \emph{(i)} gives $\Gamma, x: \tau_1 \semi \Sigma \vdashe e: \tau$
and $\Gamma \semi \Sigma \vdashv v: \tau_1$. We again get
$\Gamma \semi \Sigma \vdashe e[v/x]: \tau$ via the substitution lemma.

\tb{Case} \ts{E-IfTrue}:
We have $\sigma \semi \eif{\mf{true}}{e_1}{e_2} \mapsto e_1 \semi \sigma$.
Immediate by inversion on \emph{(i)}. Case for \ts{E-IfFalse} is analogous.

\tb{Case} \ts{E-Ref}:
We have $\sigma \semi \eref{v} \mapsto \eret{\ell} \semi \sigma[\ell \mapsto v]$
(for $\ell$ fresh).
By \emph{(i)} and enumeration of rules, 
$\Gamma \semi \Sigma \vdashe \eref{v}: \mf{Ref}\ \tau$.
Inversion gives $\Gamma \semi \Sigma \vdashv v: \tau$.

Choose $\Sigma' = \Sigma[\ell \mapsto \tau]$; 
since $\ell$ is fresh, clearly $\Sigma \sqsubseteq \Sigma'$.
Since both $\Sigma'$ and $\sigma'$ are extensions of
$\Sigma$ and $\sigma$ by fresh label $\ell$, by definition of \emph{(ii)},
to show $\Gamma \semi \Sigma' \vDash \sigma'$,
it suffices to show $\cdot \semi \Sigma' \vdashv \sigma'[\ell]: \Sigma'[\ell]$.
Now since $\sigma'[\ell]=v$ and we said $\Sigma'[\ell] = \tau$,
we need only apply the proposition from the previous paragraph.

We also need to show that $\Gamma \semi \Sigma' \vdashe \eret{\ell} : \tau$.
This follows by \ts{TExp-Ret}, \ts{TVal-Loc}, and $\Sigma'[\ell] = \tau$.

\tb{Case} \ts{E-Get}:
We have $\sigma \semi \eget{\ell} \mapsto \eret{\sigma[\ell]} \semi \sigma$.

Choose $\Sigma'=\Sigma$; 
trivially $\Gamma \semi \Sigma' \vDash \sigma$ and $\Sigma \sqsubseteq \Sigma'$.

By inversion on \emph{(i)}, $\Gamma \semi \Sigma \vdashv \ell: \mf{Ref}\ \tau$
and $\ell: \tau \in \Sigma$.
We can then apply \ts{TVal-Loc} and \ts{TExp-Ret}
to get $\Gamma \semi \Sigma \vdashe \eret{\ell}: \tau$.

\tb{Case} \ts{E-Set}:
We have $\sigma \semi \eset{\ell}{v} \mapsto \eret{()} \semi \sigma[\ell \mapsto v]$.

Again choose $\Sigma'=\Sigma$. 
It is easy to show $\Gamma \semi \Sigma \vdashe \eret{()} : \one$,
and by enumeration of rules and \emph{(i)} 
we have $\Gamma \semi \Sigma \vdashe \eset{\ell}{v}: \one$.

We also need to show $\Gamma \semi \Sigma \vDash \sigma[\ell \mapsto v]$.
By \emph{(ii)}, this holds for all $\ell' \neq \ell$.
For $\ell$, we want to show $\cdot \semi \Sigma \vdashv \sigma'[\ell]: \Sigma[\ell]$,
i.e. that $\cdot \semi \Sigma \vdashv v: \Sigma[\ell]$.
Inversion on \emph{(i)} gives us both that $\Sigma[\ell] = \tau$
and that $\Gamma \semi \Sigma \vdashv v: \tau$.
\end{proof}

\subsubsection{Soundness}

\begin{theorem}[Type soundness for expressions]
If
\begin{itemize}
\item $\cdot \semi \Sigma \vdashe e: \tau$,
\item $\Sigma \vDash \sigma$ for arbitrary $\sigma$,
\item $\sigma \semi e \mapsto^* e' \semi \sigma'$ for some $e'$ and $\sigma'$,
\item and $e' \not \mapsto$ (i.e. $e'$ cannot step further),
\end{itemize}
then $e' = \eret{v}$, and for some $\Sigma'$ with $\Sigma \subseteq \Sigma'$
and $\Sigma' \vDash \sigma'$, we get $\cdot \semi \Sigma' \vdashv v: \tau$.
\end{theorem}

\begin{proof}
Follows directly from progess \& preservation.
\end{proof}

\begin{corollary}
If $\cdot \semi \cdot \vdashe e: \tau$ 
and $\cdot \semi e \mapsto^* e' \semi \sigma$
(and $e' \not \mapsto$),
then $e' = \eret{v}$,
and for any $\Sigma$ with $\Sigma \vDash \sigma$,
we get $\cdot \semi \Sigma \vdashv v: \tau$.
\end{corollary}

\subsection{Configuration Typing}



\section{Translation}

From React-tRace Core into $\lambda_{CC}$.

This requires the addition of (mutual) recursion into $\lambda_{CC}$,
which is entirely standard.

\subsection{Types}
Judgment: $\llrr{\tau_{RTC}} = \tau_{CC}$

\begin{align*}
\llrr{\one} &= \one\\
\llrr{\mf{Bool}} &= \mf{Bool}\\
\llrr{\mf{View}} &= \one \rightarrow \mf{View}\\
\llrr{\mf{Attr}} &= \mf{Attr}\\
\llrr{\tau_1 \xrightarrow{E} \tau_2} &= \llrr{\tau_1} \rightarrow \llrr{\tau_2}\\
\llrr{\mf{Component}\ \tau} &= \llrr{\tau} \rightarrow \llrr{\mf{View}}\\
\llrr{\mf{Setter}\ \tau} &= \llrr{\tau} \rightarrow \one
\end{align*}

\subsection{Values}

Judgment: $\llrr{v_{RTC}} = v_{CC}$
\begin{align*}
\llrr{x} &= x\\
\llrr{()} &= ()\\
\llrr{\mf{true}} &= \mf{true}\\
\llrr{\mf{false}} &= \mf{false}\\
\llrr{s} &= s\\
\llrr{\lambda x: \tau . e} & = \lambda x: \llrr{\tau}.\llrr{e}\\
\llrr{\mf{text}\ v} &= \mf{text}\ \llrr{v}\\
\llrr{\mf{tag}(v, \num{v_i}, \num{v_j})} 
&= \mf{tag}(\llrr{v}, \num{\llrr{v_i}}, \num{\llrr{v_j}})\\
\llrr{\mf{attr}(v_1, v_2)} &= \mf{attr}(\llrr{v_1}, \llrr{v_2})\\
\llrr{\mf{onClick}\ v} &= \mf{onClick}\ \llrr{v}\\
\llrr{\mc{C}} &= x_\mc{C}
\end{align*}

\subsection{Expressions}
Judgment: $\llrr{e_{RTC}} = e_{CC}$
\begin{align*}
\llrr{\eret{v}} &= \eret{\llrr{v}}\\
\llrr{\elet{x}{e_1}{e_2}} &= \elet{x}{\llrr{e_1}}{\llrr{e_2}}\\
\llrr{v_1\ v_2} &= \llrr{v_1}\ \llrr{v_2}\\
\llrr{\eif{v}{e_1}{e_2}} &= \eif{\llrr{v}}{\llrr{e_1}}{\llrr{e_2}}\\
\end{align*}

\subsection{Top-level Expressions}
Judgment: $\Theta \vdash \llrr{\hat{e}_{RTC}} = \hat{e}_{CC}$
where $\Theta ::= \cdot \mid (x, \ell), \Theta$

Also define $\Theta - x = \Theta \setminus \{x, \ell\}$ for any $\ell$
\begin{mathpar}
\inferrule*{
  \Theta - x \vdash \llrr{\hat{e}} = e'
} {
  \Theta \vdash \llrr{\elet{x}{e}{\hat{e}}} = \elet{x}{e}{e'}
}
\and
\inferrule*{
  \ell_x\ \mf{fresh}
  \and
  \Theta, (x, \ell_x) \vdash \llrr{\hat{e}} = e'
} {
  \Theta \vdash \llrr{\estate{x}{\ell}{v}{\hat{e}}} =\\
  \elet{\ell_x}{\eref{\llrr{v}}}{
    \elet{x_\mf{set}}{\lambda y: ?.\eset{\ell_x}{y}}{e'}
  }
}
\and
\inferrule*{
  \Theta \vdash \mf{deref}(\llrr{e}) = e'
} {
  \Theta \vdash \llrr{\hat{e}} = e'
}
\and
\inferrule*{
  \Theta \vdash \mf{deref}(e)
} {
  \Theta, (x, \ell_x) \vdash \mf{deref}(e) 
  = \elet{x}{\eget{\ell_x}}{e'}
}
\and
\inferrule*{ } {
  \cdot \vdash \mf{deref}(e) = e
}
\end{mathpar}

\subsection{Programs}
Judgment: $\Omega \vdash \llrr{P} = e$ 
where $\Omega = \cdot \mid \four{\mc{C}}{x}{\tau}{\hat{e}}, \Omega$
\begin{mathpar}
\inferrule*{
  \Omega \vdash \llrr{P} = e'
} {
  \Omega \vdash \llrr{\elet{x}{e}{P}} = \elet{x}{\llrr{e}}{e'}
}
\and
\inferrule*{
  \Omega, \four{\mc{C}}{x}{\tau}{\hat{e}} \vdash \llrr{P} = e'
} {
  \Omega \vdash \llrr{\elet{\mc{C}(x: \tau)}{\hat{e}}{P}} = e'
}
\and
\inferrule*{
  \forall \four{\mc{C}_i}{x_i}{\tau_i}{\hat{e}_i} \in \Omega
  \qquad
  \cdot \vdash \llrr{\hat{e}_i} = e_i
} {
  \Omega \vdash \erun{e} 
  = \eletrec{\num{\mc{C}_i: \llrr{\mf{Component}\ \tau_i}}}
  {\num{\lambda x_i: \llrr{\tau_i} . e_i}}{\llrr{e}}
}
\end{mathpar}

\end{document}